% !TEX root = main.tex
\documentclass[12pt,openany]{book}

% Keep front-, main-, and back-matter transitions consistent with openany.
\let\cleardoublepage\clearpage

\usepackage{amsmath,amssymb,amsthm}
\usepackage[margin=1in]{geometry}
\usepackage[shortlabels]{enumitem}
\usepackage[nottoc]{tocbibind}
\usepackage{xcolor}
\usepackage[hidelinks]{hyperref}

\setlength{\emergencystretch}{0.5em}
\raggedbottom
\setlist[enumerate,1]{label=(\roman*)}

%------------------------------------------------
% Temporary notes
%------------------------------------------------
\newif\ifshowdraftnotes
\showdraftnotestrue % Change to \showdraftnotesfalse for the final PDF.

\ifshowdraftnotes
  \NewDocumentEnvironment{draftnote}{+b}
    {\par\smallskip\noindent\begingroup
     \setlength{\fboxsep}{6pt}%
     \fcolorbox{gray!55}{gray!12}{%
       \begin{minipage}{\dimexpr\linewidth-2\fboxsep-2\fboxrule\relax}
         \textbf{Draft note: }\ignorespaces#1
       \end{minipage}%
     }%
     \endgroup\par\smallskip}
    {}
\else
  \NewDocumentEnvironment{draftnote}{+b}{}{}
\fi

%------------------------------------------------
% Thesis metadata
%------------------------------------------------
\newcommand{\thesistitle}{The Kechris--Pestov--Todorcevic Correspondence}
\newcommand{\thesissubtitle}{Fraïssé Theory, Ramsey Theory, and Topological Dynamics}
\newcommand{\thesisauthor}{Ali Hushemian}
\newcommand{\thesisdate}{July 2026}

\hypersetup{
  unicode=true,
  pdftitle={The Kechris-Pestov-Todorcevic Correspondence},
  pdfsubject={Fraïssé Theory, Ramsey Theory, and Topological Dynamics},
  pdfauthor={Ali Hushemian}
}


%------------------------------------------------
% Theorem-like environments
%------------------------------------------------
\theoremstyle{plain}                    % default
\newtheorem{theorem}{Theorem}           % numbered within sections
\newtheorem{lemma}[theorem]{Lemma}
\newtheorem{proposition}[theorem]{Proposition}
\newtheorem{corollary}[theorem]{Corollary}
\newtheorem{axiom}[theorem]{Axiom}
\newtheorem{fact}[theorem]{Fact}
\newtheorem*{proposition*}{Proposition}
\newtheorem*{lemma*}{Lemma}
\newtheorem*{corollary*}{Corollary}
\theoremstyle{definition}
\newtheorem{definition}[theorem]{Definition}
\newtheorem{example}[theorem]{Example}
\newtheorem*{example*}{Example}
\newtheorem*{examples*}{Examples}
\newtheorem{examples}[theorem]{Examples}
\newtheorem*{theorem*}{Theorem}
\newtheorem*{definition*}{Definition}

\theoremstyle{remark}
\newtheorem*{remark*}{Remark}
\newtheorem{remark}[theorem]{Remark}
\newtheorem{claim}[theorem]{Claim}
\newtheorem{note}[theorem]{Note}
\newtheorem*{claim*}{Claim}

\theoremstyle{plain}
\newtheorem{exercise}[theorem]{Exercise}
\newtheorem*{exercise*}{Exercise}
\newtheorem{problem}[theorem]{Problem}
\newtheorem{conjecture}[theorem]{Conjecture}
\newtheorem{notation}[theorem]{Notation}
\newtheorem{warning}[theorem]{Warning}
%------------------------------------------------
% Custom macros (add as needed)
%------------------------------------------------

\DeclareMathOperator{\Age}{Age}
\newcommand{\R}{\mathbb{R}}
\newcommand{\C}{\mathbb{C}}
\newcommand{\Z}{\mathbb{Z}}
\newcommand{\Q}{\mathbb{Q}}
\newcommand{\N}{\mathbb{N}}
\newcommand{\U}{\mathbb{U}}
\newcommand{\cK}{\mathcal{K}}
\newcommand{\cLO}{\mathcal{LO}}
\DeclareMathOperator{\Sym}{Sym}
\DeclareMathOperator{\Stab}{Stab}
\DeclareMathOperator{\Orb}{Orb}
\DeclareMathOperator{\rk}{rk}
\DeclareMathOperator{\dom}{dom}
\newcommand{\nset}[2]{\left[{#1}\right]^{#2}} % n-element subsets
\newcommand{\kto}[4]{{#1}\to({#2})^{#3}_{#4}} % arrow notation
% \newcommand{\Aut}{\operatorname{Aut}}
\DeclareMathOperator{\id}{id}
\DeclareMathOperator{\Flim}{Flim}
\DeclareMathOperator{\LO}{LO}
\DeclareMathOperator{\Forb}{Forb}
\DeclareMathOperator{\Adm}{Adm}
\DeclareMathOperator{\Homeo}{Homeo}
\DeclareMathOperator{\Iso}{Iso}
\DeclareMathOperator{\pat}{pat}
\DeclareMathOperator{\tp}{tp}

\newcommand{\ncl}[2]{\left\langle\!\left\langle #1 \right\rangle\!\right\rangle^{#2}}
\DeclareMathOperator{\Hom}{Hom}
\DeclareMathOperator{\im}{im}
\DeclareMathOperator{\Aut}{Aut}
\DeclareMathOperator{\Def}{Def}
\DeclareMathOperator{\PF}{PF}
\renewcommand{\qedsymbol}{\ensuremath{\dashv}}
\DeclareMathOperator{\Fix}{Fix}

\DeclareMathOperator{\DCL}{DCL}
\DeclareMathOperator{\ACL}{ACL}
\DeclareMathOperator{\dcl}{dcl}
\DeclareMathOperator{\acl}{acl}
\DeclareMathOperator{\Ker}{Ker}
\DeclareMathOperator{\End}{End}
\DeclareMathOperator{\Rad}{Rad}

\newcommand{\bfa}{\boldsymbol{A}}
\newcommand{\bfb}{\boldsymbol{B}}
\newcommand{\bfc}{\boldsymbol{C}}
\newcommand{\bfd}{\boldsymbol{D}}
\newcommand{\bfe}{\boldsymbol{E}}
\newcommand{\bff}{\boldsymbol{F}}
\newcommand{\bfh}{\boldsymbol{H}}
\newcommand{\bfo}{\boldsymbol{O}}
\newcommand{\bfr}{\boldsymbol{R}}
\newcommand{\bfu}{\boldsymbol{U}}
\newcommand{\bfx}{\boldsymbol{X}}
\newcommand{\bfv}{\boldsymbol{V}}
\newcommand{\bffa}{\boldsymbol{a}}
\newcommand{\bffb}{\boldsymbol{b}}
\newcommand{\bfzero}{\boldsymbol{0}}
\newcommand{\bfone}{\boldsymbol{1}}
\newcommand{\norm}[1]{\left\lVert #1\right\rVert}


\numberwithin{theorem}{section}

\begin{document}

\frontmatter

\begin{titlepage}
  \thispagestyle{empty}
  \begin{center}
    \vspace*{0.12\textheight}

    {\Large\textsc{Thesis}\par}

    \vspace{1.8cm}
    \rule{0.78\textwidth}{0.6pt}\par
    \vspace{0.8cm}

    {\Huge\bfseries \thesistitle\par}

    \vspace{0.7cm}
    {\Large \thesissubtitle\par}

    \vspace{0.8cm}
    \rule{0.78\textwidth}{0.6pt}\par

    \vfill

    {\large by\par}
    \vspace{0.35cm}
    {\LARGE\thesisauthor\par}

    \vfill

    {\large\thesisdate\par}
  \end{center}
\end{titlepage}

\tableofcontents

\mainmatter

% !TEX root = main.tex
\chapter{Fraïssé Classes and Fraïssé Limits}

This chapter develops the Fraïssé-theoretic background needed for the rest of the thesis. We introduce the basic notions of age, hereditary property, joint embedding property, and amalgamation property, and we explain how these conditions characterize classes of finite structures arising from countable ultrahomogeneous structures. We then state and prove Fraïssé's theorem, which yields the existence and uniqueness of the Fraïssé limit of a Fraïssé class. These constructions provide the model-theoretic foundation for the later study of automorphism groups, Ramsey properties, and universal minimal flows.

\section{Preliminaries on structures and embeddings}

Throughout this chapter, \(L\) denotes a signature. An \(L\)-structure \(A\) consists of a set \(\mathrm{dom}(A)\) together with interpretations of the constant, relation, and function symbols of \(L\). If \(c\) is a constant symbol, \(R\) an \(n\)-ary relation symbol, and \(F\) an \(n\)-ary function symbol, their interpretations in \(A\) are denoted by
\[
    c^A\in \mathrm{dom}(A),\qquad R^A\subseteq \mathrm{dom}(A)^n,\qquad F^A:\mathrm{dom}(A)^n\to \mathrm{dom}(A).
\]
We write \(|A|\) for \(|\mathrm{dom}(A)|\). Constants may, when convenient, be regarded as \(0\)-ary function symbols.

If \(A\) and \(B\) are \(L\)-structures, a map \(f:\mathrm{dom}(A)\to \mathrm{dom}(B)\) is a \emph{homomorphism} if
\begin{enumerate}
    \item \(f(c^A)=c^B\) for every constant symbol \(c\),
    \item whenever \(\bar a\in R^A\), one has \(f(\bar a)\in R^B\),
    \item \(f(F^A(\bar a))=F^B(f(\bar a))\) for every tuple \(\bar a\).
\end{enumerate}
An \emph{embedding} is an injective homomorphism which also reflects relations, that is,
\[
    \bar a\in R^A \iff f(\bar a)\in R^B
\]
for every relation symbol \(R\). A bijective embedding is an \emph{isomorphism}. We write \(A\cong B\) if \(A\) and \(B\) are isomorphic. An \emph{automorphism} of \(A\) is an isomorphism from \(A\) onto itself.


\begin{lemma}\label{lem:isomorphism-basics}
    Let \(A,B,C\) be \(L\)-structures.
    \begin{enumerate}
        \item[(i)] The identity map \(1_A\colon A\to A\) is an isomorphism.
        \item[(ii)] If \(f\colon A\to B\) is an isomorphism, then \(f^{-1}\colon B\to A\) is an isomorphism.
        \item[(iii)] If \(f\colon A\to B\) and \(g\colon B\to C\) are isomorphisms, then \(g\circ f\colon A\to C\) is an isomorphism.
    \end{enumerate}
\end{lemma}

\begin{proof}
    Part (i) is immediate. For (ii), since \(f\) is a bijective embedding, its inverse is again a bijection preserving constants, relations, and functions, and hence is an isomorphism. For (iii), by composition of embeddings, \(g\circ f\) is an embedding, and since it is also a composition of bijections, it is bijective. Therefore \(g\circ f\) is an isomorphism.
\end{proof}

\begin{definition}[Substructure]
    Let \(A\) and \(B\) be \(L\)-structures with \(\mathrm{dom}(A)\subseteq \mathrm{dom}(B)\). We say that \(A\) is a \emph{substructure} of \(B\), and write \(A\subseteq B\), if the inclusion map is an embedding. Equivalently,
    \[
        c^A=c^B,\qquad R^A=R^B\cap \mathrm{dom}(A)^n,\qquad
        F^A=F^B\!\restriction_{\mathrm{dom}(A)^n}
    \]
    for all constant, relation, and function symbols of \(L\).
\end{definition}

\begin{lemma}\label{lem:substructure-characterization}
    Let \(B\) be an \(L\)-structure and let \(X\subseteq \mathrm{dom}(B)\). The following are equivalent:
    \begin{enumerate}
        \item[(i)] \(X=\mathrm{dom}(A)\) for some substructure \(A\subseteq B\);
        \item[(ii)] \(X\) contains all constants of \(B\) and is closed under all functions of \(B\).
    \end{enumerate}
    If these conditions hold, the substructure \(A\) is uniquely determined.
\end{lemma}

\begin{proof}
    If \(A\subseteq B\), then \(\mathrm{dom}(A)\) plainly contains the constants and is closed under the functions of \(B\). Conversely, if \(X\subseteq \mathrm{dom}(B)\) satisfies (ii), define an \(L\)-structure \(A\) on \(X\) by
    \[
        c^A=c^B,\qquad R^A=R^B\cap X^n,\qquad F^A=F^B\!\restriction_{X^n}.
    \]
    Then the inclusion \(A\hookrightarrow B\) is an embedding, so \(A\subseteq B\). Uniqueness is immediate from the definition of substructure.
\end{proof}

\begin{definition}[Generated substructure]
    Let \(B\) be an \(L\)-structure and \(Y\subseteq \mathrm{dom}(B)\). By Lemma~\ref{lem:substructure-characterization}, there is a unique smallest substructure of \(B\) whose domain contains \(Y\). It is denoted by
    \[
        \langle Y\rangle_B
    \]
    and is called the \emph{substructure generated by \(Y\)}. When the ambient structure is clear, we simply write \(\langle Y\rangle\). A structure is \emph{finitely generated} if it is generated by a finite subset of its domain.
\end{definition}

The following observation is used repeatedly: any embedding may be viewed as an inclusion after replacing the codomain by an isomorphic copy.

\begin{lemma}\label{lem:embedding-as-inclusion}
    Let \(f:A\hookrightarrow C\) be an embedding of \(L\)-structures. Then there exist an \(L\)-structure \(B\) and an isomorphism \(h:B\to C\) such that \(A\subseteq B\) and
    \[
        f=h\!\restriction_A .
    \]
\end{lemma}

\begin{proof}
    Let
    \[
        \mathrm{dom}(B):=\mathrm{dom}(A)\,\dot\cup\,\bigl(\mathrm{dom}(C)\setminus f(\mathrm{dom}(A))\bigr),
    \]
    and define a bijection \(\sigma:\mathrm{dom}(B)\to \mathrm{dom}(C)\) by
    \[
        \sigma(a)=f(a)\quad(a\in A),\qquad \sigma(x)=x\quad\bigl(x\in \mathrm{dom}(C)\setminus f(\mathrm{dom}(A))\bigr).
    \]
    Transport the \(L\)-structure of \(C\) to \(\mathrm{dom}(B)\) via \(\sigma\). Then \(\sigma\) becomes an isomorphism \(h:B\to C\). Since \(f\) is an embedding, the induced structure on \(\mathrm{dom}(A)\) is exactly the original structure \(A\), so \(A\subseteq B\) and \(f=h\!\restriction_A\).
\end{proof}

For the examples relevant to Fraïssé theory, the language is usually finite and relational. In that case finitely generated structures are automatically finite.

\begin{proposition}\label{prop:finite-relational-language-finiteness}
    Let \(L\) be a finite relational signature.
    \begin{enumerate}
        \item Every finitely generated \(L\)-structure is finite.
        \item For each \(n\geq 1\), there are only finitely many \(L\)-structures of cardinality \(n\), up to isomorphism.
    \end{enumerate}
\end{proposition}

\begin{proof}
    Since \(L\) has no function symbols and no constants, the substructure generated by a set is just that set itself. Hence every finitely generated \(L\)-structure is finite.

    Now fix \(n\geq 1\). Every \(L\)-structure of size \(n\) is isomorphic to one with domain \(\{1,\dots,n\}\). If \(R_1,\dots,R_k\) are the relation symbols of \(L\), of arities \(m_1,\dots,m_k\), then each \(R_i\) may be interpreted arbitrarily as a subset of \(\{1,\dots,n\}^{m_i}\). There are only finitely many such choices, so only finitely many \(L\)-structures on \(\{1,\dots,n\}\), and therefore only finitely many isomorphism types of size \(n\).
\end{proof}


\section{Ages, HP, and JEP}

We now introduce the first invariant associated with a structure.

\begin{definition}[Age]
    Let \(L\) be a signature and let \(D\) be an \(L\)-structure. The \emph{age} of \(D\), denoted \(\Age(D)\), is the class of all finitely generated \(L\)-structures that embed into \(D\). Equivalently, \(\Age(D)\) is the class of all finitely generated substructures of \(D\), taken up to isomorphism.
\end{definition}

In this thesis we regard ages up to isomorphism. Thus to say that \(D\) has countable age means that there are only countably many isomorphism types of finitely generated substructures of \(D\).

The following two properties are immediate for every age.

\begin{definition}
    Let \(\mathcal K\) be a class of finitely generated \(L\)-structures.
    \begin{enumerate}
        \item[\textup{(HP)}] \(\mathcal K\) has the \emph{hereditary property} if whenever \(A\in\mathcal K\) and \(B\) is a finitely generated substructure of \(A\), then \(B\in\mathcal K\) up to isomorphism.
        \item[\textup{(JEP)}] \(\mathcal K\) has the \emph{joint embedding property} if for all \(A,B\in\mathcal K\), there exists \(C\in\mathcal K\) such that both \(A\) and \(B\) embed into \(C\).
    \end{enumerate}
\end{definition}

\begin{proposition}\label{prop:age-hp-jep}
    Let \(D\) be an \(L\)-structure. Then \(\Age(D)\) is nonempty and satisfies \emph{(HP)} and \emph{(JEP)}.
\end{proposition}

\begin{proof}
    Choose any \(a\in D\). Then the substructure \(\langle a\rangle\) generated by \(a\) is finitely generated and embeds into \(D\), so \(\Age(D)\neq\emptyset\).

    To prove \emph{(HP)}, let \(A\in\Age(D)\), and let \(B\subseteq A\) be finitely generated. Since \(A\) embeds into \(D\), the restriction of any embedding \(A\hookrightarrow D\) gives an embedding \(B\hookrightarrow D\). Hence \(B\in\Age(D)\).

    To prove \emph{(JEP)}, let \(A,B\in\Age(D)\). Choose embeddings \(e_A:A\hookrightarrow D\) and \(e_B:B\hookrightarrow D\). Let
    \[
        C:=\langle e_A(A)\cup e_B(B)\rangle \subseteq D.
    \]
    Since \(A\) and \(B\) are finitely generated, the structure \(C\) is finitely generated. Thus \(C\in\Age(D)\), and both \(A\) and \(B\) embed into \(C\).
\end{proof}

Fraïssé proved a converse: under mild countability assumptions, these two properties already characterize the ages of countable structures.

\begin{proposition}\label{prop:hp-jep-age}
    Let \(L\) be a signature, and let \(\mathcal K\) be a nonempty class of finitely generated \(L\)-structures, countable up to isomorphism. Suppose that \(\mathcal K\) satisfies \emph{(HP)} and \emph{(JEP)}. Then there exists a countable \(L\)-structure \(C\) such that
    \[
        \mathcal K=\Age(C).
    \]
\end{proposition}

\begin{proof}
    Choose a sequence \((A_i)_{i<\omega}\) listing all isomorphism types in \(\mathcal K\), with each type appearing at least once. We construct a chain
    \[
        B_0\subseteq B_1\subseteq B_2\subseteq \cdots
    \]
    such that each \(B_i\in\mathcal K\) up to isomorphism and \(A_i\) embeds into \(B_{i+1}\).

    Set \(B_0=A_0\). Suppose \(B_i\) has been constructed. By \emph{(JEP)}, there is some \(C_i\in\mathcal K\) into which both \(B_i\) and \(A_{i+1}\) embed. By Lemma~\ref{lem:embedding-as-inclusion}, after replacing \(C_i\) by an isomorphic copy, we may assume that \(B_i\subseteq C_i\). Set \(B_{i+1}:=C_i\). This completes the construction.

    Now let
    \[
        C:=\bigcup_{i<\omega} B_i .
    \]
    Since \(C\) is a union of countably many finitely generated structures, it is countable.

    By construction, every member of \(\mathcal K\) embeds into \(C\), so \(\mathcal K\subseteq \Age(C)\). Conversely, let \(A\) be a finitely generated substructure of \(C\). A finite generating set for \(A\) is contained in some \(B_i\), hence \(A\subseteq B_i\). Since \(B_i\in\mathcal K\) up to isomorphism and \(\mathcal K\) satisfies \emph{(HP)}, it follows that \(A\in\mathcal K\) up to isomorphism. Therefore \(\Age(C)\subseteq \mathcal K\).

    Hence \(\Age(C)=\mathcal K\).
\end{proof}

\begin{remark}
    If \(L\) is a finite relational signature, then every finitely generated \(L\)-structure is finite. In that case the countability assumption on \(\mathcal K\) is often easy to verify by Proposition~\ref{prop:finite-relational-language-finiteness}.
\end{remark}


\section{A motivating example: finite linear orders}

The preceding discussion shows that every age satisfies \emph{(HP)} and \emph{(JEP)}. A natural question is whether these conditions are strong enough to determine a distinguished countable structure. The example of linear orders shows that they are not.

\begin{proposition}\label{prop:age-infinite-linear-order}
    Let \(L=\{<\}\), and let \((M,<)\) be an infinite linear order. Then
    \[
        \Age(M)=\{\text{finite linear orders}\}.
    \]
    In particular, any two infinite linear orders have the same age.
\end{proposition}

\begin{proof}
    Since the language \(L=\{<\}\) is purely relational, every finitely generated substructure of \(M\) is simply the induced substructure on a finite subset of \(M\). The restriction of \(<\) to a finite subset is again a finite linear order. Thus every member of \(\Age(M)\) is a finite linear order.

    Conversely, let \((F,<)\) be a finite linear order, say with
    \[
        F=\{a_1,\dots,a_n\}, \qquad a_1<\cdots<a_n .
    \]
    Since \(M\) is infinite, we may choose distinct elements \(b_1,\dots,b_n\in M\), and after reordering assume that
    \[
        b_1<\cdots<b_n .
    \]
    Then the map \(e\colon F\to M\) given by \(e(a_i)=b_i\) is an embedding. Hence every finite linear order belongs to \(\Age(M)\), and therefore
    \[
        \Age(M)=\{\text{finite linear orders}\}.
    \]
\end{proof}

Let \(\mathcal K_{<}\) denote the class of all finite linear orders, up to isomorphism. By Proposition~\ref{prop:age-infinite-linear-order},
\[
    \Age(\mathbb Q,<)=\Age(\mathbb Z,<)=\mathcal K_{<}.
\]
Moreover, \(\mathcal K_{<}\) satisfies \emph{(HP)} and \emph{(JEP)}: \emph{(HP)} is immediate, and \emph{(JEP)} holds because any two finite linear orders embed into a sufficiently long finite linear order.

Thus age, together with \emph{(HP)} and \emph{(JEP)}, does not distinguish the countable dense linear order without endpoints from the order of the integers. In particular, these conditions do not yet explain why the finite linear orders ``tend'' to \((\mathbb Q,<)\) rather than to \((\mathbb Z,<)\).

This shows that one further condition is needed. The next section introduces the \emph{amalgamation property}, which is precisely the additional condition singled out by Fraïssé.

\section{The amalgamation property and Fraïssé classes}

The previous section shows that age, together with \emph{(HP)} and \emph{(JEP)}, does not suffice to distinguish \((\mathbb Q,<)\) from \((\mathbb Z,<)\). Fraïssé's additional requirement is the amalgamation property.

\begin{definition}[Amalgamation property]
    Let \(\mathcal K\) be a class of finitely generated \(L\)-structures. We say that \(\mathcal K\) has the \emph{amalgamation property}, or \emph{(AP)}, if whenever \(A,B,C\in\mathcal K\) and
    \[
        e\colon A\hookrightarrow B,
        \qquad
        f\colon A\hookrightarrow C
    \]
    are embeddings, there exist \(D\in\mathcal K\) and embeddings
    \[
        g\colon B\hookrightarrow D,
        \qquad
        h\colon C\hookrightarrow D
    \]
    such that
    \[
        g\circ e=h\circ f.
    \]
\end{definition}

Thus \emph{(AP)} says that any two structures in \(\mathcal K\) which share a common substructure can be embedded into a larger structure in \(\mathcal K\) so that the two copies of the common part coincide.

\[
    \begin{tikzcd}[row sep=3em, column sep=3.5em, ampersand replacement=\&]
        \& B \ar[dashed,dr,"g"] \& \\
        A \ar[ur,"e"] \ar[dr,"f"']
        \& \& D \\
        \& C \ar[dashed,ur,"h"'] \&
    \end{tikzcd}
\]

We now return to the class of finite linear orders.

\begin{proposition}\label{prop:finite-linear-orders-ap}
    The class of finite linear orders has the amalgamation property.
\end{proposition}

\begin{proof}
    We first treat the special case in which \(A\subseteq B\) and \(A\subseteq C\), the maps \(e\) and \(f\) are the inclusion maps, and
    \[
        A=B\cap C.
    \]
    Write
    \[
        A=\{a_1<\cdots<a_n\}.
    \]
    For convenience, introduce formal symbols
    \[
        a_0=-\infty,\qquad a_{n+1}=+\infty.
    \]
    For each \(0\leq i\leq n\), let
    \[
        B_i:=\{b\in B\setminus A : a_i<b<a_{i+1}\},
        \qquad
        C_i:=\{c\in C\setminus A : a_i<c<a_{i+1}\},
    \]
    where the inequalities involving \(a_0\) and \(a_{n+1}\) are interpreted in the obvious way. Since \(A\) is a substructure of both \(B\) and \(C\), every element of \(B\setminus A\) belongs to exactly one \(B_i\), and every element of \(C\setminus A\) belongs to exactly one \(C_i\).

    Let
    \[
        D:=B\cup C
    \]
    as a set. We define a linear order \(<_{D}\) on \(D\) as follows:
    \begin{itemize}
        \item \(<_{D}\) agrees with the given orders on \(B\) and on \(C\);
        \item for each \(i\), every element of \(B_i\) is declared smaller than every element of \(C_i\);
        \item if \(x\in B_i\cup C_i\) and \(y\in B_j\cup C_j\) with \(i<j\), then \(x<_{D}y\);
        \item if \(x\in B_i\cup C_i\), then
              \[
                  a_i <_{D} x <_{D} a_{i+1},
              \]
              with the obvious interpretation when \(i=0\) or \(i=n\).
    \end{itemize}

    In other words, inside each interval determined by the elements of \(A\), we place first the elements coming from \(B\), in their original order, and then the elements coming from \(C\), in their original order.

    It is straightforward to verify that \(<_{D}\) is a linear order on \(D\). By construction, the inclusion maps
    \[
        j_B\colon B\hookrightarrow D,
        \qquad
        j_C\colon C\hookrightarrow D
    \]
    are embeddings of linear orders, and they agree on \(A\). Thus \(D\) is an amalgam of \(B\) and \(C\) over \(A\).
    We now consider the general case. Let
    \[
        e\colon A\hookrightarrow B,
        \qquad
        f\colon A\hookrightarrow C
    \]
    be embeddings of finite linear orders. By Lemma~\ref{lem:embedding-as-inclusion}, there exist finite linear orders \(B_1\) and \(C_1\), together with isomorphisms
    \[
        \beta\colon B_1\to B,
        \qquad
        \gamma\colon C_1\to C,
    \]
    such that \(A\subseteq B_1\), \(A\subseteq C_1\), and
    \[
        \beta\!\restriction_A=e,
        \qquad
        \gamma\!\restriction_A=f.
    \]
    Replacing \(B_1\) and \(C_1\) by isomorphic copies if necessary, we may moreover assume that
    \[
        B_1\cap C_1=A.
    \]
    By the special case already proved, there exist a finite linear order \(D_1\) and embeddings
    \[
        j_B\colon B_1\hookrightarrow D_1,
        \qquad
        j_C\colon C_1\hookrightarrow D_1
    \]
    such that
    \[
        j_B\!\restriction_A=j_C\!\restriction_A.
    \]
    Define
    \[
        g:=j_B\circ \beta^{-1}\colon B\hookrightarrow D_1,
        \qquad
        h:=j_C\circ \gamma^{-1}\colon C\hookrightarrow D_1.
    \]
    Then
    \[
        g\circ e
        =
        j_B\circ \beta^{-1}\circ e
        =
        j_B\!\restriction_A
        =
        j_C\!\restriction_A
        =
        j_C\circ \gamma^{-1}\circ f
        =
        h\circ f,
    \]
    so \(D_1\) is an amalgam of \(e\) and \(f\).

    \[
        \begin{tikzcd}[column sep=large, row sep=large, ampersand replacement=\&]
            \& B_1 \arrow[rr, "\beta", "\cong"'] \arrow[dr, "j_B"'] \& \& B \arrow[dl, dashed, "g:=j_B\circ\beta^{-1}"] \\
            A \arrow[ur, hook] \arrow[dr, hook'] \arrow[urrr, bend left=5, "e"] \arrow[drrr, bend right=5, "f"'] \& \& D_1 \& \\
            \& C_1 \arrow[rr, "\gamma"', "\cong"] \arrow[ur, "j_C"] \& \& C \arrow[ul, dashed, "h:=j_C\circ\gamma^{-1}"']
        \end{tikzcd}
    \]

    Hence the class of finite linear orders has \emph{(AP)}.

\end{proof}

This example leads to the central class of finite structures considered in Fraïssé theory.

\begin{definition}[Fraïssé class]
    Let \(L\) be a countable signature. A class \(\mathcal K\) of finitely generated \(L\)-structures is called a \emph{Fraïssé class} if it satisfies the following conditions:
    \begin{enumerate}
        \item \(\mathcal K\) is closed under isomorphism;
        \item \(\mathcal K\) has countably many isomorphism types;
        \item \(\mathcal K\) has the hereditary property \emph{(HP)};
        \item \(\mathcal K\) has the joint embedding property \emph{(JEP)};
        \item \(\mathcal K\) has the amalgamation property \emph{(AP)}.
    \end{enumerate}
\end{definition}


\section{Fraïssé limits}

We now arrive at the central notion of the chapter.

\begin{definition}[Ultrahomogeneous structure]
    Let \(D\) be an \(L\)-structure. We say that \(D\) is \emph{ultrahomogeneous} if every isomorphism between finitely generated substructures of \(D\) extends to an automorphism of \(D\).
\end{definition}

Fraïssé's theorem shows that every Fraïssé class determines, up to isomorphism, a unique countable ultrahomogeneous structure whose age is exactly that class.

\begin{theorem}[Fraïssé's theorem]\label{thm:fraisse-theorem}
    Let \(\mathcal K\) be a Fraïssé class. Then there exists a countable ultrahomogeneous \(L\)-structure \(D\) such that
    \[
        \Age(D)=\mathcal K.
    \]
    Moreover, \(D\) is unique up to isomorphism.
\end{theorem}

\begin{definition}[Fraïssé limit]
    Let \(\mathcal K\) be a Fraïssé class. The unique countable ultrahomogeneous structure \(D\) satisfying
    \[
        \Age(D)=\mathcal K
    \]
    is called the \emph{Fraïssé limit} of \(\mathcal K\).
\end{definition}

The next two sections are devoted to the proof of Theorem~\ref{thm:fraisse-theorem}: first the uniqueness of the Fraïssé limit, and then its existence.

\section{Uniqueness of the Fraïssé limit}

We begin with a weaker extension property which is sufficient, in the countable case, to determine a structure uniquely from its age.

\begin{definition}[Weak homogeneity]
    Let \(D\) be an \(L\)-structure. We say that \(D\) is \emph{weakly homogeneous} if whenever \(A\subseteq B\) are finitely generated substructures of \(D\) and
    \[
        f\colon A\hookrightarrow D
    \]
    is an embedding, there exists an embedding
    \[
        g\colon B\hookrightarrow D
    \]
    extending \(f\).
\end{definition}

\[
    \begin{tikzcd}[ampersand replacement=\&]
        A \arrow[r,"f"] \arrow[d,hook] \& D \\
        B \arrow[ur,dashed,"g"']
    \end{tikzcd}
\]

Weak homogeneity is strictly weaker than ultrahomogeneity as a formal condition, but for finite or countable structures the two notions coincide.

\begin{example}
    The ordered set \((\mathbb Z,<)\) is not weakly homogeneous.
\end{example}

\begin{proof}
    Let
    \[
        A=\{0,2\}\subseteq \mathbb Z,
        \qquad
        B=\{0,1,2\}\subseteq \mathbb Z,
    \]
    with the induced order, and define \(f\colon A\to \mathbb Z\) by
    \[
        f(0)=0,
        \qquad
        f(2)=1.
    \]
    Then \(f\) is an embedding. If \(g\colon B\to \mathbb Z\) were an embedding extending \(f\), we would have
    \[
        0=g(0)<g(1)<g(2)=1,
    \]
    which is impossible in \(\mathbb Z\). Hence no such extension exists.
\end{proof}

\begin{lemma}\label{lem:ultra-implies-weak}
    Every ultrahomogeneous structure is weakly homogeneous.
\end{lemma}

\begin{proof}
    Let \(A\subseteq B\) be finitely generated substructures of an ultrahomogeneous structure \(D\), and let \(f\colon A\hookrightarrow D\) be an embedding. Then \(f\colon A\to f(A)\) is an isomorphism between finitely generated substructures of \(D\). By ultrahomogeneity, \(f\) extends to an automorphism \(\alpha\) of \(D\). The restriction
    \[
        g:=\alpha\!\restriction_B\colon B\hookrightarrow D
    \]
    is then an embedding extending \(f\).
\end{proof}

The next lemma shows that a countable weakly homogeneous structure is universal for its age.

\begin{lemma}\label{lem:weak-age-embedding}
    Let \(C\) and \(D\) be at most countable \(L\)-structures. Suppose that
    \[
        \mathrm{Age}(C)\subseteq \mathrm{Age}(D)
    \]
    and that \(D\) is weakly homogeneous. Then every embedding of a finitely generated substructure of \(C\) into \(D\) extends to an embedding of \(C\) into \(D\).
\end{lemma}

\begin{proof}
    Let \(f_0\colon A_0\hookrightarrow D\) be an embedding, where \(A_0\subseteq C\) is finitely generated. Choose an increasing chain of finitely generated substructures
    \[
        A_0\subseteq A_1\subseteq A_2\subseteq \cdots
    \]
    with
    \[
        C=\bigcup_{n<\omega} A_n .
    \]
    For example, if \(C\setminus A_0=\{c_0,c_1,\dots\}\), we may take
    \[
        A_{n+1}:=\langle A_n\cup\{c_n\}\rangle .
    \]

    We construct inductively embeddings
    \[
        f_n\colon A_n\hookrightarrow D
        \qquad (n<\omega)
    \]
    such that \(f_n\subseteq f_{n+1}\). Suppose \(f_n\) has been constructed. Since \(A_{n+1}\) is finitely generated and \(\mathrm{Age}(C)\subseteq \mathrm{Age}(D)\), there exist a finitely generated substructure \(B\subseteq D\) and an isomorphism
    \[
        u\colon A_{n+1}\xrightarrow{\ \cong\ } B.
    \]
    Then
    \[
        f_n\circ (u\!\restriction_{A_n})^{-1}\colon u(A_n)\hookrightarrow D
    \]
    is an embedding. Since \(u(A_n)\subseteq B\) and \(D\) is weakly homogeneous, this embedding extends to an embedding
    \[
        h\colon B\hookrightarrow D.
    \]
    Define
    \[
        f_{n+1}:=h\circ u\colon A_{n+1}\hookrightarrow D.
    \]
    For \(a\in A_n\),
    \[
        f_{n+1}(a)=h(u(a))
        =\bigl(f_n\circ (u\!\restriction_{A_n})^{-1}\bigr)(u(a))
        =f_n(a),
    \]
    so \(f_{n+1}\) extends \(f_n\).

    Now let
    \[
        f:=\bigcup_{n<\omega} f_n\colon C\to D.
    \]
    This is well defined because the maps \(f_n\) are compatible. Since every finite tuple from \(C\) is contained in some \(A_n\), and each \(f_n\) is an embedding, it follows that \(f\) is an embedding of \(C\) into \(D\). By construction, \(f\) extends \(f_0\).
\end{proof}

We now obtain the uniqueness statement by the usual back-and-forth argument.

\begin{lemma}\label{lem:countable-weak-uniqueness}
    Let \(C\) and \(D\) be at most countable weakly homogeneous \(L\)-structures with
    \[
        \mathrm{Age}(C)=\mathrm{Age}(D).
    \]
    Then \(C\cong D\). More precisely, every isomorphism between finitely generated substructures of \(C\) and \(D\) extends to an isomorphism from \(C\) onto \(D\).
\end{lemma}

\begin{proof}
    Let
    \[
        f_0\colon A_0\xrightarrow{\ \cong\ } B_0
    \]
    be an isomorphism between finitely generated substructures \(A_0\subseteq C\) and \(B_0\subseteq D\). Choose increasing chains of finitely generated substructures
    \[
        A_0\subseteq A_1\subseteq A_2\subseteq \cdots,
        \qquad
        B_0\subseteq B_1\subseteq B_2\subseteq \cdots
    \]
    such that
    \[
        C=\bigcup_{n<\omega} A_n,
        \qquad
        D=\bigcup_{n<\omega} B_n .
    \]

    We construct inductively a chain of partial isomorphisms
    \[
        f_0\subseteq f_1\subseteq f_2\subseteq \cdots
    \]
    between finitely generated substructures of \(C\) and \(D\) such that for every \(n\),
    \[
        A_n\subseteq \dom(f_{2n}),
        \qquad
        B_n\subseteq \operatorname{im}(f_{2n+1}).
    \]

    Suppose \(f_{2n}\) has been constructed. Since \(D\) is weakly homogeneous and \(\mathrm{Age}(C)\subseteq \mathrm{Age}(D)\), Lemma~\ref{lem:weak-age-embedding} extends \(f_{2n}\) to an embedding
    \[
        g\colon C\hookrightarrow D.
    \]
    Let
    \[
        f_{2n+1}:=g\!\restriction_{\langle \dom(f_{2n})\cup A_{n+1}\rangle}.
    \]
    Then \(f_{2n+1}\) extends \(f_{2n}\), and its domain contains \(A_{n+1}\).

    Now \(f_{2n+1}^{-1}\) is an isomorphism from the finitely generated substructure \(\operatorname{im}(f_{2n+1})\subseteq D\) into \(C\). Since \(C\) is weakly homogeneous and \(\mathrm{Age}(D)\subseteq \mathrm{Age}(C)\), Lemma~\ref{lem:weak-age-embedding} extends \(f_{2n+1}^{-1}\) to an embedding
    \[
        h\colon D\hookrightarrow C.
    \]
    Define
    \[
        f_{2n+2}:=\bigl(h\!\restriction_{\langle \operatorname{im}(f_{2n+1})\cup B_{n+1}\rangle}\bigr)^{-1}.
    \]
    Then \(f_{2n+2}\) extends \(f_{2n+1}\), and its image contains \(B_{n+1}\).

    Thus the construction continues indefinitely. Let
    \[
        f:=\bigcup_{n<\omega} f_n.
    \]
    Because the maps \(f_n\) are compatible, \(f\) is a well-defined embedding \(C\to D\). The domain condition implies that every element of \(C\) lies in \(\dom(f)\), and the image condition implies that every element of \(D\) lies in \(\operatorname{im}(f)\). Hence \(f\) is an isomorphism from \(C\) onto \(D\). By construction, it extends \(f_0\).
\end{proof}

\begin{corollary}\label{cor:weak-equals-ultra-countable}
    For a finite or countable \(L\)-structure \(D\), the following are equivalent:
    \begin{enumerate}
        \item \(D\) is ultrahomogeneous;
        \item \(D\) is weakly homogeneous.
    \end{enumerate}
\end{corollary}

\begin{proof}
    By Lemma~\ref{lem:ultra-implies-weak}, every ultrahomogeneous structure is weakly homogeneous.

    Conversely, suppose that \(D\) is finite or countable and weakly homogeneous. Let
    \[
        f\colon A\xrightarrow{\ \cong\ } B
    \]
    be an isomorphism between finitely generated substructures of \(D\). Applying Lemma~\ref{lem:countable-weak-uniqueness} with \(C=D\), \(D=D\), and the given partial isomorphism \(f\), we obtain an automorphism of \(D\) extending \(f\). Hence \(D\) is ultrahomogeneous.
\end{proof}

In particular, if \(C\) and \(D\) are countable ultrahomogeneous structures with the same age, then \(C\cong D\). This proves the uniqueness part of Fraïssé's theorem.



\section{Existence of the Fraïssé limit}

We now prove the existence part of Fraïssé's theorem. The idea is to construct a countable structure \(D\) with age \(\mathcal K\) and with the extension property that characterizes weak homogeneity.

We first record a simple lemma on increasing unions.

\begin{lemma}\label{lem:age-of-union}
    Let
    \[
        D_0\subseteq D_1\subseteq D_2\subseteq \cdots
    \]
    be an increasing chain of \(L\)-structures, and let
    \[
        D:=\bigcup_{i<\omega} D_i.
    \]
    Let \(\mathcal K\) be a class of finitely generated \(L\)-structures.
    \begin{enumerate}
        \item[(i)] If \(\mathrm{Age}(D_i)\subseteq \mathcal K\) for all \(i<\omega\), then
              \[
                  \mathrm{Age}(D)\subseteq \mathcal K.
              \]
        \item[(ii)] If, in addition, every \(A\in\mathcal K\) embeds into some \(D_i\), then
              \[
                  \mathrm{Age}(D)=\mathcal K.
              \]
    \end{enumerate}
\end{lemma}

\begin{proof}
    Let \(A\) be a finitely generated substructure of \(D\), generated by a finite tuple \(\bar a\). Since the chain is increasing, all entries of \(\bar a\) lie in some \(D_i\). Hence
    \[
        A=\langle \bar a\rangle_D=\langle \bar a\rangle_{D_i},
    \]
    so \(A\in\mathrm{Age}(D_i)\subseteq\mathcal K\). This proves (i).

    For (ii), part (i) already gives \(\mathrm{Age}(D)\subseteq\mathcal K\). Conversely, let \(A\in\mathcal K\). By assumption, \(A\) embeds into some \(D_i\), and therefore also into \(D\). Thus \(A\in\mathrm{Age}(D)\), so \(\mathcal K\subseteq\mathrm{Age}(D)\).
\end{proof}

\begin{proof}[Proof of Fraïssé's theorem]
    Let \(\mathcal K\) be a Fraïssé class. Since only isomorphism types matter, we may assume that \(\mathcal K\) is closed under isomorphism.

    We shall construct an increasing chain
    \[
        D_0\subseteq D_1\subseteq D_2\subseteq \cdots
    \]
    of members of \(\mathcal K\) such that the following extension property holds:
    \begin{quote}
        whenever \(A,B\in\mathcal K\) with \(A\subseteq B\), and \(f\colon A\hookrightarrow D_i\) is an embedding, there exists \(j>i\) and an embedding
        \[
            g\colon B\hookrightarrow D_j
        \]
        extending \(f\).
    \end{quote}

    Assume for the moment that such a chain has been constructed, and let
    \[
        D:=\bigcup_{i<\omega} D_i.
    \]

    Since \(\mathcal K\) has countably many isomorphism types and each member of \(\mathcal K\) is finitely generated in a countable signature, each \(D_i\) is countable. Hence \(D\) is countable.

    Because each \(D_i\) belongs to \(\mathcal K\), Lemma~\ref{lem:age-of-union}(i) gives
    \[
        \mathrm{Age}(D)\subseteq \mathcal K.
    \]

    We claim that in fact \(\mathrm{Age}(D)=\mathcal K\). Let \(A\in\mathcal K\). By \emph{(JEP)}, there exist \(B\in\mathcal K\) and embeddings
    \[
        u\colon D_0\hookrightarrow B,
        \qquad
        v\colon A\hookrightarrow B.
    \]
    By Lemma~\ref{lem:embedding-as-inclusion}, after replacing \(B\) by an isomorphic copy, we may assume that
    \[
        D_0\subseteq B.
    \]
    Now apply the extension property to the identity embedding \(1_{D_0}\colon D_0\hookrightarrow D_0\). Then for some \(j>0\), the inclusion \(D_0\subseteq B\) extends to an embedding
    \[
        g\colon B\hookrightarrow D_j.
    \]
    Composing with \(v\), we obtain an embedding of \(A\) into \(D_j\). Thus every \(A\in\mathcal K\) embeds into some \(D_j\), and Lemma~\ref{lem:age-of-union}(ii) yields
    \[
        \mathrm{Age}(D)=\mathcal K.
    \]

    We next show that \(D\) is weakly homogeneous. Let \(A\subseteq B\) be finitely generated substructures of \(D\), and let
    \[
        f\colon A\hookrightarrow D
    \]
    be an embedding. Since \(A\) is finitely generated and \(f(A)\) is finite or countable, there exists \(i<\omega\) such that
    \[
        f(A)\subseteq D_i.
    \]
    Because \(\mathrm{Age}(D)=\mathcal K\), both \(A\) and \(B\) belong to \(\mathcal K\). Applying the extension property to the embedding
    \[
        f\colon A\hookrightarrow D_i,
    \]
    we obtain \(j>i\) and an embedding
    \[
        g\colon B\hookrightarrow D_j\subseteq D
    \]
    extending \(f\). Hence \(D\) is weakly homogeneous.

    Since \(D\) is countable, Corollary~\ref{cor:weak-equals-ultra-countable} implies that \(D\) is ultrahomogeneous. Therefore \(D\) is a countable ultrahomogeneous structure with age \(\mathcal K\), as required.

    It remains to construct the chain \((D_i)_{i<\omega}\).

    Let
    \[
        \mathcal P=\{(A_n,B_n): n<\omega\}
    \]
    be a countable set of representatives of the isomorphism types of pairs \((A,B)\) such that \(A,B\in\mathcal K\) and \(A\subseteq B\). Fix a bijection
    \[
        \pi\colon \omega\times\omega\to\omega
    \]
    such that \(\pi(i,j)\geq i\) for all \(i,j<\omega\). Choose any \(D_0\in\mathcal K\).

    Suppose \(D_i\) has been constructed. Since \(D_i\) is countable and each \(A_n\) is finitely generated, there are only countably many embeddings of the various \(A_n\) into \(D_i\). Enumerate all triples
    \[
        (f_{ij},A_{ij},B_{ij}) \qquad (j<\omega),
    \]
    where \(A_{ij}\subseteq B_{ij}\) is one of the pairs from \(\mathcal P\) and
    \[
        f_{ij}\colon A_{ij}\hookrightarrow D_i
    \]
    is an embedding.

    Now suppose \(D_k\) has been defined, and write
    \[
        k=\pi(i,j).
    \]
    Then
    \[
        f_{ij}\colon A_{ij}\hookrightarrow D_i\subseteq D_k
    \]
    is an embedding. Apply \emph{(AP)} to the inclusion
    \[
        A_{ij}\hookrightarrow B_{ij}
    \]
    and the embedding
    \[
        f_{ij}\colon A_{ij}\hookrightarrow D_k.
    \]
    We obtain some \(E\in\mathcal K\) and embeddings
    \[
        r\colon B_{ij}\hookrightarrow E,
        \qquad
        s\colon D_k\hookrightarrow E
    \]
    such that
    \[
        r\!\restriction_{A_{ij}}=s\circ f_{ij}.
    \]
    By Lemma~\ref{lem:embedding-as-inclusion}, after replacing \(E\) by an isomorphic copy, we may assume that
    \[
        D_k\subseteq E
    \]
    and that \(s\) is the inclusion map. Set
    \[
        D_{k+1}:=E.
    \]
    Then the embedding \(r\colon B_{ij}\hookrightarrow D_{k+1}\) extends \(f_{ij}\).

    This completes the construction of the chain. Because every embedding problem over every stage \(D_i\) appears as some triple \((f_{ij},A_{ij},B_{ij})\), and because \(\pi(i,j)\geq i\), each such problem is eventually handled at a later stage. Hence the required extension property holds.
\end{proof}


\section{The converse direction}

We now show that the conditions in Fraïssé's theorem are also necessary.

\begin{theorem}\label{thm:converse-fraisse}
    Let \(L\) be a countable signature, and let \(D\) be a finite or countable ultrahomogeneous \(L\)-structure. Then \(\mathrm{Age}(D)\) is a Fraïssé class. More explicitly, \(\mathrm{Age}(D)\) is nonempty, has countably many isomorphism types, and satisfies \emph{(HP)}, \emph{(JEP)}, and \emph{(AP)}.
\end{theorem}

\begin{proof}
    Set
    \[
        \mathcal K:=\mathrm{Age}(D).
    \]
    By Proposition~\ref{prop:age-hp-jep}, the class \(\mathcal K\) is nonempty and satisfies \emph{(HP)} and \emph{(JEP)}.

    It remains to verify that \(\mathcal K\) has countably many isomorphism types and satisfies \emph{(AP)}.

    Since \(D\) is finite or countable, there are only countably many finite tuples from \(D\). Every member of \(\mathcal K\) is finitely generated, hence is generated by some finite tuple from \(D\). Therefore \(\mathcal K\) has at most countably many isomorphism types.

    To prove \emph{(AP)}, let \(A,B,C\in\mathcal K\), and let
    \[
        e\colon A\hookrightarrow B,
        \qquad
        f\colon A\hookrightarrow C
    \]
    be embeddings. Since \(A\in\mathcal K=\mathrm{Age}(D)\), there exists an embedding
    \[
        u\colon A\hookrightarrow D.
    \]
    By Corollary~\ref{cor:weak-equals-ultra-countable}, the structure \(D\) is weakly homogeneous. Now
    \[
        u\circ e^{-1}\colon e(A)\hookrightarrow D
    \]
    is an embedding of the finitely generated substructure \(e(A)\subseteq B\) into \(D\). Since \(B\in\mathcal K\), weak homogeneity yields an embedding
    \[
        j_B\colon B\hookrightarrow D
    \]
    extending \(u\circ e^{-1}\). Thus
    \[
        j_B\circ e=u.
    \]
    Similarly, the embedding
    \[
        u\circ f^{-1}\colon f(A)\hookrightarrow D
    \]
    extends to an embedding
    \[
        j_C\colon C\hookrightarrow D
    \]
    with
    \[
        j_C\circ f=u.
    \]
    Hence
    \[
        j_B\circ e=j_C\circ f.
    \]

    Let
    \[
        E:=\langle j_B(B)\cup j_C(C)\rangle \subseteq D.
    \]
    Since \(B\) and \(C\) are finitely generated, the structure \(E\) is finitely generated, so \(E\in\mathcal K\). The maps \(j_B\colon B\hookrightarrow E\) and \(j_C\colon C\hookrightarrow E\) are embeddings and satisfy
    \[
        j_B\circ e=j_C\circ f,
    \]
    so \(E\) is an amalgam of \(B\) and \(C\) over \(A\). Therefore \(\mathcal K\) satisfies \emph{(AP)}.

    Thus \(\mathcal K\) is a Fraïssé class.
\end{proof}
% !TEX root = main.tex
\chapter{Automorphism Groups}
\label{chap:automorphism-groups}

Fraïssé theory leads naturally from countable structures to their automorphism groups. In this chapter we study automorphism groups of first-order structures as permutation groups on the underlying domain and equip them with the topology of pointwise convergence. We begin with the basic notions of stabilizers, orbits, transitivity, and rigidity, and then introduce the permutation topology on $\Sym(\Omega)$ and on its subgroups.

We next characterize those subgroups of $\Sym(\Omega)$ which arise as automorphism groups of structures on $\Omega$, in terms of closed subgroups and canonical structures. In the countable case, this shows that automorphism groups are closed subgroups of $S_\infty$, and hence naturally carry the structure of Polish groups. These results provide the topological framework needed in the following chapters, where we study continuous actions of such groups on compact spaces.

\section{Automorphism Groups as Permutation Groups}

Let $A$ be an $L$-structure with domain $\Omega=\dom(A)$. Every automorphism of $A$ is a permutation of $\Omega$, and the set of all automorphisms of $A$ forms a group under composition. We denote this group by $\Aut(A)$ and regard it as a subgroup of the symmetric group $\Sym(\Omega)$.

More generally, if $\Omega$ is any set, we write $\Sym(\Omega)$ for the group of all permutations of $\Omega$. Thus every automorphism group is naturally a permutation group. Our aim in this chapter is to study those subgroups of $\Sym(\Omega)$ which arise as automorphism groups of structures on $\Omega$, and to equip them with the topology of pointwise convergence.

Let $G\leq\Sym(\Omega)$. For a subset $X\subseteq\Omega$, the \emph{pointwise stabilizer} of $X$ in $G$ is \(G_{(X)}=\{g\in G:g(a)=a\text{ for all }a\in X\}\), and the \emph{setwise stabilizer} of $X$ in $G$ is \(G_{\{X\}}=\{g\in G:g(X)=X\}\).
Thus $G_{(X)}\leq G_{\{X\}}\leq G$. If $\bar a=(a_1,\dots,a_n)$ is a finite tuple from $\Omega$, we also write $G_{(\bar a)}$ for the pointwise stabilizer of the set $\{a_1,\dots,a_n\}$.

For $a\in\Omega$, the \emph{orbit} of $a$ under $G$ is the set \(G\cdot a=\{g(a):g\in G\}\). More generally, if $\bar a=(a_1,\dots,a_n)\in\Omega^n$, the \emph{orbit} of $\bar a$ under $G$ is \(G\cdot \bar a=\{g\bar a:g\in G\}\), where \(g\bar a=(g(a_1),\dots,g(a_n))\).
The orbits of $G$ on $\Omega$ partition $\Omega$. We say that $G$ is \emph{transitive} if it has exactly one orbit on $\Omega$.

Accordingly, a structure $A$ is called \emph{transitive} if $\Aut(A)$ acts transitively on $\dom(A)$. At the opposite extreme, $A$ is called \emph{rigid} if $\Aut(A)=\{1_{\Omega}\}$.

\begin{example}
    The following examples will be useful throughout the thesis.
    \begin{enumerate}[(i)]
        \item If $A$ is a pure set with no non-logical symbols, then every permutation of $\dom(A)$ is an automorphism, so $\Aut(A)=\Sym(\dom(A))$.
        \item Every finite linear order is rigid. Indeed, if $A=(\{a_1,\dots,a_n\},<)$ with $a_1<\cdots<a_n$, then any automorphism must preserve the least element, then the next least element, and so on; hence it fixes every point.
        \item The ordered set $(\mathbb Q,<)$ is transitive. Given $q,r\in\mathbb Q$, there exists an order automorphism of $(\mathbb Q,<)$ sending $q$ to $r$, so all elements lie in a single orbit under $\Aut(\mathbb Q,<)$.
    \end{enumerate}
\end{example}

These notions will reappear throughout the chapter. In particular, stabilizers of finite tuples will provide the basic open neighbourhoods in the permutation topology, while orbits on finite powers $\Omega^n$ will play a central role in the description of closed and dense subgroups.

\section{The Permutation Topology}

We now introduce the natural topology on permutation groups. Let $\Omega$ be a set. For $n<\omega$ and tuples $\bar a,\bar b\in\Omega^n$, set \(S(\bar a,\bar b)=\{g\in\Sym(\Omega):g\bar a=\bar b\}\).
A subset of $\Sym(\Omega)$ is called \emph{basic open} if it is of the form $S(\bar a,\bar b)$ for some finite tuples $\bar a,\bar b$ of the same length. An \emph{open} subset of $\Sym(\Omega)$ is any union of basic open sets. This is the \emph{permutation topology}, or equivalently the topology of pointwise convergence.

If $G\leq\Sym(\Omega)$, we always equip $G$ with the subspace topology induced from $\Sym(\Omega)$. Thus the basic open subsets of $G$ are precisely the sets \(G\cap S(\bar a,\bar b)\), where $\bar a,\bar b$ are finite tuples from $\Omega$ of the same length.

\begin{proposition}
    \label{prop:permutation-topology}
    Let $G\leq\Sym(\Omega)$.
    \begin{enumerate}[(a)]
        \item The sets $S(\bar a,\bar b)$ form a basis for a topology on $\Sym(\Omega)$. With this topology, $\Sym(\Omega)$ is a topological group, and hence so is every subgroup $G\leq\Sym(\Omega)$ with the induced topology.
        \item A subgroup $H\leq G$ is open if and only if it contains $G_{(\bar a)}$ for some finite tuple $\bar a$ from $\Omega$.
        \item A subset $F\subseteq G$ is closed if and only if whenever $g\in G$ has the property that every basic open neighbourhood of $g$ meets $F$, then $g\in F$.
        \item A subgroup $H\leq G$ is dense in $G$ if and only if $H$ and $G$ have the same orbits on $\Omega^n$ for every positive integer $n$.
    \end{enumerate}
\end{proposition}

\begin{proof}
    For (a), note first that \(\Sym(\Omega)=S(\emptyset,\emptyset)\).
    Moreover, if $g\in S(\bar a,\bar b)\cap S(\bar c,\bar d)$, then
    \[
        g\in S(\bar a\bar c,\bar b\bar d)\subseteq S(\bar a,\bar b)\cap S(\bar c,\bar d),
    \]
    where $\bar a\bar c$ and $\bar b\bar d$ denote concatenation of tuples. Thus the sets $S(\bar a,\bar b)$ form a basis for a topology on $\Sym(\Omega)$.

    To prove continuity of inversion, observe that \(g\in S(\bar a,\bar b)\Longleftrightarrow g^{-1}\in S(\bar b,\bar a)\). For continuity of multiplication, let $(g,h)\in\Sym(\Omega)\times\Sym(\Omega)$ and suppose that \(gh\in S(\bar a,\bar b)\).
    Set $\bar c=h\bar a$. Then
    \[
        h\in S(\bar a,\bar c)\qquad\text{and}\qquad g\in S(\bar c,\bar b),
    \]
    and whenever $g'\in S(\bar c,\bar b)$ and $h'\in S(\bar a,\bar c)$, we have \(g'h'\bar a=g'\bar c=\bar b\), so that \(g'h'\in S(\bar a,\bar b)\).
    Hence
    \[
        S(\bar c,\bar b)\cdot S(\bar a,\bar c)\subseteq S(\bar a,\bar b),
    \]
    which proves continuity of multiplication. Therefore $\Sym(\Omega)$ is a topological group. Since subgroups inherit the subspace topology, every subgroup $G\leq\Sym(\Omega)$ is also a topological group.

    For (b), if $\bar a$ is a finite tuple from $\Omega$, then \(G_{(\bar a)}=G\cap S(\bar a,\bar a)\),
    so $G_{(\bar a)}$ is open in $G$. If $H\leq G$ contains $G_{(\bar a)}$, then
    \[
        H=\bigcup_{h\in H} hG_{(\bar a)},
    \]
    and each left coset $hG_{(\bar a)}$ is open, since left translation is a homeomorphism of $G$. Thus $H$ is open.

    Conversely, suppose that $H\leq G$ is open. Since $1_G\in H$, there exists a basic open neighbourhood of $1_G$ contained in $H$. Such a neighbourhood has the form \(G\cap S(\bar a,\bar a)=G_{(\bar a)}\) for some finite tuple $\bar a$ from $\Omega$. Hence $G_{(\bar a)}\subseteq H$.

    Statement (c) is just the usual characterization of closed sets in terms of neighbourhoods: a subset $F\subseteq G$ is closed if and only if every point outside $F$ has an open neighbourhood disjoint from $F$, equivalently, whenever every basic open neighbourhood of $g$ meets $F$, one has $g\in F$.

    For (d), first suppose that $H$ is dense in $G$, and let $\bar a,\bar b\in\Omega^n$ lie in the same $G$-orbit. Choose $g\in G$ with $g\bar a=\bar b$. Then $S(\bar a,\bar b)$ is a basic open neighbourhood of $g$, so by density it meets $H$. Hence there exists $h\in H$ such that $h\bar a=\bar b$. Thus $\bar a$ and $\bar b$ lie in the same $H$-orbit. It follows that every $G$-orbit on $\Omega^n$ is also an $H$-orbit, and therefore $H$ and $G$ have the same orbits on $\Omega^n$.

    Conversely, suppose that $H$ and $G$ have the same orbits on $\Omega^n$ for every $n\geq 1$. To show that $H$ is dense in $G$, it suffices by part (c) to show that every nonempty basic open subset of $G$ meets $H$. Let \(G\cap S(\bar a,\bar b)\) be a nonempty basic open subset of $G$. Then there exists $g\in G$ such that $g\bar a=\bar b$, so $\bar a$ and $\bar b$ lie in the same $G$-orbit on $\Omega^n$. By assumption they also lie in the same $H$-orbit, so there exists $h\in H$ such that $h\bar a=\bar b$. Hence \(h\in H\cap S(\bar a,\bar b)\subseteq H\cap G\cap S(\bar a,\bar b)\), and therefore $G\cap S(\bar a,\bar b)$ meets $H$. This proves that $H$ is dense in $G$.
\end{proof}

In particular, the stabilizers of finite tuples form a neighbourhood basis at the identity in any subgroup of $\Sym(\Omega)$, and orbit structure on finite powers $\Omega^n$ detects density. These two observations will be used repeatedly in the next section.

\section{Closed Subgroups and Canonical Structures}

We now characterize those subgroups of $\Sym(\Omega)$ which arise as automorphism groups of structures on $\Omega$. In view of the preceding section, the relevant notion of closedness is topological closedness in the permutation topology.

\begin{theorem}
    \label{thm:closed-subgroups-automorphism-groups}
    Let $\Omega$ be a set, let $G\leq \Sym(\Omega)$, and let $H\leq G$. Then the following are equivalent.
    \begin{enumerate}[(a)]
        \item $H$ is closed in $G$.
        \item There exists a structure $A$ with domain $\Omega$ such that
              \[
                  H=G\cap \Aut(A).
              \]
    \end{enumerate}
    In particular, a subgroup $H\leq \Sym(\Omega)$ is of the form $\Aut(A)$ for some structure $A$ with domain $\Omega$ if and only if $H$ is closed in $\Sym(\Omega)$.
\end{theorem}

\begin{proof}
    Assume first that $H$ is closed in $G$. For each $n<\omega$ and each orbit $\Delta$ of $H$ on $\Omega^n$, introduce an $n$-ary relation symbol $R_\Delta$. Let $L$ be the language consisting of all these relation symbols, and define an $L$-structure $A$ with domain $\Omega$ by setting \(R_\Delta^A=\Delta\) for every orbit $\Delta$ of $H$ on $\Omega^n$.

    Every element of $H$ preserves each relation $R_\Delta^A$, since each $R_\Delta^A$ is an $H$-orbit. Hence \(H\subseteq G\cap \Aut(A)\).

    Conversely, let $g\in G\cap \Aut(A)$. We claim that every basic open neighbourhood of $g$ in $G$ meets $H$. Let $\bar a\in\Omega^n$, and let $\Delta$ be the $H$-orbit of $\bar a$. Since $\bar a\in \Delta=R_\Delta^A$ and $g$ is an automorphism of $A$, we have
    \[
        g\bar a\in R_\Delta^A=\Delta.
    \]
    Thus there exists $h\in H$ such that \(h\bar a=g\bar a\). It follows that \(h\in H\cap G\cap S(\bar a,g\bar a)\),
    so the basic open neighbourhood $G\cap S(\bar a,g\bar a)$ of $g$ meets $H$. Since $H$ is closed in $G$, we conclude that $g\in H$. Therefore \(G\cap \Aut(A)\subseteq H\), and hence \(H=G\cap \Aut(A)\).

    Now assume that there exists a structure $A$ with domain $\Omega$ such that \(H=G\cap \Aut(A)\).
    To show that $H$ is closed in $G$, let $g\in G$ and suppose that every basic open neighbourhood of $g$ in $G$ meets $H$. We show that $g\in \Aut(A)$.

    Let $\varphi(\bar x)$ be an atomic formula in the language of $A$, and let $\bar a$ be a tuple from $\Omega$ of the same length as $\bar x$. By assumption, the basic open neighbourhood \(G\cap S(\bar a,g\bar a)\) meets $H$. Choose \(h\in H\cap G\cap S(\bar a,g\bar a)\).
    Then $h\bar a=g\bar a$. Since $h\in H\subseteq \Aut(A)$, we have
    \[
        A\models \varphi(\bar a)\iff A\models \varphi(h\bar a).
    \]
    Since $h\bar a=g\bar a$, it follows that \(A\models \varphi(\bar a)\iff A\models \varphi(g\bar a)\). Thus $g$ preserves all atomic formulas, and since $g$ is bijective it follows that \(g\in G\cap \Aut(A)=H\).
    Hence $H$ is closed in $G$.
\end{proof}

\begin{corollary}
    \label{cor:automorphism-groups-are-closed}
    Let $A$ be a structure with domain $\Omega$. Then $\Aut(A)$ is a closed subgroup of $\Sym(\Omega)$.
\end{corollary}

\begin{proof}
    Apply Theorem~\ref{thm:closed-subgroups-automorphism-groups} with $G=\Sym(\Omega)$ and $H=\Aut(A)$.
\end{proof}

When $H$ is a closed subgroup of $\Sym(\Omega)$, the structure constructed in the proof of Theorem~\ref{thm:closed-subgroups-automorphism-groups} is called the \emph{canonical structure} of $H$. Its domain is $\Omega$, and for each $n<\omega$ its basic $n$-ary relations are precisely the $H$-orbits on $\Omega^n$. By construction, \(H=\Aut(A)\).

Thus the canonical structure records exactly the orbit structure of $H$ on finite powers of $\Omega$. In particular, closed subgroups of $\Sym(\Omega)$ and automorphism groups of structures on $\Omega$ may be regarded as two equivalent ways of encoding the same data.

\section{Countable Automorphism Groups as Polish Groups}

We now specialize to the countable case, which is the one relevant for Fra\"{\i}ss\'e limits. Throughout this section, let $\Omega$ be a countable set, and fix an enumeration \(\Omega=\{\omega_1,\omega_2,\dots\}\).
Our goal is to show that automorphism groups of countable structures are Polish. Since $\Sym(\Omega)$ is already a topological group by Proposition~\ref{prop:permutation-topology}, it remains to verify in this setting separability, complete metrizability, and the fact that closed subgroups of Polish groups are again Polish.


\begin{proposition}
    \label{prop:sym-countable-second-countable}
    The permutation topology on $\Sym(\Omega)$ is second countable. Consequently, every subgroup of $\Sym(\Omega)$, endowed with the induced topology, is also second countable.
\end{proposition}

\begin{proof}
    For each $n<\omega$ and each pair of tuples $\bar a,\bar b\in\Omega^n$, the set $S(\bar a,\bar b)$ is basic open. Since $\Omega$ is countable, each $\Omega^n$ is countable, and hence
    \[
        \bigcup_{n<\omega}\Omega^n\times\Omega^n
    \]
    is countable. Therefore there are only countably many basic open sets of the form $S(\bar a,\bar b)$, so $\Sym(\Omega)$ has a countable basis. Since second countability passes to subspaces, every subgroup of $\Sym(\Omega)$ is also second countable.
\end{proof}

\begin{corollary}
    \label{cor:sym-countable-separable}
    The space $\Sym(\Omega)$ is separable. Consequently, every subgroup of $\Sym(\Omega)$ is separable.
\end{corollary}

\begin{proof}
    Every second countable space is separable, so the conclusion follows from Proposition~\ref{prop:sym-countable-second-countable}.
\end{proof}

We now describe an explicit compatible complete metric. First define a metric $d$ on $\Sym(\Omega)$ by
\[
    d(g,h)=
    \begin{cases}
        0,            & \text{if } g=h,                                                                   \\[4pt]
        \dfrac{1}{n}, & \text{if } g\neq h,\ \text{where } n=\min\{m\geq 1:g(\omega_m)\neq h(\omega_m)\}.
    \end{cases}
\]
Thus $d(g,h)$ records the first point of the enumeration on which $g$ and $h$ disagree.

\begin{proposition}
    \label{prop:d-compatible}
    The function $d$ is a metric on $\Sym(\Omega)$, and it induces the permutation topology.
\end{proposition}

\begin{proof}
    The only nontrivial point is the triangle inequality. Let $g,h,k\in\Sym(\Omega)$, and suppose that
    \[
        d(g,k)=\frac{1}{a}, \qquad d(k,h)=\frac{1}{b}.
    \]
    Then $g(\omega_m)=k(\omega_m)=h(\omega_m)$ for every $m<\min(a,b)$. Hence
    \[
        d(g,h)\leq \frac{1}{\min(a,b)}=\max\{d(g,k),d(k,h)\},
    \]
    so $d$ is in fact an ultrametric.

    For $g\in\Sym(\Omega)$ and $n\geq 1$, the open ball of radius $1/(n+1)$ about $g$ is
    \[
        B_d\!\left(g,\frac{1}{n+1}\right)
        =
        \{h\in\Sym(\Omega):h(\omega_m)=g(\omega_m)\text{ for all }m\leq n\}.
    \]
    Thus the $d$-balls are exactly the sets of permutations agreeing with $g$ on a finite initial segment of the enumeration. These form a neighbourhood basis at $g$ for the permutation topology, so $d$ induces that topology.
\end{proof}

The metric $d$ is not complete, since a sequence of longer and longer finite cycles can be Cauchy for $d$ while converging pointwise to an injective non-surjective map. To obtain a complete compatible metric, define \(d'(g,h)=d(g,h)+d(g^{-1},h^{-1})\).

\begin{proposition}
    \label{prop:dprime-compatible-complete}
    The function $d'$ is a complete metric on $\Sym(\Omega)$, and it induces the permutation topology.
\end{proposition}

\begin{proof}
    Since $d$ is a metric, so is $d'$. Moreover, \(d(g,h)\leq d'(g,h)\),
    so every $d'$-open set is $d$-open. Conversely, if $h$ agrees with $g$ on $\omega_1,\dots,\omega_n$, then every permutation sufficiently close to $h$ in the metric $d'$ also agrees with $g$ on those points. Hence every basic $d$-open neighbourhood is $d'$-open, and the two metrics induce the same topology.

    It remains to prove completeness. Let $F=\Omega^\Omega$, and define a metric $\delta$ on $F$ by
    \[
        \delta(f_1,f_2)=
        \begin{cases}
            0,            & \text{if } f_1=f_2,                                                                       \\[4pt]
            \dfrac{1}{n}, & \text{if } f_1\neq f_2,\ \text{where } n=\min\{m\geq 1:f_1(\omega_m)\neq f_2(\omega_m)\}.
        \end{cases}
    \]
    Exactly as above, $(F,\delta)$ is complete. Equip $F\times F$ with the product metric
    \[
        \rho\bigl((f_1,h_1),(f_2,h_2)\bigr)=\delta(f_1,f_2)+\delta(h_1,h_2).
    \]
    Then $(F\times F,\rho)$ is complete.

    Now define \(\Phi:\Sym(\Omega)\to F\times F\) by \(\Phi(g)=(g,g^{-1})\). For $g,h\in\Sym(\Omega)$,
    \[
        \rho(\Phi(g),\Phi(h))
        =\delta(g,h)+\delta(g^{-1},h^{-1})
        =d'(g,h),
    \]
    so $\Phi$ is an isometric embedding.

    Its image consists exactly of the pairs $(f,h)\in F\times F$ satisfying \(f\circ h=\operatorname{id}_{\Omega}\) and \(h\circ f=\operatorname{id}_{\Omega}\). This is a closed subset of $F\times F$, since for each $\omega\in\Omega$ the conditions \(f(h(\omega))=\omega\) and \(h(f(\omega))=\omega\) define closed sets. Therefore $\Phi(\Sym(\Omega))$ is complete, and hence $(\Sym(\Omega),d')$ is complete.
\end{proof}

\begin{definition}
    A \emph{Polish group} is a topological group whose topology is separable and completely metrizable.
\end{definition}

\begin{proposition}
    \label{prop:closed-subgroup-polish}
    Let $G$ be a Polish group and let $H\leq G$ be a closed subgroup. Then $H$, equipped with the subspace topology, is again a Polish group.
\end{proposition}

\begin{proof}
    Since multiplication and inversion on $H$ are restrictions of the corresponding continuous maps on $G$, the group $H$ is a topological group with the subspace topology.

    Let $d$ be a compatible complete metric on $G$. Then the restriction $d|_{H\times H}$ is a compatible metric on $H$, and it is complete because $H$ is closed in the complete metric space $(G,d)$.

    Since $G$ is Polish, it is in particular separable and metrizable, hence second countable. Therefore the subspace $H$ is also second countable, and thus separable. It follows that $H$ is separable and completely metrizable, so $H$ is Polish.
\end{proof}
\begin{theorem}
    \label{thm:aut-countable-polish}
    Let $A$ be a structure with countable domain $\Omega$. Then $\Aut(A)$, endowed with the permutation topology, is a Polish group.
\end{theorem}

\begin{proof}
    By Corollary~\ref{cor:automorphism-groups-are-closed}, the group $\Aut(A)$ is a closed subgroup of $\Sym(\Omega)$. By Proposition~\ref{prop:dprime-compatible-complete}, the group $\Sym(\Omega)$ is completely metrizable, and by Corollary~\ref{cor:sym-countable-separable} it is separable. Since $\Sym(\Omega)$ is a topological group by Proposition~\ref{prop:permutation-topology}, it is therefore a Polish group. Now Proposition~\ref{prop:closed-subgroup-polish} shows that $\Aut(A)$ is Polish.
\end{proof}

\begin{proposition}
    \label{prop:countable-local-basis-open-subgroups}
    Let $G\leq\Sym(\Omega)$. Then the family
    \[
        \{G_{(F)}:F\subseteq \Omega \text{ finite}\}
    \]
    is a countable neighbourhood basis at the identity consisting of open subgroups.
\end{proposition}

\begin{proof}
    Each $G_{(F)}$ is a subgroup of $G$, and by Proposition~\ref{prop:d-compatible} it is open. By Proposition~\ref{prop:sym-countable-second-countable}, the finite subsets of $\Omega$ form a countable family. Finally, every basic open neighbourhood of the identity in $G$ contains $G_{(F)}$ for some finite set $F\subseteq\Omega$, so these subgroups form a neighbourhood basis at the identity.
\end{proof}

In particular, automorphism groups of countable structures are non-archimedean Polish groups: they admit a countable local basis at the identity consisting of open subgroups. This is the topological setting in which we will study flows, minimality, and universal minimal flows in the following chapters.

% !TEX root = main.tex
\chapter[Extreme Amenability and Ramsey Theory]{Extreme Amenability and the Ramsey Property}
\label{chap:extreme-amenability-ramsey}

The aim of this chapter is to connect the dynamical notion of extreme amenability introduced in the previous chapter with the Ramsey property for classes of finite structures. This is the first decisive step in the Kechris--Pestov--Todorcevic correspondence. For closed subgroups of \(S_\infty\), the topology is controlled by stabilisers of finite sets, and this makes it possible to reformulate a fixed-point property for flows in terms of finite colourings of finite configurations. The result is a bridge from topological dynamics to structural Ramsey theory.

We begin by recalling the basic notions of structural Ramsey theory and some standard examples. We then prove a characterisation of extremely amenable closed subgroups of \(S_\infty\) in terms of a Ramsey property for the corresponding finite types. Finally, we translate this characterisation into Fraïssé-theoretic language and obtain the fundamental criterion that if \(\bff\) is the Fraïssé limit of a Fraïssé order class \(\cK\), then \(\Aut(\bff)\) is extremely amenable if and only if \(\cK\) has the Ramsey property. This criterion will be the main tool in the next chapter, where we construct concrete examples of extremely amenable automorphism groups.


\section{Structural Ramsey theory}
\label{sec:structural-ramsey-theory}

We begin by recalling the basic notions of structural Ramsey theory that will be used in the sequel. The point of this section is to formulate the combinatorial property that later characterises extreme amenability for automorphism groups of Fraïssé limits.

Let \(\bfa\) and \(\bfb\) be \(L\)-structures. If \(\bfa\hookrightarrow \bfb\), we let
\[
    \begin{pmatrix} \bfb \\ \bfa \end{pmatrix}
    =
    \{\bfa' \leq \bfb : \bfa' \cong \bfa\},
\]
the set of all substructures of \(\bfb\) isomorphic to \(\bfa\).

If \(\bfa\hookrightarrow \bfb\hookrightarrow \bfc\) and \(k \geq 2\), we write \(\bfc \to (\bfb)^{\bfa}_k\) to mean that for every colouring
\(c \colon \begin{pmatrix} \bfc \\ \bfa \end{pmatrix} \to \{1,\dots,k\}\), there is some
\(\bfb' \in \begin{pmatrix} \bfc \\ \bfb \end{pmatrix}\) such that
\(c \restriction \begin{pmatrix} \bfb' \\ \bfa \end{pmatrix}\) is constant. In this situation we say that \(\bfb'\) is \emph{homogeneous} for the colouring \(c\).

\begin{definition}
    Let \(\cK\) be a class of finite \(L\)-structures. We say that \(\cK\) has the \emph{Ramsey property} if \(\cK\) is hereditary and for every \(\bfa,\bfb \in \cK\) with \(\bfa\hookrightarrow \bfb\) and every integer \(k \geq 2\), there exists \(\bfc \in \cK\) such that
    \[
        \bfc \to (\bfb)^{\bfa}_k.
    \]
\end{definition}

As usual, it is enough to verify the Ramsey property for two colours.

\begin{proposition}
    A hereditary class \(\cK\) of finite structures has the Ramsey property if and only if for every \(\bfa,\bfb \in \cK\) with \(\bfa\hookrightarrow \bfb\), there exists \(\bfc \in \cK\) such that
    \[
        \bfc \to (\bfb)^{\bfa}_2.
    \]
\end{proposition}

\begin{proof}
    Only one implication needs proof. Assume the two-colour condition holds. We show by induction on \(k \geq 2\) that for all \(\bfa\hookrightarrow \bfb\) in \(\cK\) there exists \(\bfc \in \cK\) with \(\bfc \to (\bfb)^{\bfa}_k\).

    The case \(k=2\) is our assumption. Suppose the statement has been proved for \(k-1\), where \(k \geq 3\), and fix \(\bfa\hookrightarrow \bfb\) in \(\cK\). Choose \(\bfd \in \cK\) such that \(\bfd \to (\bfb)^{\bfa}_{k-1}\), and then choose \(\bfc \in \cK\) such that \(\bfc \to (\bfd)^{\bfa}_2\).

    Now let \(c \colon \begin{pmatrix} \bfc \\ \bfa \end{pmatrix} \to \{1,\dots,k\}\) be a \(k\)-colouring. Define a two-colouring \(c' \colon \begin{pmatrix} \bfc \\ \bfa \end{pmatrix} \to \{0,1\}\) by
    \[
        c'(\bfa')=
        \begin{cases}
            0, & c(\bfa')=k,                 \\
            1, & c(\bfa')\in\{1,\dots,k-1\}.
        \end{cases}
    \]
    Since \(\bfc \to (\bfd)^{\bfa}_2\), there is \(\bfd' \in \begin{pmatrix} \bfc \\ \bfd \end{pmatrix}\) homogeneous for \(c'\). If \(c'\) is constantly \(0\) on \(\begin{pmatrix} \bfd' \\ \bfa \end{pmatrix}\), then every copy of \(\bfa\) in \(\bfd'\) has colour \(k\), so \(\bfd'\) contains a homogeneous copy of \(\bfb\). If \(c'\) is constantly \(1\), then the original colouring \(c\) uses only the colours \(1,\dots,k-1\) on \(\begin{pmatrix} \bfd' \\ \bfa \end{pmatrix}\). Since \(\bfd' \cong \bfd\) and \(\bfd \to (\bfb)^{\bfa}_{k-1}\), there is a copy of \(\bfb\) inside \(\bfd'\) homogeneous for \(c\).

    Thus \(\bfc \to (\bfb)^{\bfa}_k\), as required.
\end{proof}

We now record several standard examples.

\begin{example}
    The class \(\cLO\) of finite linear orders has the Ramsey property.
\end{example}

\begin{proof}
    Two finite linear orders are isomorphic if and only if they have the same cardinality. Thus, for finite linear orders \(\bfa\), \(\bfb\), and \(\bfc\), the set
    \(\begin{pmatrix} \bfc \\ \bfa \end{pmatrix}\) may be identified with the family of \(|A|\)-element subsets of \(C\). Under this identification, the statement \(\bfc \to (\bfb)^{\bfa}_k\) is exactly the classical finite Ramsey theorem.
\end{proof}

\begin{example}
    Let \(L=\{<,E\}\), where \(<\) is a binary relation symbol and \(E\) is a binary relation symbol interpreted as a symmetric irreflexive graph relation. Then the class \(\vec{\mathcal{GR}}\) of finite ordered graphs has the Ramsey property.
\end{example}

This is the theorem of Ne\v{s}et\v{r}il and R\"odl~\cite{NR77,NR83}.

\begin{example}
    Let \(F\) be a finite field, and let \(\mathcal{V}_F\) be the class of finite-dimensional vector spaces over \(F\), regarded as structures in a language coding addition and scalar multiplication. Then \(\mathcal{V}_F\) has the Ramsey property.
\end{example}

This is the Graham--Leeb--Rothschild theorem~\cite{GLR}.

\begin{example}
    Let \(\mathcal{BA}\) be the class of finite Boolean algebras in the usual algebraic language. Then \(\mathcal{BA}\) has the Ramsey property.
\end{example}

This is an equivalent formulation of the Dual Ramsey Theorem of Graham and Rothschild; see, for example,~\cite[\S 4.13]{N95}.

\begin{definition}
    A structure \(\bfa\) is said to be \emph{rigid} if every automorphism of \(\bfa\) is the identity. Equivalently, \(\Aut(\bfa)=\{\id_{\bfa}\}\).
    A class of finite structures is called rigid if every member of the class is rigid.
\end{definition}

For our purposes, ordered classes are especially important, since finite linearly ordered structures are rigid. This rigidity ties the Ramsey property to amalgamation.

\begin{proposition}
    \label{prop:rigid-rp-implies-ap}
    Let \(\cK\) be a hereditary class of finite rigid \(L\)-structures. If \(\cK\) has the joint embedding property and the Ramsey property, then \(\cK\) has the amalgamation property.
\end{proposition}

\begin{proof}
    Let \(\bfa,\bfb,\bfc \in \cK\) and let \(f \colon \bfa \hookrightarrow \bfb\) and \(g \colon \bfa \hookrightarrow \bfc\) be embeddings. Since \(\cK\) has the joint embedding property, there exists \(\bfe \in \cK\) into which both \(\bfb\) and \(\bfc\) embed. Fix embeddings \(u \colon \bfb \hookrightarrow \bfe\) and \(v \colon \bfc \hookrightarrow \bfe\). By the Ramsey property, there is \(\bfd \in \cK\) such that \(\bfd \to (\bfe)^{\bfa}_4\).

    We define a colouring of \(\begin{pmatrix} \bfd \\ \bfa \end{pmatrix}\) with values in the four-element set \(\mathcal{P}(\{\bfb,\bfc\})\). For \(\bfa' \in \begin{pmatrix} \bfd \\ \bfa \end{pmatrix}\), put
    \[
        \bfb \in c(\bfa')
        \quad\Longleftrightarrow\quad
        \text{there exists an embedding } r \colon \bfb \hookrightarrow \bfd \text{ with } r(f(A))=A',
    \]
    and define similarly
    \[
        \bfc \in c(\bfa')
        \quad\Longleftrightarrow\quad
        \text{there exists an embedding } s \colon \bfc \hookrightarrow \bfd \text{ with } s(g(A))=A'.
    \]

    Choose \(\bfe' \in \begin{pmatrix} \bfd \\ \bfe \end{pmatrix}\) homogeneous for this colouring. Since \(\bfe' \cong \bfe\) contains copies of both \(\bfb\) and \(\bfc\), there is some copy of \(\bfa\) inside \(\bfe'\) whose colour contains \(\bfb\), and also some copy of \(\bfa\) inside \(\bfe'\) whose colour contains \(\bfc\). Because \(\bfe'\) is homogeneous, the common colour of all members of \(\begin{pmatrix} \bfe' \\ \bfa \end{pmatrix}\) must therefore be exactly \(\{\bfb,\bfc\}\).

    Now choose any \(\bfa' \in \begin{pmatrix} \bfe' \\ \bfa \end{pmatrix}\). Since \(c(\bfa')=\{\bfb,\bfc\}\), there are embeddings \(r \colon \bfb \hookrightarrow \bfd\) and \(s \colon \bfc \hookrightarrow \bfd\) such that \(r(f(A))=A'=s(g(A))\). Hence both \(r \circ f\) and \(s \circ g\) are isomorphisms from \(\bfa\) onto \(\bfa'\). Since \(\bfa\) is rigid, this isomorphism is unique, so \(r \circ f = s \circ g\).
    Thus \(\bfd\), together with the embeddings \(r\) and \(s\), is an amalgam of \(\bfb\) and \(\bfc\) over \(\bfa\). Therefore \(\cK\) has the amalgamation property.
\end{proof}

In particular, if \(\cK\) is a nonempty countable hereditary class of finite rigid structures, has the joint embedding property, contains structures of arbitrarily large cardinality, and has the Ramsey property, then \(\cK\) is a Fraïssé class.

The importance of these observations for us is that many of the classes arising in the Kechris--Pestov--Todorcevic correspondence are order classes. For such classes, rigidity is automatic, and Proposition~\ref{prop:rigid-rp-implies-ap} shows that once the Ramsey property is established, the amalgamation property follows under mild additional hypotheses.



\section{Extreme amenability for closed subgroups of \texorpdfstring{\(S_\infty\)}{S-infinity}}

\label{sec:ea-closed-subgroups-sinf}

We now reformulate extreme amenability for closed subgroups of \(S_\infty\) in combinatorial terms. The basic idea is that for such groups, open subgroups are determined by finite information, and this allows one to express a fixed-point property for flows as a Ramsey-type statement about finite configurations.

We begin with a general criterion for the existence of a fixed point.

\begin{lemma}
    \label{lem:fixed-point-criterion}
    Let \(G\) be a topological group and let \(X\) be a \(G\)-flow. Then the following are equivalent.
    \begin{enumerate}
        \item[(i)] \(X\) has a fixed point.
        \item[(ii)] For every \(n \geq 1\), every continuous map \(f \colon X \to \R^n\), every \(\epsilon >0\), and every finite set \(F \subseteq G\), there exists \(x \in X\) such that
              \[
                  \norm{f(x)-f(g\cdot x)}\leq \epsilon
                  \qquad\text{for all } g\in F,
              \]
              where \(\norm{\cdot}\) denotes the Euclidean norm.
    \end{enumerate}
\end{lemma}

\begin{proof}
    The implication (i) \(\Rightarrow\) (ii) is immediate.

    For (ii) \(\Rightarrow\) (i), define, for each continuous \(f \colon X \to \R^n\), each \(\epsilon>0\), and each finite \(F\subseteq G\),
    \[
        A_{f,\epsilon,F}
        =
        \{x\in X : \norm{f(x)-f(g\cdot x)}\leq \epsilon \text{ for all } g\in F\}.
    \]
    Each set \(A_{f,\epsilon,F}\) is closed.
    We claim that the family \(\{A_{f,\epsilon,F}\}\) has the finite intersection property.
    Indeed, given finitely many triples \((f_j,\epsilon_j,F_j)\), \(j=1,\dots,m\), let \(\bar F=F_1\cup\cdots\cup F_m\) and \(\bar\epsilon=\min\{\epsilon_1,\dots,\epsilon_m\}\), and define \(\bar f=(f_1,\dots,f_m)\colon X \to \R^{n_1+\cdots+n_m}\), where \(f_j\colon X\to\R^{n_j}\).
    Then
    \[
        A_{\bar f,\bar\epsilon,\bar F}
        \subseteq
        \bigcap_{j=1}^m A_{f_j,\epsilon_j,F_j},
    \]
    and \(A_{\bar f,\bar\epsilon,\bar F}\) is nonempty by (ii).
    Since \(X\) is compact, it follows that
    \[
        \bigcap_{f,\epsilon,F} A_{f,\epsilon,F}\neq\emptyset.
    \]
    Choose
    \[
        x\in \bigcap_{f,\epsilon,F} A_{f,\epsilon,F}.
    \]
    We show that \(x\) is fixed. Suppose not. Then there exists \(g\in G\) such that \(g\cdot x\neq x\). By complete regularity of the compact Hausdorff space \(X\), there is a continuous function \(f\colon X\to\R\) such that \(f(x)=0\) and \(f(g\cdot x)=1\).
    But then \(x\notin A_{f,1/2,\{g\}}\), a contradiction. Hence \(x\) is a fixed point.
\end{proof}

We now specialize to closed subgroups of \(S_\infty\).

\begin{proposition}
    \label{prop:coset-characterization}
    Let \(G\leq S_\infty\) be a closed subgroup. Then the following are equivalent.
    \begin{enumerate}
        \item[(i)] \(G\) is extremely amenable.
        \item[(ii)] For every open subgroup \(V\leq G\), every finite colouring \(c\colon G/V\to\{1,\dots,k\}\), and every finite set \(A\subseteq G/V\), there exist \(g\in G\) and \(i\in\{1,\dots,k\}\) such that
              \[
                  c(g\cdot a)=i
                  \qquad\text{for all } a\in A,
              \]
              where \(G\) acts on \(G/V\) by left translation.
    \end{enumerate}
\end{proposition}

\begin{proof}
    Assume first that \(G\) is extremely amenable. Let \(V\leq G\) be open, let \(c\colon G/V\to\{1,\dots,k\}\) be a colouring, and let \(A\subseteq G/V\) be finite. Consider the compact space \(Y=\{1,\dots,k\}^{G/V}\) with the product topology, endowed with the shift action
    \[
        (g\cdot p)(x)=p(g^{-1}\cdot x)
        \qquad
        (g\in G,\ p\in Y,\ x\in G/V).
    \]
    Let \(X=\overline{G\cdot c}\subseteq Y\). Since \(G\) is extremely amenable, \(X\) has a fixed point \(\gamma\). Because the action of \(G\) on \(G/V\) is transitive,
    every fixed point of \(Y\) is constant.
    Thus there is \(i\in\{1,\dots,k\}\) such that \(\gamma(x)=i\) for all \(x\in G/V\). Since \(\gamma\in\overline{G\cdot c}\), there is \(g\in G\) such that \((g^{-1}\cdot c)(a)=\gamma(a)=i\) for all \(a\in A\).
    Equivalently, \(c(g\cdot a)=i\) for all \(a\in A\).

    Conversely, assume (ii), and let \(X\) be a \(G\)-flow. By Lemma~\ref{lem:fixed-point-criterion}, it is enough to verify condition (ii) of that lemma. So fix a continuous map \(f\colon X\to\R^n\), an \(\epsilon>0\), and a finite set \(F\subseteq G\).

    By continuity of the action and compactness of \(X\), there exists an open neighbourhood \(U\) of the identity in \(G\) such that
    \[
        \norm{f(x)-f(u\cdot x)}\leq \epsilon/3
        \qquad\text{for all } x\in X,\ u\in U.
    \]
    Since \(G\) is a subgroup of \(S_\infty\) equipped with the subspace topology, it has a neighbourhood basis at the identity consisting of open subgroups.
    Choose an open subgroup \(V\leq U\).

    Partition the compact set \(f(X)\subseteq\R^n\) into finitely many sets \(A_1,\dots,A_k\) of diameter at most \(\epsilon/3\). Fix \(x_0\in X\) and define \(U_i=\{g\in G : f(g\cdot x_0)\in A_i\}\). Let \(W_i=V U_i=\{vg : v\in V,\ g\in U_i\}\).
    Each \(W_i\) is a union of right cosets of \(V\). Equivalently, each \(W_i^{-1}\) is a union of left cosets of \(V\). Since \(G=U_1\cup\cdots\cup U_k\), we also have \(G=W_1\cup\cdots\cup W_k\).
    Define a colouring \(c\colon G/V\to\{1,\dots,k\}\) so that
    \[
        c(gV)=i
        \quad\Longrightarrow\quad
        g^{-1}\in W_i.
    \]
    This is possible because the sets \(W_i^{-1}\) cover \(G\) and each is a union of left cosets of \(V\).

    Apply (ii) to the finite subset \(A'=\{hV : h\in F\cup\{1_G\}\}\subseteq G/V\). There exist \(g\in G\) and \(i\in\{1,\dots,k\}\) such that \(c(ghV)=i\) for all \(h\in F\cup\{1_G\}\). By the definition of \(c\), this means that \(h^{-1}g^{-1}\in W_i=V U_i\) for all \(h\in F\cup\{1_G\}\), or equivalently, \(hg\in U_i^{-1}V\).
    Replacing \(U_i\) by \(U_i^{-1}\) in the construction if necessary, we may assume that for each \(h\in F\cup\{1_G\}\) there exists \(v_h\in V\) such that \(v_h h g \in U_i\). Set \(x=g\cdot x_0\). Then for each \(h\in F\cup\{1_G\}\) we have \(f(v_h h\cdot x)=f(v_h h g\cdot x_0)\in A_i\). Since \(v_h\in V\subseteq U\), it follows that \(\norm{f(h\cdot x)-f(v_h h\cdot x)}\leq \epsilon/3\).
    Hence \(f(h\cdot x)\) lies in the \(\epsilon/3\)-neighbourhood of \(A_i\) for every \(h\in F\cup\{1_G\}\). As the diameter of \(A_i\) is at most \(\epsilon/3\), we obtain
    \[
        \norm{f(x)-f(h\cdot x)}\leq \epsilon
        \qquad\text{for all } h\in F.
    \]
    Thus \(X\) has a fixed point, and \(G\) is extremely amenable.
\end{proof}

For closed subgroups of \(S_\infty\), it is enough to consider pointwise stabilisers of finite sets. For each nonempty finite \(F\subseteq \N\), let
\[
    G_{(F)}=\{g\in G : g(n)=n \text{ for all } n\in F\}.
\]
These subgroups form a neighbourhood basis at the identity.

We now recast the preceding proposition in terms of finite subsets of \(\N\). The action of \(G\) on \(\N\) induces an action on the set of all finite subsets of \(\N\) by \(g\cdot F=\{g(n): n\in F\}\).
For a finite nonempty set \(F\subseteq\N\), let
\[
    G_{\{F\}}=\{g\in G : g\cdot F=F\}
\]
be its setwise stabiliser. The orbit \(G\cdot F\) will be called the \emph{\(G\)-type} of \(F\). If \(\rho\) and \(\sigma\) are \(G\)-types, we write \(\rho\leq \sigma\) if some, equivalently every, member of \(\sigma\) contains some member of \(\rho\).

We also need the logic action on the space of structures on \(\N\). If \(L\) is a countable signature, let \(X_L\) denote the space of all \(L\)-structures with universe \(\N\). When \(L\) is relational, \(X_L\) is compact, and the natural action of \(S_\infty\) on \(X_L\) is continuous.

\begin{draftnote}
    Suppose that \(L=\{R_i:i\in I\}\) is relational, where \(R_i\) has arity
    \(n_i\). An \(L\)-structure \(\bfa\) with universe \(\N\) is completely
    determined by the truth values
    \[
        R_i^{\bfa}(\bar a),
        \qquad
        i\in I,\quad \bar a\in\N^{n_i}.
    \]
    Hence
    \[
        X_L
        \cong
        \prod_{i\in I}\{0,1\}^{\N^{n_i}},
    \]
    where \(\{0,1\}\) carries the discrete topology. A basic open set specifies
    the truth values of only finitely many atomic statements. Since
    \(\{0,1\}\) is compact, Tychonoff's theorem implies that \(X_L\) is compact.

    The logic action of \(S_\infty\) on \(X_L\) is given by
    \[
        R_i^{g\cdot\bfa}(\bar a)
        \quad\Longleftrightarrow\quad
        R_i^{\bfa}(g^{-1}\bar a).
    \]
    To see that this action is continuous, fix \(g\in S_\infty\),
    \(\bfa\in X_L\), and a basic neighbourhood of \(g\cdot\bfa\), determined
    by finitely many atomic statements \(R_i(\bar a)\). Let \(A\subseteq\N\)
    be the finite set of all entries occurring in these tuples. Require
    \(h\in S_\infty\) to agree with \(g\) on \(g^{-1}(A)\), and require
    \(\bfb\in X_L\) to agree with \(\bfa\) on the corresponding finitely many
    statements \(R_i(g^{-1}\bar a)\). Then
    \(h^{-1}\bar a=g^{-1}\bar a\), and therefore
    \[
        R_i^{h\cdot\bfb}(\bar a)
        \iff
        R_i^{\bfb}(g^{-1}\bar a)
        \iff
        R_i^{\bfa}(g^{-1}\bar a)
        \iff
        R_i^{g\cdot\bfa}(\bar a).
    \]
    Thus a neighbourhood of \((g,\bfa)\) is mapped into the prescribed
    neighbourhood of \(g\cdot\bfa\), proving joint continuity.

    The relational hypothesis is essential for compactness. A function symbol
    of arity \(n\) would be interpreted by an element of
    \(\N^{\N^n}\), and a constant symbol by an element of the infinite discrete
    space \(\N\); these spaces are not compact. In particular,
    \(\LO(\N)\) is compact because it is a closed subspace of
    \(X_{\{<\}}\): every failure of the axioms of a linear order is witnessed
    on finitely many elements.
\end{draftnote}

In the special case \(L=\{<\}\), let \(\LO(\N)\subseteq X_L\) be the compact space of all linear orderings on \(\N\). We say that \(G\) \emph{preserves an ordering} if the \(G\)-flow \(\LO(\N)\) has a fixed point, that is, if there is a linear order on \(\N\) invariant under every element of \(G\).
\begin{proposition}
    \label{prop:type-characterization}
    Let \(G\leq S_\infty\) be a closed subgroup. Then the following are equivalent.
    \begin{enumerate}
        \item[(i)] \(G\) is extremely amenable.
        \item[(ii)]
              \begin{enumerate}
                  \item[(a)] For every finite nonempty \(F\subseteq\N\),
                        \[
                            G_{(F)}=G_{\{F\}}.
                        \]
                  \item[(b)] Whenever \(\rho\leq \sigma\) are \(G\)-types and
                        \[
                            c\colon \rho\to\{1,\dots,k\}
                        \]
                        is a finite colouring, there exist \(i\in\{1,\dots,k\}\) and \(F\in \sigma\) such that
                        \[
                            c(F')=i
                            \qquad\text{for every } F'\subseteq F \text{ with } F'\in \rho.
                        \]
              \end{enumerate}
        \item[(iii)]
              \begin{enumerate}
                  \item[(a')] \(G\) preserves an ordering.
                  \item[(b)] The condition in (ii)(b) holds.\end{enumerate}
    \end{enumerate}
\end{proposition}

\begin{proof}


    Assume first that \(G\) is extremely amenable. Then \(G\) preserves an ordering because \(\LO(\N)\) is a compact \(G\)-flow. This gives (iii)(a').

    To prove (iii)(b), let \(\rho\leq \sigma\) be \(G\)-types and let \(c\colon \rho\to\{1,\dots,k\}\) be a colouring.
    Choose \(F'\in\rho\) and let \(V=G_{(F')}\).
    Since \(G\) preserves an ordering, every finite setwise stabiliser equals the corresponding pointwise stabiliser, so \(G_{(F')}=G_{\{F'\}}\).
    Hence \(G/V \cong G\cdot F'=\rho\).
    Now fix \(F_\sigma\in\sigma\) and consider the finite subset \(A=\{D\subseteq F_\sigma : D\in\rho\}\subseteq \rho\).
    Applying Proposition~\ref{prop:coset-characterization}, we obtain \(g\in G\) and \(i\in\{1,\dots,k\}\) such that \(c(g\cdot D)=i\) for all \(D\in A\).
    Let \(F=g\cdot F_\sigma\in \sigma\).
    If \(F'\subseteq F\) and \(F'\in\rho\), then \(g^{-1}\cdot F'\subseteq F_\sigma\) and \(g^{-1}\cdot F'\in \rho\), so \(c(F')=i\).
    Thus (iii)(b) holds.

    The implication (iii) \(\Rightarrow\) (ii) is immediate, since an invariant linear order on \(\N\) forces every finite setwise stabiliser to equal the corresponding pointwise stabiliser.

    Finally, assume (ii).
    We verify condition (ii) of Proposition~\ref{prop:coset-characterization} for open subgroups of the form \(G_{(F)}\).
    By (ii)(a), these equal the setwise stabilisers \(G_{\{F\}}\).
    Fix a colouring \(c\colon G/G_{\{F\}} \to \{1,\dots,k\}\) and a finite subset \(A\subseteq G/G_{\{F\}}\).
    Identify \(G/G_{\{F\}} \cong G\cdot F = \rho\).
    Let \(F_\sigma=\bigcup A\) and let \(\sigma\) be the \(G\)-type of \(F_\sigma\).
    Then \(\rho\leq \sigma\).
    By (ii)(b), there exist \(i\in\{1,\dots,k\}\) and \(E\in \sigma\) such that \(c(F')=i\) for every \(F'\subseteq E\) with \(F'\in \rho\).
    Since \(E\) and \(F_\sigma\) lie in the same \(G\)-type, we may write \(E=g\cdot F_\sigma\) for some \(g\in G\).
    Every element of \(A\) is a member of \(\rho\) contained in \(F_\sigma\), so after applying \(g\) it becomes a member of \(\rho\) contained in \(E\).
    Hence \(c(g\cdot a)=i\) for all \(a\in A\).
    By Proposition~\ref{prop:coset-characterization}, \(G\) is extremely amenable.
\end{proof}
We now isolate the corresponding Ramsey property for the group itself.

\begin{definition}
    \label{def:group-ramsey-property}
    Let \(G\leq S_\infty\). Suppose that \(\rho\leq \sigma\) are \(G\)-types. For \(F\in \sigma\), define
    \[
        \begin{pmatrix} F \\ \rho \end{pmatrix}
        =
        \{F'\subseteq F : F'\in \rho\}.
    \]
    If \(\rho\leq \sigma\leq \tau\) are \(G\)-types and \(k\geq 2\), we write
    \[
        \tau \to (\sigma)^\rho_k
    \]
    if for every \(F\in\tau\) and every colouring
    \(c\colon \begin{pmatrix} F \\ \rho \end{pmatrix}\to\{1,\dots,k\}\), there exists
    \(F_\sigma\in \begin{pmatrix} F \\ \sigma \end{pmatrix}\) such that
    \(c\restriction \begin{pmatrix} F_\sigma \\ \rho \end{pmatrix}\) is constant.

    We say that \(G\) has the \emph{Ramsey property} if for all \(G\)-types \(\rho\leq \sigma\) and all \(k\geq 2\), there exists a \(G\)-type \(\tau\geq \sigma\) such that
    \[
        \tau \to (\sigma)^\rho_k.
    \]
\end{definition}

\begin{theorem}
    \label{thm:ea-iff-ordering-plus-group-ramsey}
    Let \(G\leq S_\infty\) be a closed subgroup. Then the following are equivalent.
    \begin{enumerate}
        \item[(i)] \(G\) is extremely amenable.
        \item[(ii)] \(G\) preserves an ordering and has the Ramsey property.
    \end{enumerate}
\end{theorem}

\begin{proof}

    Assume that \(G\) is extremely amenable. Then \(G\) preserves an ordering by Proposition~\ref{prop:type-characterization}.

    We show that \(G\) has the Ramsey property. As in the usual Ramsey argument, it is enough to treat the case of two colours. Suppose instead that there exist \(G\)-types \(\rho\leq \sigma\) such that no \(G\)-type \(\tau\geq \sigma\) satisfies
    \[
        \tau\to (\sigma)^\rho_2.
    \]
    Fix \(F_\sigma\in \sigma\). Then for every finite set \(E\supseteq F_\sigma\) there exists a colouring \(c_E\colon \begin{pmatrix} E \\ \rho \end{pmatrix}\to\{1,2\}\) with no homogeneous member of \(\begin{pmatrix} E \\ \sigma \end{pmatrix}\).

    Let \(I\) be the set of all finite nonempty subsets of \(\N\), and choose an ultrafilter \(\mathcal U\) on \(I\) containing every set of the form \(\{E\in I : F\subseteq E\}\), where \(F\subseteq\N\) is finite. For each \(D\in \rho\), exactly one of the sets
    \[
        \{E\supseteq D\cup F_\sigma : c_E(D)=1\},
        \qquad
        \{E\supseteq D\cup F_\sigma : c_E(D)=2\}
    \]
    belongs to \(\mathcal U\).
    Define \(c(D)=i\) if the \(i\)th set belongs to \(\mathcal U\).
    This gives a colouring \(c\colon \rho\to\{1,2\}\).
    By Proposition~\ref{prop:type-characterization}, there exists \(F\in \sigma\) such that \(c\restriction \begin{pmatrix} F \\ \rho \end{pmatrix}\) is constant, say constantly equal to \(i\).
    For each \(D\in \begin{pmatrix} F \\ \rho \end{pmatrix}\), the set \(A_D=\{E\supseteq F\cup F_\sigma : c_E(D)=i\}\) belongs to \(\mathcal U\).
    Since there are only finitely many such \(D\), their intersection is nonempty.
    Choose
    \[
        E\in \bigcap_{D\in \begin{pmatrix} F \\ \rho \end{pmatrix}} A_D.
    \]
    Then \(E\supseteq F_\sigma\) and \(c_E(D)=i\) for all \(D\in \begin{pmatrix} F \\ \rho \end{pmatrix}\).
    Thus \(F\) is a homogeneous member of \(\begin{pmatrix} E \\ \sigma \end{pmatrix}\) for the colouring \(c_E\), contradicting the choice of \(c_E\).
    Therefore \(G\) has the Ramsey property.

    Conversely, assume that \(G\) preserves an ordering and has the Ramsey property. To prove extreme amenability, by Proposition~\ref{prop:type-characterization} it is enough to verify condition (ii)(b) there. Let \(\rho\leq \sigma\) be \(G\)-types and let \(c\colon \rho\to\{1,\dots,k\}\) be a colouring. Since \(G\) has the Ramsey property, there exists a \(G\)-type \(\tau\geq \sigma\) such that
    \[
        \tau\to (\sigma)^\rho_k.
    \]
    Choose \(F\in\tau\). Then the restriction of \(c\) to \(\begin{pmatrix} F \\ \rho \end{pmatrix}\) has a homogeneous member of \(\begin{pmatrix} F \\ \sigma \end{pmatrix}\), which is exactly the statement required in Proposition~\ref{prop:type-characterization}. Hence \(G\) is extremely amenable.
\end{proof}



\begin{remark}
    \label{rem:cofinal-types}
    A family \(T\) of \(G\)-types is called \emph{cofinal} if for every \(G\)-type \(\rho\) there exists \(\sigma\in T\) such that \(\rho\leq \sigma\).
    It is enough in Definition~\ref{def:group-ramsey-property} to consider \(\rho\), \(\sigma\), and \(\tau\) ranging over any fixed cofinal family of \(G\)-types.
\end{remark}



\section[Fraïssé order classes]{Fraïssé order classes and extremely amenable\\automorphism groups}
\label{sec:fraisse-order-classes-and-ea-groups}

We now translate the characterisation obtained in the previous section into Fraïssé-theoretic language.

Let \(L\) be a signature containing a distinguished binary relation symbol \(<\). An \emph{order structure} for \(L\) is an \(L\)-structure \(\bfa\) such that its interpretation \(\prec^{\bfa}\) of \(<\) is a linear ordering on \(A\). A class \(\cK\) of finite \(L\)-structures is called an \emph{order class} if every member of \(\cK\) is an order structure.

The next theorem identifies the extremely amenable closed subgroups of \(S_\infty\) with automorphism groups of Fraïssé limits of Ramsey order classes.

\begin{theorem}
    \label{thm:ea-iff-fraisse-order-ramsey}
    Let \(G\leq S_\infty\) be a closed subgroup. Then the following are equivalent.
    \begin{enumerate}
        \item[(i)] \(G\) is extremely amenable.
        \item[(ii)] There is a Fraïssé order class \(\cK\) with the Ramsey property such that \(G=\Aut(\bff)\), where \(\bff=\Flim(\cK)\).
    \end{enumerate}
\end{theorem}

\begin{proof}


    Assume first that \(G\) is extremely amenable. Let \(\bfa_G\) be the canonical structure on \(\N\) associated with \(G\); recall from Chapter~\ref{chap:fraisse-automorphism-groups} that its basic \(n\)-ary relations are the \(G\)-orbits on \(\N^n\), and that \(\Aut(\bfa_G)=G\).
    By construction, the structure \(\bfa_G\) is ultrahomogeneous and has a relational signature.

    Since \(G\) is extremely amenable, Theorem~\ref{thm:ea-iff-ordering-plus-group-ramsey} shows that \(G\) preserves an ordering on \(\N\).
    Let \(\prec^{\bfa}\) be a linear ordering of \(\N\) fixed by every element of \(G\).
    Expand \(\bfa_G\) by a new binary relation symbol \(<\) interpreted as \(\prec^{\bfa}\), and denote the resulting structure by \(\bfa=\langle \bfa_G,\prec^{\bfa}\rangle\).
    Then every element of \(G\) preserves \(\prec^{\bfa}\), so \(G\leq \Aut(\bfa)\).
    Conversely, every automorphism of \(\bfa\) is in particular an automorphism of \(\bfa_G\), hence belongs to \(G\).
    Therefore \(\Aut(\bfa)=G\).

    Since \(\bfa_G\) is ultrahomogeneous and \(\prec^{\bfa}\) is \(G\)-invariant, the expansion \(\bfa\) is also ultrahomogeneous: if \(\varphi\colon \bfb\hookrightarrow \bfc\) is an isomorphism between finite substructures of \(\bfa\), then \(\varphi\) is in particular an isomorphism between finite substructures of \(\bfa_G\), so it extends to some \(g\in G=\Aut(\bfa_G)\); because \(g\) preserves \(\prec^{\bfa}\), it is an automorphism of \(\bfa\) extending \(\varphi\).
    Because the signature is relational, finitely generated substructures of \(\bfa\) are finite, so \(\bfa\) is locally finite.
    It follows that \(\cK:=\Age(\bfa)\) is a Fraïssé order class with Fraïssé limit \(\bfa\).

    It remains to show that \(\cK\) has the Ramsey property.
    Let \(\bfb,\bfc\in\cK\) with \(\bfb\hookrightarrow \bfc\), and let \(k\geq 2\).
    Consider the \(G\)-types \(\rho_{\bfb}\) and \(\rho_{\bfc}\) consisting of all subsets of \(\N\) which, with the induced structure from \(\bfa\), are isomorphic to \(\bfb\) and \(\bfc\), respectively.
    Since \(\bfa\) is ultrahomogeneous, copies of a finite substructure in \(\bfa\) form a single \(G\)-type.
    By Theorem~\ref{thm:ea-iff-ordering-plus-group-ramsey}, the group \(G\) has the Ramsey property.
    Hence there exists a \(G\)-type \(\tau\geq \rho_{\bfc}\) such that \(\tau\to (\rho_{\bfc})^{\rho_{\bfb}}_k\).
    Choose a finite substructure \(\bfd\leq \bfa\) whose underlying set belongs to \(\tau\).
    Then \(\bfd\in\cK\), and the statement \(\bfd\to (\bfc)^{\bfb}_k\) is exactly the same as \(\tau\to (\rho_{\bfc})^{\rho_{\bfb}}_k\).
    Thus \(\cK\) has the Ramsey property.

    This proves (i) \(\Rightarrow\) (ii).

    Now assume (ii). Let \(\bff=\Flim(\cK)\) and \(G=\Aut(\bff)\), where \(\cK\) is a Fraïssé order class with the Ramsey property.

    Since \(\bff\) is an order structure, its linear order \(\prec^{\bff}\) on \(F\) is preserved by every automorphism of \(\bff\). Thus \(G\) preserves an ordering.

    To show that \(G\) is extremely amenable, it is enough by Theorem~\ref{thm:ea-iff-ordering-plus-group-ramsey} to verify that \(G\) has the Ramsey property. By Remark~\ref{rem:cofinal-types}, it suffices to check this on any cofinal family of \(G\)-types. We take the family of \(G\)-types arising from finite substructures of \(\bff\).

    Let \(\bfb,\bfc\in\Age(\bff)=\cK\) with \(\bfb\hookrightarrow \bfc\), and let \(\rho_{\bfb},\rho_{\bfc}\) be the corresponding \(G\)-types of copies of \(\bfb\) and \(\bfc\) in \(\bff\). Since \(\cK\) has the Ramsey property, for every \(k\geq 2\) there exists \(\bfd\in\cK\) such that \(\bfd\to (\bfc)^{\bfb}_k\). Let \(\tau\) be the \(G\)-type of copies of \(\bfd\) in \(\bff\). Then \(\tau\to (\rho_{\bfc})^{\rho_{\bfb}}_k\).
    Therefore \(G\) has the Ramsey property on this cofinal family of types, hence \(G\) has the Ramsey property. Since \(G\) also preserves an ordering, Theorem~\ref{thm:ea-iff-ordering-plus-group-ramsey} yields that \(G\) is extremely amenable.
\end{proof}

The preceding theorem immediately gives the Fraïssé-theoretic form that will be used in the next chapter.

\begin{corollary}
    \label{cor:fraisse-order-class-ea}
    Let \(\cK\) be a Fraïssé order class and let \(\bff=\Flim(\cK)\).
    Then the following are equivalent.
    \begin{enumerate}
        \item[(i)] \(\Aut(\bff)\) is extremely amenable.
        \item[(ii)] \(\cK\) has the Ramsey property.
    \end{enumerate}
\end{corollary}

\begin{proof}
    Apply Theorem~\ref{thm:ea-iff-fraisse-order-ramsey} with \(G=\Aut(\bff)\).
\end{proof}

The preceding corollary is the form of the Kechris--Pestov--Todorcevic correspondence that will be used in the next chapter. It shows that, for a Fraïssé order class, the extreme amenability of the automorphism group of its Fraïssé limit is exactly equivalent to the Ramsey property of the class. In Chapter~\ref{chap:examples-extreme-amenability} we apply this criterion to concrete examples.

% !TEX root = main.tex
\chapter[Examples of Extreme Amenability]{Examples of Extremely Amenable Automorphism Groups}
\label{chap:examples-extreme-amenability}
In Chapter~\ref{chap:extreme-amenability-ramsey} we proved that if \(\cK\) is a Fraïssé order class with Fraïssé limit \(\bff\), then \(\Aut(\bff)\) is extremely amenable if and only if \(\cK\) has the Ramsey property. The aim of this chapter is to apply that criterion to several standard Fraïssé classes and obtain concrete examples of extremely amenable automorphism groups.

We begin with classes of finite ordered graphs and hypergraphs, where the required Ramsey property is provided by the theorems of Ne\v{s}et\v{r}il and R\"odl and their extensions. We then consider finite-dimensional vector spaces over a fixed finite field and finite Boolean algebras, which yield further examples of extremely amenable automorphism groups. Finally, we turn to finite metric spaces and the ordered rational Urysohn space, leading to Pestov's theorem on the extreme amenability of the isometry group of the Urysohn space.

The purpose of the chapter is not to give an exhaustive catalogue of Ramsey classes, but to show how the criterion from Chapter~\ref{chap:extreme-amenability-ramsey} works in a range of important settings. These examples will also prepare the ground for the later computation of universal minimal flows via ordered expansions.

\section{Graphs and equivalence relations}

We begin with finite ordered graphs and some closely related classes. Let \(L_0=\{E\}\) and \(L=\{E,<\}\), where \(E\) is a binary relation symbol and \(<\) is intended to denote a linear order. An \(L_0\)-structure \(\bfa_0\) is called a \emph{graph} if \(E^{\bfa_0}\) is irreflexive and symmetric. An \emph{ordered graph} is an \(L\)-structure \(\bfa=\langle \bfa_0,\prec^{\bfa}\rangle\) such that \(\bfa_0\) is a graph and \(\prec^{\bfa}\), the interpretation of \(<\), is a linear ordering of \(A\).

The most important Fraïssé classes of finite graphs for our purposes are the class \(\mathcal{GR}\) of all finite graphs and, for each \(n\geq 3\), the class \(\Forb(K_n)\) of all finite \(K_n\)-free graphs. Their ordered expansions will give the first basic examples of extremely amenable automorphism groups. We will also consider finite equivalence relations, which provide a useful contrast: some ordered expansions are not Ramsey, while a more carefully chosen ordered expansion does lead to extreme amenability.

\subsection[Ordered graphs and \texorpdfstring{\(K_n\)}{K sub n}-free ordered graphs]{Ordered graphs and \(K_n\)-free ordered graphs}

Let \(\vec{\mathcal{GR}}\) denote the class of all finite ordered graphs, and for each \(n\geq 3\) let \(\vec{\Forb}(K_n)\) denote the class of all finite ordered \(K_n\)-free graphs.

Both classes are hereditary and clearly contain structures of arbitrarily large cardinality. Since \(\mathcal{GR}\) and \(\Forb(K_n)\) have strong amalgamation and strong joint embedding, it follows that \(\vec{\mathcal{GR}}\) and \(\vec{\Forb}(K_n)\) are Fraïssé order classes. The Ramsey property for these classes is due to Ne\v{s}et\v{r}il and R\"odl~\cite{NR77,NR83}; see also~\cite{N89}. Therefore Chapter~\ref{chap:extreme-amenability-ramsey} applies immediately.

\begin{theorem}
    \label{thm:random-ordered-graph-ea}
    Let \(\bfo=\Flim(\vec{\mathcal{GR}})\) be the random ordered graph. Then \(\Aut(\bfo)\) is extremely amenable.
\end{theorem}

\begin{proof}
    The class \(\vec{\mathcal{GR}}\) is a Fraïssé order class with the Ramsey property, so the conclusion follows from Corollary~\ref{cor:fraisse-order-class-ea}.
\end{proof}

\begin{theorem}
    \label{thm:kn-free-ordered-graph-ea}
    For each \(n\geq 3\), let \(\bfo_n=\Flim(\vec{\Forb}(K_n))\) be the random ordered \(K_n\)-free graph. Then \(\Aut(\bfo_n)\) is extremely amenable.
\end{theorem}

\begin{proof}
    The class \(\vec{\Forb}(K_n)\) is a Fraïssé order class with the Ramsey property, so again the result follows from Corollary~\ref{cor:fraisse-order-class-ea}.
\end{proof}

These are the basic graph-theoretic examples. In contrast, the unordered automorphism groups \(\Aut(\Flim(\mathcal{GR}))\) and \(\Aut(\Flim(\Forb(K_n)))\) are not extremely amenable, since they do not preserve a linear ordering.


\subsection{A failure of the Ramsey property: equivalence relations}

We now turn to finite equivalence relations. Strictly speaking, an equivalence relation is not a graph, since it is reflexive. Nevertheless, it is convenient to treat it in parallel with the graph examples by suppressing the diagonal.

Let \(\mathcal{EQ}\) denote the class of finite equivalence relations, viewed as \(L_0\)-structures, and let \(\vec{\mathcal{EQ}}\) be the class of all finite linearly ordered equivalence relations. This class is a Fraïssé order class. Its Fraïssé limit is, up to isomorphism,
\[
    \langle \Q,E,\prec\rangle,
\]
where \(\prec\) is the usual order on \(\Q\) and \(E\) is an equivalence relation with infinitely many classes, each of which is dense in \(\Q\). Unlike the graph classes above, \(\vec{\mathcal{EQ}}\) does \emph{not} have the Ramsey property.

\begin{theorem}
    \label{thm:eq-not-ramsey}
    The class \(\vec{\mathcal{EQ}}\) does not have the Ramsey property. Consequently, the automorphism group of
    \[
        \langle \Q,E,\prec\rangle,
    \]
    where \(E\) has infinitely many dense classes, is not extremely amenable.
\end{theorem}

\begin{proof}
    Let \(\bfa\) be the ordered structure with underlying set \(A=\{a,b\}\), order \(a\prec^{\bfa} b\), and \(E^{\bfa}\)-classes \(\{a\}\) and \(\{b\}\). Let \(\bfb\) be the ordered structure with underlying set \(B=\{a,b,c\}\), order \(a\prec^{\bfb} b\prec^{\bfb} c\), and \(E^{\bfb}\)-classes \(\{a,c\}\) and \(\{b\}\). We claim that there is no \(\bfc\in \vec{\mathcal{EQ}}\) such that
    \[
        \bfc\to (\bfb)^{\bfa}_2.
    \]
    Let \(\bfc\) be any member of \(\vec{\mathcal{EQ}}\). Order the \(E^{\bfc}\)-classes by the least element they contain. Colour a copy \(\{x,y\}\) of \(\bfa\) in \(\bfc\) red if the \(E^{\bfc}\)-class of \(x\) is lower than the \(E^{\bfc}\)-class of \(y\), and colour it blue otherwise.

    Now consider any copy \(\{x,y,z\}\) of \(\bfb\) in \(\bfc\), with \(x\prec^{\bfc} y\prec^{\bfc} z\).
    Since \(x\) and \(z\) lie in the same equivalence class while \(y\) lies in a different one, the pair \(\{x,y\}\) and the pair \(\{y,z\}\) must receive different colours. Hence no copy of \(\bfb\) is monochromatic. Therefore \(\vec{\mathcal{EQ}}\) fails the Ramsey property.

    The final statement follows from Corollary~\ref{cor:fraisse-order-class-ea}.
\end{proof}

This example is useful because it shows that simply adding a linear order to a Fraïssé class does not automatically produce a Ramsey class.

\subsection{Convexly ordered equivalence relations}

There is, however, a natural ordered expansion of finite equivalence relations that \emph{does} have the Ramsey property.

Let \(\cK\) be the class of all finite \(L\)-structures
\(\bfa=\langle A,E^{\bfa},\prec^{\bfa}\rangle\) such that \(E^{\bfa}\) is an
equivalence relation on \(A\), \(\prec^{\bfa}\) is the interpretation of \(<\) and is a linear ordering on \(A\), and
each \(E^{\bfa}\)-class is convex with respect to \(\prec^{\bfa}\); that is,
whenever
\[
    a\prec^{\bfa} b\prec^{\bfa} c
    \quad\text{and}\quad
    a\,E^{\bfa}\,c,
\]
we also have \(a\,E^{\bfa}\,b\).

It is straightforward to check that \(\cK\) is a Fraïssé order class. Its Fraïssé limit is, up to isomorphism,
\[
    \bff=\langle \Q^2,E,\prec_{\ell}\rangle,
\]
where \(\prec_{\ell}\) is the lexicographic ordering on \(\Q^2\) and
\[
    (r,s)\,E\,(r',s')
    \quad\Longleftrightarrow\quad
    r=r'.
\]

To prove that \(\Aut(\bff)\) is extremely amenable, we use the following standard closure properties.

\begin{lemma}
    \label{lem:ea-closure-properties}
    \leavevmode
    \begin{enumerate}
        \item[(i)] Let \(G\) be a topological group and let \(\mathcal{H}\) be a directed family of extremely amenable subgroups of \(G\). If \(\bigcup \mathcal{H}\) is dense in \(G\), then \(G\) is extremely amenable.
        \item[(ii)] Let \(N\) be a closed normal subgroup of a topological group \(G\). If \(N\) and \(G/N\) are extremely amenable, then \(G\) is extremely amenable.
        \item[(iii)] An arbitrary direct product of extremely amenable groups is extremely amenable.
    \end{enumerate}
\end{lemma}

\begin{proof}

    For (i), let \(X\) be a compact \(G\)-flow. For each \(H\in\mathcal{H}\), the set \(X^H=\{x\in X : h\cdot x=x \text{ for all } h\in H\}\)
    is nonempty and closed. Since \(\mathcal{H}\) is directed, the family \(\{X^H : H\in\mathcal{H}\}\) has the finite intersection property, so compactness gives
    \[
        \bigcap_{H\in\mathcal{H}} X^H\neq\emptyset.
    \]
    Any point in this intersection is fixed by \(\bigcup\mathcal{H}\), hence by all of \(G\) because the action is continuous and \(\bigcup\mathcal{H}\) is dense.

    For (ii), let \(X\) be a compact \(G\)-flow. Since \(N\) is extremely amenable, the fixed-point set \(X^N=\{x\in X : n\cdot x=x \text{ for all } n\in N\}\)
    is nonempty. It is closed and \(G\)-invariant because \(N\) is normal. The induced action of \(G/N\) on \(X^N\) is continuous, and since \(G/N\) is extremely amenable, there is a point of \(X^N\) fixed by \(G/N\), hence by \(G\).

    For (iii), every finite subproduct is extremely amenable by repeated use of (ii). These finite subproducts form a directed family whose union is dense in the full product, so (i) applies.
\end{proof}

\begin{theorem}
    \label{thm:convex-eq-ea}
    Let
    \[
        \bff=\langle \Q^2,E,\prec_{\ell}\rangle
    \]
    be as above. Then \(\Aut(\bff)\) is extremely amenable.
\end{theorem}

\begin{proof}


    For each \(r\in \Q\), let \(I_r=\{r\}\times \Q\).
    Any automorphism \(\pi\in \Aut(\bff)\) permutes the fibres \(I_r\), since these are exactly the \(E\)-classes. Thus \(\pi\) induces an order automorphism \(f_{\pi}\in \Aut(\langle \Q,\prec\rangle)\) given by \(\pi(I_r)=I_{f_{\pi}(r)}\).
    Moreover, for each \(r\in \Q\), the restriction of \(\pi\) to the fibre \(I_r\) induces an order automorphism \((g_{\pi})_r\in \Aut(\langle \Q,\prec\rangle)\).
    Therefore we obtain a homomorphism
    \[
        \Theta\colon \Aut(\bff)\to \Aut(\langle \Q,\prec\rangle)\ltimes \Aut(\langle \Q,\prec\rangle)^{\Q},
    \]
    where the first factor acts on the second by coordinate shift.

    It is straightforward to check that \(\Theta\) is a topological group isomorphism. Hence it is enough to show that
    \[
        \Aut(\langle \Q,\prec\rangle)\ltimes \Aut(\langle \Q,\prec\rangle)^{\Q}
    \]
    is extremely amenable.

    Since the class of finite linear orders has the Ramsey property,
    Corollary~\ref{cor:fraisse-order-class-ea} shows that
    \(\Aut(\langle \Q,\prec\rangle)\) is extremely amenable. By
    Lemma~\ref{lem:ea-closure-properties}(iii), so is the product
    \(\Aut(\langle \Q,\prec\rangle)^{\Q}\).
    Applying Lemma~\ref{lem:ea-closure-properties}(ii), we conclude that the semidirect product is extremely amenable, and therefore so is \(\Aut(\bff)\).
\end{proof}

\begin{corollary}
    \label{cor:convex-eq-ramsey}
    The class of finite convexly ordered equivalence relations has the Ramsey property.
\end{corollary}

\begin{proof}
    This follows from Theorem~\ref{thm:convex-eq-ea} and Corollary~\ref{cor:fraisse-order-class-ea}.
\end{proof}

The contrast between Theorem~\ref{thm:eq-not-ramsey} and Corollary~\ref{cor:convex-eq-ramsey} is instructive. In the first case, an arbitrary linear ordering does not interact well with the equivalence relation. In the second, convexity aligns the order with the equivalence classes, and the resulting ordered expansion fits the Ramsey framework.


\section{Hypergraphs}

The graph examples extend naturally to hypergraphs. Let \(L_0=\{R_i : i\in I\}\) be a finite relational signature, where each \(R_i\) has arity \(n_i\geq 2\). An \(L_0\)-structure \(\bfa_0\) will be called a \emph{hypergraph} of type \(L_0\) if, for each \(i\in I\), the relation \(R_i^{\bfa_0}\) is irreflexive and symmetric. An \emph{ordered hypergraph} is an \(L\)-structure \(\bfa=\langle \bfa_0,\prec^{\bfa}\rangle\), where \(L=L_0\cup\{<\}\), \(\bfa_0\) is a hypergraph of type \(L_0\), and \(\prec^{\bfa}\), the interpretation of \(<\), is a linear ordering of \(A\).

We first consider the class of all finite ordered hypergraphs of a fixed type. We then record a more general family obtained by forbidding finite irreducible substructures.

\subsection{Ordered hypergraphs}

Let \(\vec{\mathcal H}_{L_0}\) denote the class of all finite ordered hypergraphs of type \(L_0\). This is a Fraïssé order class. Indeed, heredity is immediate, and joint embedding and amalgamation are obtained by taking disjoint unions over the underlying ordered hypergraphs.

The Ramsey property for such classes was proved by Ne\v{s}et\v{r}il and R\"odl, extending the graph case~\cite{NR77,NR83}. Hence Chapter~\ref{chap:extreme-amenability-ramsey} applies exactly as before.

\begin{theorem}
    \label{thm:ordered-hypergraph-ea}
    Let
    \[
        \bff=\Flim(\vec{\mathcal H}_{L_0}).
    \]
    Then \(\Aut(\bff)\) is extremely amenable.
\end{theorem}

\begin{proof}
    The class \(\vec{\mathcal H}_{L_0}\) is a Fraïssé order class with the Ramsey property, so the conclusion follows from Corollary~\ref{cor:fraisse-order-class-ea}.
\end{proof}

Thus every fixed finite hypergraph type yields an extremely amenable automorphism group once a linear ordering is added.

\subsection{Forbidden irreducible hypergraphs}

The previous example admits a substantial generalisation. Let \(\mathcal A\) be a family of finite hypergraphs of type \(L_0\). We write \(\Forb(\mathcal A)\) for the class of all finite hypergraphs of type \(L_0\) that admit no embedding of a member of \(\mathcal A\).

To obtain a Fraïssé class in this setting, one usually assumes that the forbidden structures are \emph{irreducible}. Recall that a finite hypergraph \(\bfa\) is called \emph{irreducible} if any two distinct vertices of \(A\) lie together in some relation of \(\bfa\). Equivalently, \(\bfa\) cannot be written as a free amalgam of two proper substructures.

Let \(\vec{\Forb}(\mathcal A)\) denote the class of all finite linearly ordered members of \(\Forb(\mathcal A)\).

\begin{theorem}
    \label{thm:forb-irreducible-hypergraph-ea}
    Assume that \(\mathcal A\) is a family of finite irreducible hypergraphs of type \(L_0\). Then \(\vec{\Forb}(\mathcal A)\) is a Fraïssé order class with the Ramsey property. Consequently, if
    \[
        \bff=\Flim(\vec{\Forb}(\mathcal A)),
    \]
    then \(\Aut(\bff)\) is extremely amenable.
\end{theorem}

\begin{proof}
    It is a standard result of Ne\v{s}et\v{r}il and R\"odl, in the form generalised by Abramson and Harrington~\cite{AH78}, that \(\vec{\Forb}(\mathcal A)\) is a Ramsey class whenever the forbidden structures are finite and irreducible. The hereditary property is immediate, and the Fraïssé properties are standard for this class. Therefore \(\vec{\Forb}(\mathcal A)\) is a Fraïssé order class with the Ramsey property. The conclusion then follows from Corollary~\ref{cor:fraisse-order-class-ea}.
\end{proof}

This theorem contains the ordered graph examples from the previous section as special cases. Indeed, if \(L_0=\{E\}\) and \(\mathcal A=\{K_n\}\), then \(\vec{\Forb}(\mathcal A)=\vec{\Forb}(K_n)\).
The hypergraph setting is therefore not merely parallel to the graph case, but a genuine extension of it.

\section{Vector spaces over finite fields}

We now turn to finite-dimensional vector spaces over a fixed finite field. This provides a different kind of example, in which the ordering is not arbitrary but is induced by an ordered basis.

Fix a finite field \(F\). Let \(\mathcal{V}_F\) denote the class of all finite-dimensional vector spaces over \(F\), regarded as structures in the language of vector spaces over \(F\). This is a Fraïssé class, and its Fraïssé limit is the countably infinite-dimensional vector space over \(F\).

To obtain an ordered expansion, we use the anti-lexicographic ordering associated with an ordered basis. Let \(\bfv\) be a finite-dimensional vector-space structure over \(F\), with underlying set \(V\), and let \(b_1\prec^{\bfv}\cdots\prec^{\bfv}b_n\) be an ordered basis of \(\bfv\). Every element \(v\in V\) may be written uniquely in the form
\[
    v=\alpha_1b_1+\cdots+\alpha_n b_n,
    \qquad
    \alpha_i\in F.
\]
Fix once and for all a linear ordering of the field \(F\) with \(0\) as its least element. This induces a linear ordering \(\prec^{\bfv}\) on \(V\) by declaring that \(v\prec^{\bfv}w\) if, at the largest index \(i\) for which the coefficients of \(v\) and \(w\) differ, the coefficient of \(v\) is smaller than the coefficient of \(w\). This is the anti-lexicographic ordering determined by the ordered basis \(b_1\prec^{\bfv}\cdots\prec^{\bfv}b_n\).

An ordering of a finite-dimensional vector space obtained in this way will be called a \emph{natural ordering}. Let \(\mathcal{OV}_F\) denote the class of all finite-dimensional vector spaces over \(F\) equipped with a natural ordering.

The class \(\mathcal{OV}_F\) is a Fraïssé order class. Indeed, heredity is immediate, and the amalgamation property follows from the corresponding property for finite-dimensional vector spaces together with the fact that natural orderings are determined by ordered bases.

The Ramsey property for \(\mathcal{OV}_F\) is a consequence of the Graham--Leeb--Rothschild theorem; see also Thomas's formulation for naturally ordered finite vector spaces~\cite{GLR,Th}. Hence Chapter~\ref{chap:extreme-amenability-ramsey} applies exactly as in the previous examples.

\begin{theorem}
    \label{thm:finite-field-vector-spaces-ea}
    Let
    \[
        \bff_F=\Flim(\mathcal{OV}_F).
    \]
    Then \(\Aut(\bff_F)\) is extremely amenable.
\end{theorem}

\begin{proof}
    The class \(\mathcal{OV}_F\) is a Fraïssé order class with the Ramsey property. Therefore Corollary~\ref{cor:fraisse-order-class-ea} yields that \(\Aut(\bff_F)\) is extremely amenable.
\end{proof}

Thus naturally ordered finite-dimensional vector spaces over a finite field give another family of extremely amenable automorphism groups. Unlike the graph and hypergraph examples, the ordering here is tied to the underlying algebraic structure, and this compatibility is exactly what makes the Ramsey property available.

\section{Boolean algebras}

We next consider finite Boolean algebras. Let \(L_0=\{0,1,\neg,\wedge,\vee\}\), where \(0\) and \(1\) are constant symbols, \(\neg\) is unary, and \(\wedge,\vee\) are binary. Let \(\mathcal{BA}\) denote the class of all finite Boolean algebras in this language. This is a Fraïssé class, and its Fraïssé limit is the countable atomless Boolean algebra.

To obtain an ordered expansion, we define a natural ordering on each finite Boolean algebra in terms of its atoms. Let \(\bfb\) be a finite Boolean algebra with underlying set \(B\), and let \(a_1\prec^{\bfb}\cdots\prec^{\bfb}a_n\) be a linear ordering of its atoms. Every element \(x\in B\) has a unique representation
\[
    x=\delta_1 a_1\vee\cdots\vee \delta_n a_n,
    \qquad
    \delta_i\in\{0,1\},
\]
where \(\delta_i a_i\) means \(a_i\) if \(\delta_i=1\) and \(0\) if \(\delta_i=0\). This ordering of the atoms induces an anti-lexicographic ordering on \(B\): for
\[
    x=\delta_1 a_1\vee\cdots\vee \delta_n a_n
    \qquad\text{and}\qquad
    y=\epsilon_1 a_1\vee\cdots\vee \epsilon_n a_n,
\]
we declare that \(x\prec^{\bfb}y\) if at the largest index \(i\) for which \(\delta_i\neq \epsilon_i\), one has \(\delta_i<\epsilon_i\).
An ordering obtained in this way will be called a \emph{natural ordering} of \(\bfb\).

Let \(\mathcal{OBA}\) denote the class of all finite Boolean algebras equipped with a natural ordering.

\begin{proposition}
    \label{prop:oba-fraisse-order-class}
    The class \(\mathcal{OBA}\) is a Fraïssé order class.
\end{proposition}

\begin{proof}
    Heredity is immediate. If \(\bfb\) is a finite Boolean algebra with a natural ordering and \(\bfc\leq \bfb\) is a subalgebra, then the atoms of \(\bfc\) are unions of consecutive blocks of atoms of \(\bfb\), and the induced ordering on \(\bfc\) is again natural.

    Joint embedding is clear. For amalgamation, one uses the atom structure of finite Boolean algebras: embeddings of finite Boolean algebras are determined by the images of atoms, and an amalgam may be constructed by refining the atom decompositions of the two larger algebras over the common subalgebra. Ordering the new atoms compatibly with the given orderings yields a natural ordering on the amalgam. Hence \(\mathcal{OBA}\) is a Fraïssé order class.
\end{proof}

The Ramsey property for naturally ordered finite Boolean algebras is equivalent to the Dual Ramsey Theorem of Graham and Rothschild~\cite{GR}. Therefore the criterion from Chapter~\ref{chap:extreme-amenability-ramsey} applies once again.

\begin{theorem}
    \label{thm:ordered-boolean-algebra-ea}
    Let \(\bfb_\infty^{*}=\Flim(\mathcal{OBA})\) be the countable atomless Boolean algebra with its canonical natural ordering. Then \(\Aut(\bfb_\infty^{*})\) is extremely amenable.
\end{theorem}

\begin{proof}
    By Proposition~\ref{prop:oba-fraisse-order-class}, the class \(\mathcal{OBA}\) is a Fraïssé order class, and by the Dual Ramsey Theorem it has the Ramsey property. Corollary~\ref{cor:fraisse-order-class-ea} therefore gives that \(\Aut(\bfb_\infty^{*})\) is extremely amenable.
\end{proof}

This example is parallel to the vector-space case. The ordering is not chosen arbitrarily, but is determined by an ordering of the atoms, and this compatibility is what places the class within the Ramsey framework.

\section{Metric spaces}

We conclude with metric spaces, which provide one of the most striking applications of the correspondence developed in Chapter~\ref{chap:extreme-amenability-ramsey}.

Let \(L_0=\{R_q : q\in \Q_{>0}\}\), where each \(R_q\) is a binary relation symbol. A metric space \((X,d)\) determines an \(L_0\)-structure \(\bfx=\langle X,(R_q^{\bfx})_{q\in \Q_{>0}}\rangle\) by setting
\[
    R_q^{\bfx}(x,y)
    \quad\Longleftrightarrow\quad
    d(x,y)<q.
\]
In this way finite metric spaces with rational distances may be viewed as finite relational structures.

Let \(\mathcal{M}_{\Q}\) denote the class of all finite metric spaces with rational distances. This is a Fraïssé class; its Fraïssé limit will be denoted by \(\bfu_{\Q}\) and called the \emph{rational Urysohn space}. Let \(\mathcal{OM}_{\Q}\) be the class of all finite linearly ordered metric spaces with rational distances. Since \(\mathcal{M}_{\Q}\) has strong amalgamation and strong joint embedding, it follows that \(\mathcal{OM}_{\Q}\) is a Fraïssé order class.

The Fraïssé limit
\[
    \bfu_{\Q}^{<}=\Flim(\mathcal{OM}_{\Q})
\]
will be called the \emph{ordered rational Urysohn space}. A result of Ne\v{s}et\v{r}il implies that \(\mathcal{OM}_{\Q}\) has the Ramsey property~\cite{N04}. Hence the criterion from Chapter~\ref{chap:extreme-amenability-ramsey} yields the following.

\begin{theorem}
    \label{thm:ordered-rational-urysohn-ea}
    The automorphism group of the ordered rational Urysohn space is extremely amenable.
\end{theorem}

\begin{proof}
    The class \(\mathcal{OM}_{\Q}\) is a Fraïssé order class with the Ramsey property, so the conclusion follows from Corollary~\ref{cor:fraisse-order-class-ea}.
\end{proof}

We now explain how this yields Pestov's theorem on the full Urysohn space.

Let \(\U\) denote the Urysohn space, that is, the unique complete separable metric space that is universal for separable metric spaces and ultrahomogeneous for isometries. Urysohn showed that the completion of \(\bfu_{\Q}\) is isometric to \(\U\)~\cite{U}, so we may regard its underlying metric space \(U_{\Q}\) as a dense subspace of \(\U\). Uspenskij proved that \(\Iso(\U)\), equipped with the pointwise convergence topology, is a universal Polish group~\cite{Usp90}. Pestov showed that this group is extremely amenable~\cite{P02}.

\begin{theorem}[Pestov]
    \label{thm:pestov-urysohn}
    The topological group \(\Iso(\U)\) is extremely amenable.
\end{theorem}

\begin{proof}
    We first record a simple general fact.

    \begin{lemma}
        \label{lem:dense-image-ea}
        Let \(\pi\colon G\to H\) be a continuous homomorphism of topological groups with dense range. If \(G\) is extremely amenable, then \(H\) is extremely amenable.
    \end{lemma}

    \begin{proof}
        Let \(X\) be a compact \(H\)-flow. Via \(\pi\), this becomes a \(G\)-flow, so by extreme amenability there is \(x\in X\) fixed by every element of \(G\). Since \(\pi(G)\) is dense in \(H\) and the action is continuous, \(x\) is fixed by every element of \(H\).
    \end{proof}

    Let \(\bfu_{\Q}^{<}=\langle \bfu_{\Q},\prec^{\bfu_{\Q}^{<}}\rangle\) be the ordered rational Urysohn space. Every automorphism of \(\bfu_{\Q}^{<}\) is in particular an isometry of \(\bfu_{\Q}\), and therefore extends uniquely to an isometry of the completion \(\U\). Thus we obtain a continuous homomorphism
    \[
        \Phi\colon \Aut(\bfu_{\Q}^{<})\to \Iso(\U).
    \]
    By Theorem~\ref{thm:ordered-rational-urysohn-ea} and Lemma~\ref{lem:dense-image-ea}, it is enough to show that \(\Phi(\Aut(\bfu_{\Q}^{<}))\) is dense in \(\Iso(\U)\).

    Fix \(h\in \Iso(\U)\), \(x_1,\dots,x_n\in \U\), and \(\epsilon>0\). We must find \(g\in \Aut(\bfu_{\Q}^{<})\) such that
    \[
        d(\Phi(g)(x_i),h(x_i))<\epsilon
        \qquad\text{for } i=1,\dots,n.
    \]

    For this we use the following approximation lemma.

    \begin{lemma}
        \label{lem:ordered-rational-approximation}
        Let \(x_1,\dots,x_n,y_1,\dots,y_n\in \U\) be such that \(x_i\mapsto y_i\) for \(i=1,\dots,n\) is an isometry between the two finite subspaces they generate. Then for every \(\epsilon>0\) there exist points
        \[
            x_1',\dots,x_n',y_1',\dots,y_n'\in U_{\Q}
        \]
        such that:
        \begin{enumerate}
            \item[(i)] the map \(x_i'\mapsto y_i'\) is an isometry;
            \item[(ii)] this isometry is order-preserving with respect to \(\prec^{\bfu_{\Q}^{<}}\);
            \item[(iii)]
                  \[
                      d(x_i,x_i')<\epsilon
                      \qquad\text{and}\qquad
                      d(y_i,y_i')<\epsilon
                      \qquad (i=1,\dots,n).
                  \]
        \end{enumerate}
    \end{lemma}

    \begin{proof}
        This is proved by induction on \(n\), using the density of \(U_{\Q}\) in \(\U\) together with the one-point extension property of the rational Urysohn space. At each stage one chooses rational points sufficiently close to the given ones and realises the required rational distance pattern so that the resulting partial isometry is order-preserving.
    \end{proof}

    Apply Lemma~\ref{lem:ordered-rational-approximation} to the tuples
    \[
        x_1,\dots,x_n
        \qquad\text{and}\qquad
        h(x_1),\dots,h(x_n).
    \]
    We obtain points \(x_1',\dots,x_n',y_1',\dots,y_n'\in U_{\Q}\) such that the map \(x_i'\mapsto y_i'\) is an order-preserving isometry and
    \[
        d(x_i,x_i')<\epsilon/2,
        \qquad
        d(h(x_i),y_i')<\epsilon/2
        \qquad (i=1,\dots,n).
    \]

    Since \(\bfu_{\Q}^{<}\) is ultrahomogeneous, this finite order-preserving isometry extends to some \(g\in \Aut(\bfu_{\Q}^{<})\).
    Then \(\Phi(g)\) extends the same map on \(U_{\Q}\), so for each \(i\),
    \[
        \begin{aligned}
            d(\Phi(g)(x_i),h(x_i))
             & \leq d(\Phi(g)(x_i),\Phi(g)(x_i')) + d(\Phi(g)(x_i'),h(x_i)) \\
             & = d(x_i,x_i') + d(y_i',h(x_i))                               \\
             & < \epsilon.
        \end{aligned}
    \]
    Thus \(\Phi(g)\) belongs to the prescribed neighbourhood of \(h\), and the range of \(\Phi\) is dense in \(\Iso(\U)\). By Lemma~\ref{lem:dense-image-ea}, \(\Iso(\U)\) is extremely amenable.
\end{proof}

The metric-space case is the culmination of this chapter. It shows that the criterion from Chapter~\ref{chap:extreme-amenability-ramsey} applies not only to combinatorial structures such as graphs, hypergraphs, vector spaces, and Boolean algebras, but also to metric structures, and it leads to one of the most important examples of an extremely amenable Polish group.

The examples in this chapter show that the criterion from Chapter~\ref{chap:extreme-amenability-ramsey} applies to a broad range of Fraïssé order classes. In each case, the Ramsey property of a suitable ordered class yields the extreme amenability of the automorphism group of its Fraïssé limit. In the next chapter we turn from extreme amenability to universal minimal flows and ordered expansions, where the full scope of the Kechris--Pestov--Todorcevic correspondence becomes visible.

% !TEX root = main.tex
\chapter{Reasonable Ordered Expansions}
\label{chap:reasonable-expansions}

In the previous chapters, ordered Fraïssé classes were used to obtain extremely amenable automorphism groups and to prepare the description of universal minimal flows. To pass from these ordered expansions back to the underlying unordered structures, one must understand how an ordered class is related to the class obtained by forgetting the linear order. The aim of this chapter is to isolate the conditions under which this passage behaves well.

We first introduce the notion of a reasonable ordered expansion and show that, for a Fraïssé order class \(\cK\), reasonableness is exactly the condition ensuring that the reduct class \(\cK_0\) is again a Fraïssé class and that the reduct of the Fraïssé limit of \(\cK\) is the Fraïssé limit of \(\cK_0\). We then discuss freely ordered expansions and strong amalgamation, and conclude with a brief remark on order-forgetful expansions, where Ramsey properties pass directly between the ordered and unordered settings.

\section{Reducts and expansions}

Let \(L\) be a signature containing a distinguished binary relation symbol \(<\), and let \(L_0=L\setminus\{<\}\). If \(\bfa\) is an \(L\)-structure with underlying set \(A\), we write \(\bfa_0=\bfa\!\restriction\!L_0\) for the reduct of \(\bfa\) to the signature \(L_0\). Thus \(\bfa_0\) is obtained from \(\bfa\) by forgetting the interpretation of the symbol \(<\). Conversely, if \(\bfa_0\) is an \(L_0\)-structure and \(\prec^{\bfa}\) is a linear ordering of \(A\), we write \(\langle \bfa_0,\prec^{\bfa}\rangle\) for the corresponding \(L\)-expansion of \(\bfa_0\) in which \(<\) is interpreted by \(\prec^{\bfa}\).

If \(\cK\) is a class of finite \(L\)-structures, we write \(\cK_0=\{\bfa_0 : \bfa\in\cK\}\) for the class of all reducts of members of \(\cK\) to the signature \(L_0\). In the situations that matter for us, \(\cK\) will be an order class, so each structure in \(\cK\) carries a linear ordering, while \(\cK_0\) is the corresponding unordered class obtained by forgetting that ordering.

The basic question is then the following. Suppose that \(\cK\) is a Fraïssé order class with Fraïssé limit \(\bff\). Under what conditions is the reduct class \(\cK_0\) again a Fraïssé class, and when is its Fraïssé limit simply the reduct \(\bff_0=\bff\!\restriction\!L_0\)?
The key condition is the notion of a reasonable ordered expansion, introduced in the next section.

\section{Reasonable Fraïssé order classes}

We now isolate the condition that controls when a Fraïssé order class behaves well after the linear order is forgotten. Intuitively, the issue is whether embeddings in the reduct can be lifted to embeddings in the ordered class by choosing a suitable ordering on the larger structure. This leads to the following definition.

\begin{definition}
    \label{def:reasonable}
    Let \(L\) be a signature with \(L\supseteq\{<\}\), and put \(L_0=L\setminus\{<\}\). Let \(\cK\) be a class of finite \(L\)-structures, and let \(\cK_0=\{\bfa_0 : \bfa\in\cK\}\). We say that \(\cK\) is \emph{reasonable} if whenever \(\bfa_0,\bfb_0\in\cK_0\), \(\pi\colon \bfa_0\hookrightarrow \bfb_0\) is an embedding, and \(\prec^{\bfa}\) is a linear ordering of \(A\) such that \(\bfa=\langle \bfa_0,\prec^{\bfa}\rangle\in\cK\), there exists a linear ordering \(\prec^{\bfb}\) of \(B\) such that \(\bfb=\langle \bfb_0,\prec^{\bfb}\rangle\in\cK\) and \(\pi\colon \bfa\hookrightarrow \bfb\) is an embedding.
\end{definition}

In other words, \(\cK\) is reasonable if every embedding in the reduct can be made compatible with the order after choosing a suitable ordering on the target structure.

The importance of this condition is that it exactly characterises when the reduct of a Fraïssé order class is again a Fraïssé class, with the expected Fraïssé limit.

\begin{proposition}
    \label{prop:reasonable-characterization}
    Let \(L\supseteq\{<\}\) be a signature, let \(L_0=L\setminus\{<\}\), and let \(\cK\) be a Fraïssé order class of finite \(L\)-structures. Put
    \[
        \cK_0=\{\bfa_0 : \bfa\in\cK\},
        \qquad
        \bff=\Flim(\cK),
        \qquad
        \bff_0=\bff\!\restriction\!L_0.
    \]
    Then the following are equivalent.
    \begin{enumerate}
        \item[(i)] \(\cK_0\) is a Fraïssé class and
              \[
                  \bff_0=\Flim(\cK_0).
              \]
        \item[(ii)] \(\cK\) is reasonable.
    \end{enumerate}
\end{proposition}

\begin{proof}
    Assume first that \(\cK_0\) is a Fraïssé class and that \(\bff_0=\Flim(\cK_0)\).
    We show that \(\cK\) is reasonable.

    Let \(\bfa_0,\bfb_0\in\cK_0\), let \(\pi\colon \bfa_0\hookrightarrow \bfb_0\) be an embedding, and let \(\prec^{\bfa}\) be a linear ordering of \(A\) such that \(\bfa=\langle \bfa_0,\prec^{\bfa}\rangle\in\cK\). Since \(\bff=\Flim(\cK)\), there is an embedding \(\varphi\colon \bfa\hookrightarrow \bff\). Forgetting the order, we may also regard \(\varphi\) as an embedding \(\varphi\colon \bfa_0\hookrightarrow \bff_0\). Because \(\bff_0\) is the Fraïssé limit of \(\cK_0\), it has the extension property. Hence there is an embedding \(\psi\colon \bfb_0\hookrightarrow \bff_0\) such that \(\psi\circ\pi=\varphi\).
    Now let \(\prec^{\bfb}\) be the linear ordering on \(B\) pulled back from \(\bff\) by \(\psi\), that is,
    \[
        x\prec^{\bfb} y
        \quad\Longleftrightarrow\quad
        \psi(x)\prec^{\bff}\psi(y).
    \]
    If we put \(\bfb=\langle \bfb_0,\prec^{\bfb}\rangle\), then \(\bfb\) is isomorphic to the substructure of \(\bff\) induced on \(\psi(B)\), so \(\bfb\in\Age(\bff)=\cK\). Moreover, since \(\psi\circ\pi=\varphi\) and \(\varphi\) preserves \(\prec^{\bfa}\), the map \(\pi\colon \bfa\hookrightarrow \bfb\) is an embedding. Thus \(\cK\) is reasonable.

    Conversely, assume that \(\cK\) is reasonable. We show first that \(\cK_0\) is a Fraïssé class.

    Since \(\cK_0=\Age(\bff_0)\),
    the class \(\cK_0\) is hereditary, countable up to isomorphism, and has structures of arbitrarily large cardinality. It therefore remains only to verify the amalgamation property.

    Let \(\bfa_0,\bfb_0,\bfc_0\in\cK_0\), and let \(f\colon \bfa_0\hookrightarrow \bfb_0\) and \(g\colon \bfa_0\hookrightarrow \bfc_0\) be embeddings. Choose a linear ordering \(\prec^{\bfa}\) on \(A\) such that \(\bfa=\langle \bfa_0,\prec^{\bfa}\rangle\in\cK\).
    Since \(\cK\) is reasonable, there exist linear orderings \(\prec^{\bfb}\) and \(\prec^{\bfc}\) on \(B\) and \(C\) such that \(\bfb=\langle \bfb_0,\prec^{\bfb}\rangle\in\cK\) and \(\bfc=\langle \bfc_0,\prec^{\bfc}\rangle\in\cK\), and both \(f\colon \bfa\hookrightarrow \bfb\) and \(g\colon \bfa\hookrightarrow \bfc\) are embeddings. Since \(\cK\) is a Fraïssé class, it has the amalgamation property. Hence there exist \(\bfd\in\cK\) and embeddings \(r\colon \bfb\hookrightarrow \bfd\) and \(s\colon \bfc\hookrightarrow \bfd\) such that \(r\circ f=s\circ g\). Passing to reducts, if \(\bfd_0=\bfd\!\restriction\!L_0\), then \(r\) and \(s\) are also embeddings \(r\colon \bfb_0\hookrightarrow \bfd_0\) and \(s\colon \bfc_0\hookrightarrow \bfd_0\), and still satisfy \(r\circ f=s\circ g\).
    Thus \(\cK_0\) has the amalgamation property, and therefore \(\cK_0\) is a Fraïssé class.

    It remains to show that
    \[
        \bff_0=\Flim(\cK_0).
    \]
    Since \(\Age(\bff_0)=\cK_0\),
    it is enough to verify that \(\bff_0\) is ultrahomogeneous, or equivalently that it has the extension property.

    Let \(\bfa_0,\bfb_0\in\cK_0\), let \(\pi\colon \bfa_0\hookrightarrow \bfb_0\) be an embedding, and let \(\varphi\colon \bfa_0\hookrightarrow \bff_0\) be an embedding. Pull back the order on \(\bff\) along \(\varphi\) and define a linear ordering \(\prec^{\bfa}\) on \(A\) by
    \[
        x\prec^{\bfa} y
        \quad\Longleftrightarrow\quad
        \varphi(x)\prec^{\bff}\varphi(y).
    \]
    Then \(\bfa=\langle \bfa_0,\prec^{\bfa}\rangle\in\cK\), and \(\varphi\) is in fact an embedding \(\varphi\colon \bfa\hookrightarrow \bff\).
    Since \(\cK\) is reasonable, there exists a linear ordering \(\prec^{\bfb}\) on \(B\) such that \(\bfb=\langle \bfb_0,\prec^{\bfb}\rangle\in\cK\) and \(\pi\colon \bfa\hookrightarrow \bfb\) is an embedding. Because \(\bff=\Flim(\cK)\), it has the extension property, so there is an embedding \(\psi\colon \bfb\hookrightarrow \bff\) such that \(\psi\circ\pi=\varphi\). Forgetting the order, \(\psi\) is also an embedding \(\psi\colon \bfb_0\hookrightarrow \bff_0\) with \(\psi\circ\pi=\varphi\).
    Thus \(\bff_0\) has the extension property. It follows that \(\bff_0\) is ultrahomogeneous and \(\Age(\bff_0)=\cK_0\),
    so
    \[
        \bff_0=\Flim(\cK_0).
    \]
\end{proof}
The proposition shows that for Fraïssé order classes, reasonableness is exactly the condition that allows one to pass from the ordered class back to its reduct without losing the Fraïssé structure. In the next section we consider an important source of reasonable order classes, namely freely ordered expansions.



\section[Freely ordered expansions and strong amalgamation]{Freely ordered expansions and\\strong amalgamation}

A particularly important source of ordered expansions is obtained by imposing no compatibility condition between the additional linear order and the underlying \(L_0\)-structure.

\begin{definition}
    Let \(\cK_0\) be a class of finite \(L_0\)-structures. We define
    \[
        \cK_0 * \cLO
        =
        \{\langle \bfa_0,\prec^{\bfa}\rangle : \bfa_0\in\cK_0 \text{ and } \prec^{\bfa} \text{ is a linear ordering of } A\}.
    \]
    We call \(\cK_0 * \cLO\) the \emph{freely ordered expansion} of \(\cK_0\).
\end{definition}

The point is that in a freely ordered expansion every linear order is allowed. In particular, such classes are automatically reasonable.

\begin{lemma}
    Let \(\cK=\cK_0 * \cLO\). Then \(\cK\) is reasonable.
\end{lemma}

\begin{proof}
    Let \(\bfa_0,\bfb_0\in\cK_0\), let \(\pi\colon \bfa_0\hookrightarrow \bfb_0\) be an embedding, and suppose that \(\bfa=\langle \bfa_0,\prec^{\bfa}\rangle\in\cK\).
    Choose any linear ordering \(\prec^{\bfb}\) of \(B\) extending the order transported from \(\prec^{\bfa}\) along \(\pi\). Then \(\bfb=\langle \bfb_0,\prec^{\bfb}\rangle\in\cK\),
    and \(\pi\colon \bfa\hookrightarrow \bfb\) is an embedding.
\end{proof}

For freely ordered expansions, amalgamation in the ordered class is controlled exactly by strong amalgamation in the reduct.
\begin{proposition}
    \label{prop:free-orders-sap}
    Let \(L\supseteq\{<\}\) be a signature, let \(L_0=L\setminus\{<\}\), let \(\cK_0\) be a class of finite \(L_0\)-structures, and put \(\cK=\cK_0 * \cLO\).
    Then the following are equivalent.
    \begin{enumerate}
        \item[(i)] \(\cK\) has the amalgamation property.
        \item[(ii)] \(\cK\) has the strong amalgamation property.
        \item[(iii)] \(\cK_0\) has the strong amalgamation property.
    \end{enumerate}
\end{proposition}

\begin{proof}
    The implication (ii) \(\Rightarrow\) (i) is immediate.

    We prove (i) \(\Rightarrow\) (iii). Let \(\bfa_0,\bfb_0,\bfc_0\in\cK_0\), and let \(f\colon \bfa_0\hookrightarrow \bfb_0\) and \(g\colon \bfa_0\hookrightarrow \bfc_0\) be embeddings. Choose a linear ordering \(\prec^{\bfa}\) on \(A\). Since \(\cK=\cK_0 * \cLO\),
    we may choose linear orderings \(\prec^{\bfb}\) on \(B\) and \(\prec^{\bfc}\) on \(C\) such that
    \[
        \bfa=\langle \bfa_0,\prec^{\bfa}\rangle,\qquad
        \bfb=\langle \bfb_0,\prec^{\bfb}\rangle,\qquad
        \bfc=\langle \bfc_0,\prec^{\bfc}\rangle
    \]
    all belong to \(\cK\), the maps \(f\colon \bfa\hookrightarrow \bfb\) and \(g\colon \bfa\hookrightarrow \bfc\) are embeddings, and in addition
    \[
        f(A)\prec^{\bfb} (B\setminus f(A))
        \qquad\text{and}\qquad
        (C\setminus g(A))\prec^{\bfc} g(A).
    \]
    Since \(\cK\) has the amalgamation property, there exist \(\bfd\in\cK\) and embeddings \(r\colon \bfb\hookrightarrow \bfd\) and \(s\colon \bfc\hookrightarrow \bfd\) such that \(r\circ f=s\circ g\).
    We claim that
    \[
        r(B)\cap s(C)=r(f(A)).
    \]
    Suppose, towards a contradiction, that there is \(d\in r(B)\cap s(C)\setminus r(f(A))\).
    Since \(d\in r(B\setminus f(A))\), we have
    \[
        r(f(A))\prec^{\bfd} d.
    \]
    On the other hand, since \(d\in s(C\setminus g(A))\), we have
    \[
        d\prec^{\bfd} s(g(A))=r(f(A)).
    \]
    This is impossible. Hence the amalgam is strong in the reduct, and \(\cK_0\) has the strong amalgamation property.

    Finally, we prove (iii) \(\Rightarrow\) (ii). Let \(\bfa,\bfb,\bfc\in\cK\), and let \(f\colon \bfa\hookrightarrow \bfb\) and \(g\colon \bfa\hookrightarrow \bfc\) be embeddings. Forgetting the order, we obtain embeddings \(f\colon \bfa_0\hookrightarrow \bfb_0\) and \(g\colon \bfa_0\hookrightarrow \bfc_0\). By strong amalgamation in \(\cK_0\), there exist \(\bfd_0\in\cK_0\) and embeddings \(r\colon \bfb_0\hookrightarrow \bfd_0\) and \(s\colon \bfc_0\hookrightarrow \bfd_0\)
    such that
    \[
        r\circ f=s\circ g
        \qquad\text{and}\qquad
        r(B)\cap s(C)=r(f(A)).
    \]
    Since the images of \(\bfb_0\) and \(\bfc_0\) intersect only in the common copy of \(\bfa_0\), and the orders on \(\bfb\) and \(\bfc\) agree there, the two induced linear orders on \(r(B)\subseteq D\) and \(s(C)\subseteq D\) can be amalgamated to a linear ordering \(\prec^{\bfd}\) of \(D\). If we now put \(\bfd=\langle \bfd_0,\prec^{\bfd}\rangle\), then \(\bfd\in\cK\) because \(\cK=\cK_0 * \cLO\), and the maps \(r\colon \bfb\hookrightarrow \bfd\) and \(s\colon \bfc\hookrightarrow \bfd\) are embeddings witnessing strong amalgamation in \(\cK\).
\end{proof}
The same argument gives the corresponding statement for strong joint embedding.

\begin{corollary}
    \label{cor:free-orders-sjep}
    Let \(L\), \(L_0\), \(\cK_0\), and \(\cK\) be as in Proposition~\ref{prop:free-orders-sap}. Then \(\cK\) has the strong joint embedding property if and only if \(\cK_0\) has the strong joint embedding property.
\end{corollary}

\begin{proof}
    Assume first that \(\cK\) has the strong joint embedding property. Let \(\bfb_0,\bfc_0\in\cK_0\). Choose arbitrary linear orderings \(\prec^{\bfb}\) and \(\prec^{\bfc}\) on \(B\) and \(C\), respectively, and form \(\bfb=\langle \bfb_0,\prec^{\bfb}\rangle\) and \(\bfc=\langle \bfc_0,\prec^{\bfc}\rangle\) in \(\cK\). By strong joint embedding in \(\cK\), there exist \(\bfd\in\cK\) and embeddings \(r\colon \bfb\hookrightarrow \bfd\) and \(s\colon \bfc\hookrightarrow \bfd\)
    with disjoint images. Passing to reducts shows that \(\cK_0\) has the strong joint embedding property.

    Conversely, assume that \(\cK_0\) has the strong joint embedding property. Let \(\bfb,\bfc\in\cK\). Then there exist \(\bfd_0\in\cK_0\) and embeddings \(r\colon \bfb_0\hookrightarrow \bfd_0\) and \(s\colon \bfc_0\hookrightarrow \bfd_0\)
    with disjoint images. Since the images are disjoint, the linear orders on \(\bfb\) and \(\bfc\) may be combined into a linear ordering \(\prec^{\bfd}\) of \(D\). Setting \(\bfd=\langle \bfd_0,\prec^{\bfd}\rangle\), we have \(\bfd\in\cK\), and the embeddings \(r\colon \bfb\hookrightarrow \bfd\) and \(s\colon \bfc\hookrightarrow \bfd\) have disjoint images. Thus \(\cK\) has the strong joint embedding property.
\end{proof}
In particular, if \(\cK_0\) is a Fraïssé class with the strong amalgamation property, then \(\cK_0 * \cLO\) is a reasonable Fraïssé order class. Many of the ordered graph, hypergraph, and metric-space classes considered earlier arise in exactly this way.

\section{Order-forgetful expansions and Ramsey transfer}

We conclude with a special situation in which the passage between an ordered class and its reduct is particularly simple. Roughly speaking, this happens when the additional order does not create new isomorphism types beyond those already present in the reduct.

\begin{definition}
    \label{def:order-forgetful}
    Let \(L\supseteq\{<\}\) be a signature, let \(L_0=L\setminus\{<\}\), and let \(\cK\) be a class of finite \(L\)-structures. We say that \(\cK\) is \emph{order-forgetful} if whenever \(\bfa,\bfb\in\cK\) satisfy \(\bfa_0\cong \bfb_0\), it follows that \(\bfa\cong \bfb\).
    Equivalently, for every isomorphism type in the reduct class \(\cK_0\), there is at most one isomorphism type in \(\cK\) lying above it.
\end{definition}

Thus the added order carries no additional isomorphism-type information.
Naturally ordered finite-dimensional vector spaces over a fixed finite field
are a basic example: once the underlying vector space is fixed, all natural
orderings are isomorphic.

For hereditary order-forgetful classes, the Ramsey property passes directly between the ordered and unordered settings.

\begin{proposition}
    \label{prop:order-forgetful-ramsey}
    Let \(L\supseteq\{<\}\) be a signature, let \(L_0=L\setminus\{<\}\), and let \(\cK\) be a hereditary order-forgetful class of finite \(L\)-structures. Put \(\cK_0=\{\bfa_0 : \bfa\in\cK\}\).
    Then \(\cK\) has the Ramsey property if and only if \(\cK_0\) has the Ramsey property.
\end{proposition}

\begin{proof}
    Assume first that \(\cK\) has the Ramsey property. Let \(\bfa_0\hookrightarrow \bfb_0\) be structures in \(\cK_0\), and let \(k\geq 2\). Choose an expansion \(\bfb\in\cK\) with reduct \(\bfb_0\). Since \(\bfa_0\hookrightarrow \bfb_0\), there is a copy \(\bfa_0'\leq \bfb_0\) of \(\bfa_0\). Let \(\bfa'\) be the ordered structure induced on \(\bfa_0'\) by the order of \(\bfb\). Because \(\cK\) is hereditary, we have \(\bfa'\in\cK\).
    Since \(\cK\) is order-forgetful and \(\bfa_0'\cong \bfa_0\), every expansion of \(\bfa_0\) in \(\cK\) is isomorphic to \(\bfa'\). In particular, if \(\bfa\in\cK\) is any expansion of \(\bfa_0\), then \(\bfa\cong \bfa'\leq \bfb\).
    Hence, replacing \(\bfa\) by an isomorphic copy if necessary, we may apply the
    Ramsey property of \(\cK\) to the pair \(\bfa\hookrightarrow \bfb\). Choose
    \(\bfc\in\cK\) such that
    \[
        \bfc\to (\bfb)^{\bfa}_k,
    \]
    and let \(\bfc_0=\bfc\!\restriction\!L_0\). We claim that
    \[
        \bfc_0\to (\bfb_0)^{\bfa_0}_k.
    \]

    Let
    \(c_0\colon \begin{pmatrix} \bfc_0 \\ \bfa_0 \end{pmatrix}\to\{1,\dots,k\}\)
    be a colouring. Every copy \(\bfa_0'\leq \bfc_0\) inherits from \(\bfc\) a linear ordering, and since \(\cK\) is hereditary, the induced ordered structure on \(\bfa_0'\) belongs to \(\cK\). Because \(\cK\) is order-forgetful and \(\bfa_0'\cong \bfa_0\), this induced ordered structure is isomorphic to \(\bfa\). Hence \(c_0\) induces a colouring
    \(c\colon \begin{pmatrix} \bfc \\ \bfa \end{pmatrix}\to\{1,\dots,k\}\)
    by setting \(c(\bfa')=c_0(\bfa_0')\),
    where \(\bfa_0'\) is the underlying reduct of the copy \(\bfa'\).

    Choose \(\bfb'\in \begin{pmatrix} \bfc \\ \bfb \end{pmatrix}\)
    homogeneous for \(c\). Then the reduct \(\bfb_0'\) is a copy of \(\bfb_0\) in \(\bfc_0\), and every copy of \(\bfa_0\) inside \(\bfb_0'\) inherits an ordered structure isomorphic to \(\bfa\). Since \(\bfb'\) is homogeneous for \(c\), it follows that \(\bfb_0'\) is homogeneous for \(c_0\). Thus
    \[
        \bfc_0\to (\bfb_0)^{\bfa_0}_k,
    \]
    and \(\cK_0\) has the Ramsey property.

    Conversely, assume that \(\cK_0\) has the Ramsey property. Let \(\bfa\hookrightarrow \bfb\) be structures in \(\cK\), and let \(k\geq 2\). Since \(\cK_0\) has the Ramsey property, there exists \(\bfc_0\in\cK_0\)
    such that
    \[
        \bfc_0\to (\bfb_0)^{\bfa_0}_k.
    \]
    Choose any expansion \(\bfc\in\cK\) with reduct \(\bfc_0\).

    We claim that
    \[
        \bfc\to (\bfb)^{\bfa}_k.
    \]
    Let \(c\colon \begin{pmatrix} \bfc \\ \bfa \end{pmatrix}\to\{1,\dots,k\}\)
    be a colouring. Define
    \(c_0\colon \begin{pmatrix} \bfc_0 \\ \bfa_0 \end{pmatrix}\to\{1,\dots,k\}\)
    by \(c_0(\bfa_0')=c(\bfa')\),
    where \(\bfa'\) is the ordered structure induced on \(\bfa_0'\) by the order of \(\bfc\). This is well defined because \(\cK\) is hereditary, and \(\bfa'\) is isomorphic to \(\bfa\) by order-forgetfulness.

    Choose \(\bfb_0'\in \begin{pmatrix} \bfc_0 \\ \bfb_0 \end{pmatrix}\)
    homogeneous for \(c_0\). Let \(\bfb'\) be the ordered structure induced on \(\bfb_0'\) by the order of \(\bfc\). Again, since \(\cK\) is hereditary and order-forgetful, we have \(\bfb'\cong \bfb\).
    Moreover, every copy of \(\bfa\) inside \(\bfb'\) has the same \(c\)-colour, because the corresponding copies of \(\bfa_0\) inside \(\bfb_0'\) all have the same \(c_0\)-colour. Hence \(\bfb'\) is homogeneous for \(c\), and therefore
    \[
        \bfc\to (\bfb)^{\bfa}_k.
    \]
    Thus \(\cK\) has the Ramsey property.
\end{proof}
The results of this chapter provide the structural link between an ordered Fraïssé class and its reduct. In the next chapter we use this link to study spaces of admissible orderings and to formulate the additional conditions under which ordered expansions determine the universal minimal flow.

% !TEX root = main.tex
\chapter{Universal Minimal Flows and the Ordering Property}
\label{chap:umf-ordering-property}
Throughout this and the remaining chapters, bold capital letters denote
structures, the corresponding plain capitals denote their underlying sets,
and calligraphic letters denote classes of structures.

In the previous chapter we studied reasonable ordered expansions and the passage from an ordered Fraïssé class to its reduct. We now return to the dynamical side of the Kechris--Pestov--Todorcevic correspondence~\cite{KPT05}. Starting from a reasonable Fraïssé order class \(\cK\), we associate to it a compact flow \(X_\cK\) of admissible orderings on the Fraïssé limit of the reduct class \(\cK_0\). The aim of this chapter is to show that the ordering property characterises the minimality of this flow, and that together with the Ramsey property it yields the universal minimal flow.

\section{Admissible orderings and the flow \texorpdfstring{\(X_\cK\)}{X K}}

Let \(L\supseteq\{<\}\) be a signature, and put \(L_0=L\setminus\{<\}\).
Let \(\cK\) be a reasonable Fraïssé order class in the language \(L\), and let
\(\cK_0=\cK\!\restriction\!L_0\) be its reduct to the language \(L_0\). By the results of Chapter~\ref{chap:reasonable-expansions},
the class \(\cK_0\) is a Fraïssé class, and if
\(\bff=\Flim(\cK)\) and \(\bff_0=\Flim(\cK_0)\), then
\(\bff_0=\bff\!\restriction\!L_0\). Let \(\prec_0\) denote the interpretation of the language symbol
\(<\) in \(\bff\), and put
\(G_0=\Aut(\bff_0)\).

Let \(\LO(F_0)\) denote the compact space of all linear orderings on the underlying set \(F_0\) of \(\bff_0\), equipped with the logic action of \(G_0\).
We define
\[
    X_\cK=\overline{G_0\cdot \prec_0}\subseteq \LO(F_0).
\]
Thus \(X_\cK\) is the orbit closure of the distinguished ordering \(\prec_0\) induced from the Fraïssé limit \(\bff\).

The first point is that membership in \(X_\cK\) can be recognised purely in terms of finite substructures.

\begin{proposition}
    \label{prop:admissible-orderings-characterization}
    Let \(\prec\) be a linear ordering on \(F_0\). Then \(\prec\in X_\cK\) if and only if for every finite substructure \(\bfb_0\) of \(\bff_0\), the ordered structure
    \(\langle \bfb_0,\prec\restriction B_0\rangle\) belongs to \(\cK\).
\end{proposition}

\begin{proof}
    Assume first that \(\prec\in X_\cK\), and let \(\bfb_0\) be a finite substructure
    of \(\bff_0\). Since \(\prec\in \overline{G_0\cdot \prec_0}\), there exists
    \(g\in G_0\) such that
    \(\prec\restriction B_0=(g\cdot \prec_0)\restriction B_0\). Let
    \(\bfa_0=g^{-1}(\bfb_0)\), which is again a finite substructure of \(\bff_0\),
    and let \(\bfa=\langle \bfa_0,\prec_0\restriction A_0\rangle\).
    Since \(\bfa\) is a finite substructure of \(\bff\), we have \(\bfa\in\cK\). Moreover, the map
    \(g\restriction A_0\colon A_0\to B_0\) is an isomorphism from \(\bfa\) onto
    \(\bfb=\langle \bfb_0,\prec\restriction B_0\rangle\).
    Hence \(\bfb\in\cK\).

    Conversely, assume that for every finite substructure \(\bfb_0\) of \(\bff_0\), the ordered structure
    \(\bfb=\langle \bfb_0,\prec\restriction B_0\rangle\) belongs to \(\cK\). We show that
    \(\prec\in X_\cK\).

    Fix a finite substructure \(\bfb_0\) of \(\bff_0\). Since \(\bfb\in\cK=\Age(\bff)\), there exists an embedding
    \(\pi\colon \bfb\hookrightarrow \bff\). Let \(\bfa=\pi(\bfb)\) and
    \(\bfa_0=\bfa\!\restriction\!L_0\).
    Then \(\pi\) is an isomorphism from \(\bfb\) onto \(\bfa\), and in particular an isomorphism from \(\bfb_0\) onto \(\bfa_0\). Since \(\bfa_0\) and \(\bfb_0\) are finite substructures of the ultrahomogeneous structure \(\bff_0\), there exists
    \(g\in \Aut(\bff_0)=G_0\) extending \(\pi^{-1}\). It follows that
    \(\prec\restriction B_0=(g\cdot \prec_0)\restriction B_0\).
    As \(\bfb_0\) was arbitrary, this shows that every finite restriction of \(\prec\) agrees with a finite restriction of some translate of \(\prec_0\). Therefore
    \(\prec\in \overline{G_0\cdot \prec_0}=X_\cK\).
\end{proof}

\begin{definition}
    \label{def:k-admissible-ordering}
    A linear ordering \(\prec\) on \(F_0\) is called \emph{\(\cK\)-admissible} if
    \(\prec\in X_\cK\).
    Equivalently, \(\prec\) is \(\cK\)-admissible if every finite substructure of \(\bff_0\), equipped with the induced ordering from \(\prec\), belongs to \(\cK\).
\end{definition}

In other words, the points of \(X_\cK\) are exactly the linear orderings on \(F_0\) whose finite restrictions are compatible with the class \(\cK\). Since \(X_\cK\) is a closed \(G_0\)-invariant subspace of the compact space \(\LO(F_0)\), it is a compact \(G_0\)-flow. In the next section we determine exactly when this flow is minimal.

\section{The ordering property and minimality of \texorpdfstring{\(X_\cK\)}{X K}}

We now introduce the combinatorial condition that characterises the minimality of the flow \(X_\cK\).

\begin{definition}
    \label{def:ordering-property}
    We say that \(\cK\) has the \emph{ordering property} if for every
    \(\bfa_0\in\cK_0\) there exists \(\bfb_0\in\cK_0\)
    such that whenever
    \[
        \bfa=\langle \bfa_0,\prec\rangle\in\cK
        \qquad\text{and}\qquad
        \bfb=\langle \bfb_0,\prec'\rangle\in\cK,
    \]
    one has
    \[
        \bfa\hookrightarrow \bfb.
    \]
\end{definition}

Thus the ordering property says that for every finite reduct \(\bfa_0\), there is a larger finite reduct \(\bfb_0\) such that every admissible ordering of \(\bfb_0\) contains every admissible ordering of \(\bfa_0\).

The classes considered in Chapter~\ref{chap:examples-extreme-amenability} provide both positive and negative examples. For instance, the freely ordered graph and hypergraph classes satisfy the ordering property, whereas the class of all linearly ordered finite equivalence relations does not; see~\cite{KPT05}.

The key step is the following finite reformulation of the condition
\(\prec_0\in\overline{G_0\cdot \prec}\).

\begin{lemma}
    \label{lem:orbit-closure-characterization}
    Let \(\prec\) be a linear ordering on \(F_0\). Then the following are equivalent.
    \begin{enumerate}
        \item[(i)]
              \[
                  \prec_0\in\overline{G_0\cdot \prec}.
              \]
        \item[(ii)] For every \(\bfa\in\cK\),
              there exists a finite substructure \(\bfc_0\) of \(\bff_0\) such that
              \[
                  \langle \bfc_0,\prec\restriction C_0\rangle\in\cK
                  \qquad\text{and}\qquad
                  \bfa\hookrightarrow \langle \bfc_0,\prec\restriction C_0\rangle.
              \]
    \end{enumerate}
    In particular, if \(\prec\in X_\cK\), then condition~(ii) says exactly
    that every member of \(\cK\) embeds into a finite substructure of the ordered
    expansion \(\langle \bff_0,\prec\rangle\).
\end{lemma}

\begin{proof}
    Assume first that
    \(\prec_0\in\overline{G_0\cdot \prec}\), and let \(\bfa\in\cK\).
    Since \(\cK=\Age(\bff)\), we may identify \(\bfa\) with a finite substructure of \(\bff\), so that
    \(\bfa=\langle \bfa_0,\prec_0\restriction A_0\rangle\)
    for some finite substructure \(\bfa_0\) of \(\bff_0\). By the assumption on \(\prec\), there exists
    \(g\in G_0\) such that
    \((g\cdot \prec)\restriction A_0=\prec_0\restriction A_0\). Let
    \(\bfc_0=g^{-1}(\bfa_0)\).
    Then
    \(g\colon \langle \bfc_0,\prec\restriction C_0\rangle\to \bfa\)
    is an isomorphism. Hence
    \(\langle \bfc_0,\prec\restriction C_0\rangle\in\cK\) and
    \(\bfa\hookrightarrow \langle \bfc_0,\prec\restriction C_0\rangle\).

    Conversely, assume (ii). To show that
    \(\prec_0\in\overline{G_0\cdot \prec}\),
    it suffices to verify that for every finite substructure \(\bfa_0\) of \(\bff_0\), there exists
    \(g\in G_0\) such that
    \((g\cdot \prec)\restriction A_0=\prec_0\restriction A_0\).
    Fix such a finite substructure \(\bfa_0\), and let
    \(\bfa=\langle \bfa_0,\prec_0\restriction A_0\rangle\in\cK\).
    By (ii), there exists a finite substructure \(\bfc_0\) of \(\bff_0\) such that
    \(\bfc=\langle \bfc_0,\prec\restriction C_0\rangle\in\cK\) and
    \(\bfa\hookrightarrow \bfc\).
    Let \(\pi\colon \bfa\hookrightarrow \bfc\) be an embedding. Forgetting the order,
    \(\pi\) is an embedding \(\pi\colon \bfa_0\hookrightarrow \bff_0\).
    Since \(\bff_0\) is ultrahomogeneous, \(\pi^{-1}\) extends to some
    \(g\in G_0\).
    For \(x,y\in A_0\) we then have
    \[
        x\,(g\cdot \prec)\,y
        \iff
        g^{-1}(x)\prec g^{-1}(y)
        \iff
        \pi(x)\prec \pi(y)
        \iff
        x\prec_0 y,
    \]
    because \(\pi\) is an embedding of ordered structures. Thus
    \((g\cdot \prec)\restriction A_0=\prec_0\restriction A_0\), as required.
\end{proof}
We can now characterise the minimality of \(X_\cK\).

\begin{theorem}
    \label{thm:minimality-ordering-property}
    The flow \(X_\cK\) is minimal if and only if \(\cK\) has the ordering property.
\end{theorem}

\begin{proof}
    Assume first that \(\cK\) has the ordering property. Let \(\prec\in X_\cK\).
    We show that \(\prec_0\in\overline{G_0\cdot \prec}\).
    By Lemma~\ref{lem:orbit-closure-characterization}, it is enough to verify condition (ii) of that lemma.

    So let \(\bfa\in\cK\), and write
    \(\bfa=\langle \bfa_0,\prec^{\bfa}\rangle\).
    By the ordering property, there exists \(\bfb_0\in\cK_0\)
    such that whenever
    \(\bfb=\langle \bfb_0,\prec^{\bfb}\rangle\in\cK\), one has
    \(\bfa\hookrightarrow \bfb\). Since \(\Age(\bff_0)=\cK_0\),
    we may realise \(\bfb_0\) as a finite substructure of \(\bff_0\). Because \(\prec\in X_\cK\), Proposition~\ref{prop:admissible-orderings-characterization} implies that
    \(\langle \bfb_0,\prec\restriction B_0\rangle\in\cK\).
    Hence
    \(\bfa\hookrightarrow \langle \bfb_0,\prec\restriction B_0\rangle\).
    By Lemma~\ref{lem:orbit-closure-characterization},
    \(\prec_0\in\overline{G_0\cdot \prec}\).

    Now let \(\prec\in X_\cK\).
    The orbit closure \(\overline{G_0\cdot \prec}\) is a closed \(G_0\)-invariant subset of \(X_\cK\) containing \(\prec_0\). Since
    \(X_\cK=\overline{G_0\cdot \prec_0}\),
    it follows that
    \(X_\cK\subseteq \overline{G_0\cdot \prec}\).
    The reverse inclusion is automatic, so
    \(\overline{G_0\cdot \prec}=X_\cK\).
    Thus every orbit is dense, and \(X_\cK\) is minimal.

    Conversely, assume that \(X_\cK\) is minimal. Let \(\bfa_0\in\cK_0\).
    We must find \(\bfb_0\in\cK_0\)
    witnessing the ordering property for \(\bfa_0\).

    There are only finitely many linear orderings on the finite set \(A_0\), hence only finitely many ordered structures of the form
    \(\langle \bfa_0,\prec\rangle\) that belong to \(\cK\). List them as
    \(\bfa^1,\dots,\bfa^m\).

    Fix \(\prec\in X_\cK\).
    Since \(X_\cK\) is minimal,
    \(\prec_0\in\overline{G_0\cdot \prec}\).
    Applying Lemma~\ref{lem:orbit-closure-characterization} to each \(\bfa^j\), we obtain finite substructures
    \[
        \bfc^{\,j}_{\prec,0}\leq \bff_0
        \qquad (j=1,\dots,m)
    \]
    such that
    \[
        \bfc^{\,j}_{\prec}
        =
        \langle \bfc^{\,j}_{\prec,0},\prec\restriction C^{\,j}_{\prec,0}\rangle\in\cK
    \]
    and \(\bfa^j\hookrightarrow \bfc^{\,j}_{\prec}\). Let
    \[
        \bfc_{\prec,0}
        =
        \left\langle
            C^{\,1}_{\prec,0}\cup\cdots\cup C^{\,m}_{\prec,0}
        \right\rangle_{\bff_0},
        \qquad
        C_{\prec,0}=\dom(\bfc_{\prec,0}).
    \]
    Since \(\prec\in X_\cK\), Proposition~\ref{prop:admissible-orderings-characterization} gives
    \(\langle \bfc_{\prec,0},\prec\restriction C_{\prec,0}\rangle\in\cK\),
    and by construction every \(\bfa^j\) embeds into this ordered structure.

    Now define
    \[
        U_{\prec}
        =
        \{\prec'\in X_\cK \mathrel{:} \prec'\restriction C_{\prec,0}=\prec\restriction C_{\prec,0}\}.
    \]
    This is an open neighbourhood of \(\prec\) in \(X_\cK\). As \(\prec\) ranges over \(X_\cK\), the sets \(U_{\prec}\) form an open cover of the compact space \(X_\cK\). Hence there exist
    \(\prec_1,\dots,\prec_n\in X_\cK\)
    such that
    \[
        X_\cK=U_{\prec_1}\cup\cdots\cup U_{\prec_n}.
    \]

    Choose a finite substructure \(\bfb_0\leq \bff_0\) containing each of the
    finite substructures
    \[
        \bfc_{\prec_1,0},\dots,\bfc_{\prec_n,0}.
    \]
    We claim that \(\bfb_0\) witnesses the ordering property for \(\bfa_0\).

    Let
    \[
        \bfa=\langle \bfa_0,\prec^{\bfa}\rangle\in\cK
        \qquad\text{and}\qquad
        \bfb=\langle \bfb_0,\prec^{\bfb}\rangle\in\cK.
    \]
    Since \(\bfb\in\Age(\bff)=\cK\), there exists a finite ordered substructure
    \(\widetilde{\bfb}\leq \bff\) isomorphic to \(\bfb\). Let
    \(h\colon \widetilde{\bfb}\to \bfb\) be an isomorphism, and let
    \(g\in G_0\) be an automorphism of \(\bff_0\) extending
    \(h^{-1}\colon B_0\to \widetilde{B}_0\). Then
    \(\widehat{\prec}=g\cdot \prec_0\)
    belongs to \(X_\cK\), and its restriction to \(B_0\) is exactly \(\prec^{\bfb}\).

    Since \(X_\cK=U_{\prec_1}\cup\cdots\cup U_{\prec_n}\),
    there exists \(i\in\{1,\dots,n\}\) such that
    \(\widehat{\prec}\in U_{\prec_i}\).
    Hence
    \(\widehat{\prec}\restriction C_{\prec_i,0}=\prec_i\restriction C_{\prec_i,0}\).
    By construction of \(C_{\prec_i,0}\), every ordered expansion of \(\bfa_0\) belonging to \(\cK\) embeds into
    \(\langle \bfc_{\prec_i,0},\prec_i\restriction C_{\prec_i,0}\rangle\).
    In particular,
    \[
        \bfa\hookrightarrow \langle \bfc_{\prec_i,0},\prec_i\restriction C_{\prec_i,0}\rangle
        =
        \langle \bfc_{\prec_i,0},\widehat{\prec}\restriction C_{\prec_i,0}\rangle.
    \]
    Since \(C_{\prec_i,0}\subseteq B_0\)
    and \(\widehat{\prec}\restriction B_0=\prec^{\bfb}\), the latter is a substructure of \(\bfb\). Therefore
    \(\bfa\hookrightarrow \bfb\).
    This proves that \(\cK\) has the ordering property.
\end{proof}
The theorem shows that the minimality of the flow \(X_\cK\) is controlled by a purely finite combinatorial condition on the class \(\cK\). In the next section we add the Ramsey property and obtain the universal minimal flow.

\section{Ramsey property and universal minimal flows}

We now add the Ramsey property and obtain the main dynamical consequence of the Kechris--Pestov--Todorcevic correspondence.

Let \(G=\Aut(\bff)\). Since \(\bff\) is an ordered expansion of \(\bff_0\),
we have \(G\leq G_0\),
and every element of \(G\) fixes the distinguished ordering \(\prec_0\).

The first result shows that, under the Ramsey property, the pointed flow
\((X_\cK,\prec_0)\)
is universal among \(G_0\)-ambits whose distinguished point is fixed by \(G\).

\begin{lemma}
    \label{lem:graph-uniqueness}
    Assume that \(\cK\) has the Ramsey property. Let \((X,x_0)\) be a \(G_0\)-ambit
    such that \(g\cdot x_0=x_0\) for all \(g\in G\).
    Let
    \[
        \Phi
        =
        \overline{\{(g\cdot \prec_0,g\cdot x_0):g\in G_0\}}
        \subseteq X_\cK\times X.
    \]
    If \((\prec,x),(\prec,x')\in \Phi\), then \(x=x'\).
\end{lemma}

\begin{proof}
    Suppose that \((\prec,x),(\prec,x')\in \Phi\) and \(x\neq x'\).
    Choose disjoint open neighbourhoods \(U\ni x\) and \(U'\ni x'\).
    Since \(X\) is compact Hausdorff, there exists an open neighbourhood \(D\) of the diagonal in \(X\times X\) such that
    \(D[U]\cap U'=\emptyset\), where
    \[
        D[U]=\{y\in X \mathrel{:} \text{there exists } z\in U \text{ with } (z,y)\in D\}.
    \]
    By replacing \(D\) with \(D\cap D^{-1}\),
    we may assume that \(D\) is symmetric.

    Because the action of \(G_0\) on the compact space \(X\) is continuous, there exists an open neighbourhood \(V\) of the identity in \(G_0\) such that
    \((v\cdot z,z)\in D\) for all \(z\in X\) and all \(v\in V\).
    Since \(G_0\) is non-Archimedean, we may choose a finite substructure
    \(\bfa_0\leq \bff_0\)
    such that the pointwise stabiliser
    \(G_{0,(A_0)}\subseteq V\).

    Because \((\prec,x),(\prec,x')\in \Phi\),
    there exist
    \(g,h\in G_0\)
    such that
    \(g\cdot x_0\in U\) and \(h\cdot x_0\in U'\),
    and
    \((g\cdot \prec_0)\restriction A_0=(h\cdot \prec_0)\restriction A_0=\prec\restriction A_0\).
    Consider the map
    \(\theta=h^{-1}g\colon g^{-1}(A_0)\to h^{-1}(A_0)\).
    Since
    \((g\cdot \prec_0)\restriction A_0=(h\cdot \prec_0)\restriction A_0\),
    the map \(\theta\) is an isomorphism of the finite ordered structures
    \(\langle g^{-1}(\bfa_0),\prec_0\restriction g^{-1}(A_0)\rangle\) and
    \(\langle h^{-1}(\bfa_0),\prec_0\restriction h^{-1}(A_0)\rangle\).
    Because \(\Age(\bff)=\cK\)
    and \(\bff\) is ultrahomogeneous, \(\theta\) extends to some
    \(\gamma\in G=\Aut(\bff)\).
    Then for every \(a\in A_0\),
    \(h\gamma g^{-1}(a)=a\),
    so
    \(v:=h\gamma g^{-1}\in G_{0,(A_0)}\subseteq V\).
    Since \(\gamma\in G\) fixes \(x_0\), we obtain
    \[
        h\cdot x_0=v\cdot g\cdot \gamma^{-1}\cdot x_0=v\cdot g\cdot x_0.
    \]
    Hence
    \((h\cdot x_0,g\cdot x_0)\in D\).
    Because \(D\) is symmetric, this implies
    \(h\cdot x_0\in D[U]\).
    But
    \(g\cdot x_0\in U\) and \(h\cdot x_0\in U'\),
    so
    \(h\cdot x_0\in D[U]\cap U'\),
    contrary to the choice of \(D\). Therefore
    \(x=x'\).
\end{proof}
\begin{theorem}
    \label{thm:ramsey-universal-ambit}
    Assume that \(\cK\) has the Ramsey property. Let \((X,x_0)\) be a \(G_0\)-ambit
    such that \(g\cdot x_0=x_0\) for all \(g\in G\).
    Then there exists a unique \(G_0\)-map \(\pi\colon X_\cK\to X\) such that
    \(\pi(\prec_0)=x_0\).
    In particular, \((X_\cK,\prec_0)\) is universal among \(G_0\)-ambits whose distinguished point is fixed by \(G\).
\end{theorem}

\begin{proof}
    Let
    \[
        \Phi
        =
        \overline{\{(g\cdot \prec_0,g\cdot x_0):g\in G_0\}}
        \subseteq X_\cK\times X.
    \]
    Since the projection \(p_1\colon X_\cK\times X\to X_\cK\)
    is a closed map and \(p_1(\Phi)\) contains the dense set
    \(G_0\cdot \prec_0\),
    we have
    \(p_1(\Phi)=X_\cK\). Thus for every \(\prec\in X_\cK\) there exists
    \(x\in X\) such that \((\prec,x)\in \Phi\).
    By Lemma~\ref{lem:graph-uniqueness}, this \(x\) is unique. Hence \(\Phi\) is the graph of a function
    \(\pi\colon X_\cK\to X\).

    The map
    \(p_1\restriction \Phi\colon \Phi\to X_\cK\)
    is a continuous bijection from a compact space to a Hausdorff space, hence a homeomorphism. Therefore
    \(\pi=p_2\circ (p_1\restriction \Phi)^{-1}\)
    is continuous.

    The set
    \(\{(g\cdot \prec_0,g\cdot x_0):g\in G_0\}\)
    is \(G_0\)-invariant, and so is its closure \(\Phi\). It follows that \(\pi\) is \(G_0\)-equivariant. Also,
    \((\prec_0,x_0)\in \Phi\),
    so
    \(\pi(\prec_0)=x_0\).

    To prove uniqueness, let
    \(\pi'\colon X_\cK\to X\)
    be another \(G_0\)-map with
    \(\pi'(\prec_0)=x_0\).
    For every \(g\in G_0\) we then have
    \[
        \pi'(g\cdot \prec_0)
        =g\cdot \pi'(\prec_0)
        =g\cdot x_0
        =\pi(g\cdot \prec_0).
    \]
    Since \(G_0\cdot \prec_0\)
    is dense in \(X_\cK\) and both maps are continuous, it follows that
    \(\pi'=\pi\).

    Finally, because \((X,x_0)\) is an ambit, the orbit
    \(G_0\cdot x_0\)
    is dense in \(X\). Since \(\pi\) is continuous and \(G_0\)-equivariant, its image is a closed \(G_0\)-invariant subset of \(X\) containing
    \(\pi(G_0\cdot \prec_0)=G_0\cdot x_0\).
    Hence
    \(\pi(X_\cK)=X\).
    Thus \((X_\cK,\prec_0)\) is universal among such ambits.
\end{proof}

We can now combine this with the results of Chapters~\ref{chap:extreme-amenability-ramsey} and~\ref{chap:reasonable-expansions} and the minimality criterion from the previous section.

\begin{corollary}
    \label{cor:ramsey-ordering-umf}
    Assume that \(\cK\) has both the Ramsey property and the ordering property. Then \(X_\cK\) is the universal minimal flow of \(G_0=\Aut(\bff_0)\).
\end{corollary}

\begin{proof}
    By Corollary~\ref{cor:fraisse-order-class-ea}, the Ramsey property of the Fraïssé order class \(\cK\) implies that
    \(G=\Aut(\bff)\)
    is extremely amenable. Let \(Y\) be a minimal compact \(G_0\)-flow. Restricting the action to the subgroup \(G\), we obtain a compact \(G\)-flow, so extreme amenability yields a point
    \(y_0\in Y\)
    fixed by every element of \(G\). Since \(Y\) is minimal as a \(G_0\)-flow, the orbit
    \(G_0\cdot y_0\)
    is dense in \(Y\), and therefore \((Y,y_0)\) is a \(G_0\)-ambit with \(G\)-fixed distinguished point.

    By Theorem~\ref{thm:ramsey-universal-ambit}, there exists a surjective
    \(G_0\)-map \(\pi\colon X_\cK\to Y\).
    On the other hand, by Theorem~\ref{thm:minimality-ordering-property}, the ordering property implies that \(X_\cK\) is minimal. Hence every minimal compact \(G_0\)-flow is a factor of \(X_\cK\), which is therefore the universal minimal flow of \(G_0\).
\end{proof}

The preceding corollary shows how the Ramsey property and admissible orderings determine the universal minimal flow of the automorphism group of the reduct. In the following chapters we turn to the converse direction, uniqueness questions for such expansions, and the role of Ramsey degrees in the computation of universal minimal flows.

% !TEX root = main.tex
\chapter{Examples of Universal Minimal Flows}
\label{chap:examples-umf}

In Chapter~\ref{chap:umf-ordering-property} we proved that if a reasonable Fra\"iss\'e order class $\k$ has both the Ramsey property and the ordering property, then the flow $X_\k$ of admissible orderings is the universal minimal flow of the automorphism group of the reduct. The aim of this chapter is to apply that theorem to concrete classes and to see how different kinds of universal minimal flows arise in practice.

The examples naturally divide into three types. In some cases every linear ordering is admissible, so the universal minimal flow is the full space of linear orderings. In other cases admissible orderings form a proper closed subspace determined by the underlying structure. Finally, there are examples in which the universal minimal flow is finite. We begin with the first phenomenon.

\section{Cases where the universal minimal flow is the full ordering space}

In the examples considered in this section, the ordered expansion is the free one: no compatibility condition is imposed between the additional linear order and the underlying structure. As a result, every linear ordering is admissible, so the flow $X_\k$ coincides with the full ordering space. Once the Ramsey and ordering properties are known, the universal minimal flow follows immediately from Chapter~\ref{chap:umf-ordering-property}.

For a countable structure $\bff_0$ with underlying set $F_0$, we write
$\mathrm{LO}(F_0)$ for the compact space of all linear orderings on $F_0$.

\begin{theorem}
\label{thm:umf-full-ordering-space}
The following groups have universal minimal flow equal to the full space of linear orderings on their underlying countable structure.
\begin{enumerate}
    \item[(i)] Let $\Gamma$ be the random graph. Then the universal minimal flow of $\Aut(\Gamma)$ is $\mathrm{LO}(\Gamma)$.

    \item[(ii)] For each $n\geq 3$, let $\Gamma_n$ be the random $K_n$-free graph. Then the universal minimal flow of $\Aut(\Gamma_n)$ is $\mathrm{LO}(\Gamma_n)$.

    \item[(iii)] The universal minimal flow of $S_\infty$ is $\mathrm{LO}(\mathbb N)$.

    \item[(iv)] Let $F$ be the Fra\"iss\'e limit of the class of all finite hypergraphs of a fixed finite relational type $L_0$. Then the universal minimal flow of $\Aut(F)$ is $\mathrm{LO}(F)$.

    \item[(v)] Let $\mathcal A$ be a family of finite irreducible hypergraphs of type $L_0$, and let $F_{\mathcal A}$ be the Fra\"iss\'e limit of $\mathrm{Forb}(\mathcal A)$. Then the universal minimal flow of $\Aut(F_{\mathcal A})$ is $\mathrm{LO}(F_{\mathcal A})$.

    \item[(vi)] Let $U_{\mathbb Q}$ be the rational Urysohn space. Then the universal minimal flow of $\Aut(U_{\mathbb Q})$ is $\mathrm{LO}(U_{\mathbb Q})$.
\end{enumerate}
\end{theorem}

\begin{proof}
In each case, let $\k_0$ be the relevant Fra\"iss\'e class, let
\(\bff_0=\Flim(\k_0)\) with underlying set $F_0$, and let
\(\k=\k_0*\mathcal{LO}\) be its free ordered expansion. By construction, every
linear ordering on a finite member of $\k_0$ gives a member of $\k$, so every
linear ordering on $F_0$ is $\k$-admissible. Therefore
\(X_\k=\mathrm{LO}(F_0)\).

It remains to verify that $\k$ has the Ramsey property and the ordering property. The Ramsey property was established in the preceding chapters: for graphs and $K_n$-free graphs by the theorem of Ne\v{s}et\v{r}il and R\"odl, for finite hypergraphs and $\mathcal A$-free hypergraphs by the corresponding hypergraph Ramsey theorems, and for finite rational metric spaces by Ne\v{s}et\v{r}il's theorem. The ordering property holds for these free ordered expansions. Hence Chapter~\ref{chap:umf-ordering-property} applies and gives that the universal minimal flow of the automorphism group of the reduct is exactly $X_\k$.

This yields (i), (ii), (iv), (v), and (vi). For (iii), take $\k_0$ to be the class of finite sets in the empty language. Then its Fra\"iss\'e limit is a countably infinite set, whose automorphism group is $S_\infty$, and again \(X_\k=\mathrm{LO}(\mathbb N)\).
Thus the universal minimal flow of $S_\infty$ is $\mathrm{LO}(\mathbb N)$.
\end{proof}

These examples show that when the ordered expansion is completely free, no information is lost in passing from admissible orderings to all linear orderings. In the next section we consider the opposite situation, where admissibility imposes genuine structural restrictions and the universal minimal flow becomes a proper closed subspace of the full ordering space.

\section{Structured spaces of admissible orderings}

We now consider examples in which not every linear ordering is admissible. In these cases the universal minimal flow is still obtained from Chapter~\ref{chap:umf-ordering-property}, but it is a proper closed subspace of the full ordering space, determined by compatibility conditions between the order and the underlying structure.

\begin{theorem}
\label{thm:umf-structured-ordering-spaces}
The following groups have universal minimal flow equal to a proper space of admissible orderings.
\begin{enumerate}
    \item[(i)] Let \(\bff_0=\langle \mathbb Q,E\rangle\) be the countable equivalence relation with infinitely many equivalence classes, each of which is dense in $\mathbb Q$. Let $X_{\mathcal K}$ be the space of all convex linear orderings on $\mathbb Q$, where $\mathcal K$ is the Fra\"iss\'e order class of all finite convexly ordered equivalence relations. Then the universal minimal flow of \(\Aut(\bff_0)\) is $X_{\mathcal K}$.

    \item[(ii)] Let $V_F$ be the countably infinite-dimensional vector space over a finite field $F$, and let $X_{\mathcal K}$ be the space of all natural orderings on $V_F$, where $\mathcal K$ is the class of naturally ordered finite-dimensional vector spaces over $F$. Then the universal minimal flow of \(\Aut(V_F)\) is $X_{\mathcal K}$.

    \item[(iii)] Let $\bfb_\infty$ be the countable atomless Boolean algebra, and let $X_{\mathcal K}$ be the space of all natural orderings on $\bfb_\infty$, where $\mathcal K$ is the class of naturally ordered finite Boolean algebras. Then the universal minimal flow of \(\Aut(\bfb_\infty)\) is $X_{\mathcal K}$.
\end{enumerate}
\end{theorem}

\begin{proof}
In each case, the relevant ordered class $\mathcal K$ was introduced in Chapter~\ref{chap:examples-extreme-amenability}. We verify that it fits the framework of Chapter~\ref{chap:umf-ordering-property}.

For (i), let $\mathcal K$ be the class of finite convexly ordered equivalence relations. By Chapter~\ref{chap:examples-extreme-amenability}, $\mathcal K$ is a Fra\"iss\'e order class with the Ramsey property, and its reduct $\mathcal K_0$ is the class of finite equivalence relations. Moreover, Chapter~\ref{chap:examples-extreme-amenability} showed that arbitrary linear orderings do not work here: the admissible orderings are exactly the convex ones. Thus the flow $X_{\mathcal K}$ is the space of convex linear orderings on $\mathbb Q$. Since $\mathcal K$ also has the ordering property, Chapter~\ref{chap:umf-ordering-property} implies that $X_{\mathcal K}$ is the universal minimal flow of $\Aut(\bff_0)$.

For (ii), let $\mathcal K$ be the class of naturally ordered finite-dimensional vector spaces over $F$. By Chapter~\ref{chap:examples-extreme-amenability}, this is a Fra\"iss\'e order class with the Ramsey property, and its reduct is the class of finite-dimensional vector spaces over $F$. The admissible orderings are precisely the natural orderings induced by ordered bases. Hence the corresponding flow $X_{\mathcal K}$ is the space of natural orderings on $V_F$. Since $\mathcal K$ has the ordering property, Chapter~\ref{chap:umf-ordering-property} yields that $X_{\mathcal K}$ is the universal minimal flow of $\Aut(V_F)$.

For (iii), let $\mathcal K$ be the class of naturally ordered finite Boolean algebras. Again, Chapter~\ref{chap:examples-extreme-amenability} shows that $\mathcal K$ is a Fra\"iss\'e order class with the Ramsey property, with reduct the class of finite Boolean algebras. The admissible orderings are precisely the natural orderings induced by orderings of the atoms. Therefore $X_{\mathcal K}$ is the space of natural orderings on $\bfb_\infty$. Since $\mathcal K$ also has the ordering property, Chapter~\ref{chap:umf-ordering-property} implies that $X_{\mathcal K}$ is the universal minimal flow of $\Aut(\bfb_\infty)$.
\end{proof}

The significance of these examples is that the universal minimal flow is no longer the full space of all linear orderings. Instead, admissibility imposes a structural constraint: convexity with respect to the equivalence classes, or naturality with respect to a basis or an atomic decomposition. In the Boolean algebra case, this space of admissible orderings admits a more concrete topological description, which we discuss next.

\section{Boolean algebras and maximal chains}

We now return to the Boolean algebra example and record a more concrete description of the corresponding universal minimal flow. Besides the space of natural orderings on the countable atomless Boolean algebra, there is a second model in terms of maximal chains of closed subsets of the Cantor space. The two descriptions are equivalent.

Let \(\mathcal C=2^{\mathbb N}\) be the Cantor space, and let \(\Phi(\mathcal C)\) denote the compact space of all maximal chains of closed subsets of $\mathcal C$, ordered by inclusion. The group \(\mathrm{Homeo}(\mathcal C)\) acts continuously on $\Phi(\mathcal C)$ by \(g\cdot \mathcal M=\{g(M):M\in\mathcal M\}\).

By Stone duality, the countable atomless Boolean algebra $\bfb_\infty$ may be identified with the Boolean algebra of clopen subsets of $\mathcal C$, and under this identification the group \(\Aut(\bfb_\infty)\) is naturally isomorphic to \(\mathrm{Homeo}(\mathcal C)\).
The next theorem identifies the Boolean-algebra universal minimal flow from Theorem~\ref{thm:umf-structured-ordering-spaces}(iii) with the maximal-chain flow on $\mathcal C$.

\begin{theorem}
\label{thm:boolean-maximal-chains}
Under the natural identification
\[
\Aut(\bfb_\infty)\cong \mathrm{Homeo}(2^{\mathbb N}),
\]
the universal minimal flow of $\Aut(\bfb_\infty)$ is isomorphic to
\[
\Phi(2^{\mathbb N}).
\]
Equivalently, the space of natural orderings on $\bfb_\infty$ is $\,\Aut(\bfb_\infty)$-equivariantly homeomorphic to the space of maximal chains of closed subsets of $2^{\mathbb N}$.
\end{theorem}

\begin{proof}
We sketch the construction. By Theorem~\ref{thm:umf-structured-ordering-spaces}(iii), the universal minimal flow of $\Aut(\bfb_\infty)$ is the space $X_{\mathcal K}$ of natural orderings on $\bfb_\infty$.

Given a maximal chain \(\mathcal M\in \Phi(2^{\mathbb N})\), one obtains a linear ordering on the atoms of each finite subalgebra of $\bfb_\infty$ by comparing the unique points at which the members of $\mathcal M$ first separate those atoms. Passing to the induced anti-lexicographic orderings, this defines a natural ordering on every finite subalgebra, and these finite orderings are coherent. Hence $\mathcal M$ determines an element of $X_{\mathcal K}$.

Conversely, let \(\prec\in X_{\mathcal K}\).
The ordering $\prec$ induces, on each finite clopen partition of $2^{\mathbb N}$, a linear ordering of its atoms compatible with refinement. From these finite orderings one reconstructs a unique maximal chain of closed subsets of $2^{\mathbb N}$ by taking the nested family of closed sets determined by initial segments of the induced orderings. This produces an element of $\Phi(2^{\mathbb N})$.

These two constructions are inverse to one another, continuous, and
equivariant for the action of $\Aut(\bfb_\infty)$. Therefore
\[
X_{\mathcal K}\cong \Phi(2^{\mathbb N})
\]
as $\Aut(\bfb_\infty)$-flows. Since $X_{\mathcal K}$ is the universal minimal flow, the same is true of $\Phi(2^{\mathbb N})$.
\end{proof}

This identification is useful because it shows that the same universal minimal flow may appear in two rather different forms: combinatorially, as a space of natural orderings on a Fra\"iss\'e limit, and topologically, as a space of maximal chains in the Cantor space. In the next section we consider a different phenomenon, where the universal minimal flow is finite.

\section{A finite universal minimal flow}

We conclude with an example in which the universal minimal flow is finite. This shows that, although many of the flows considered so far are infinite compact spaces of orderings, the universal minimal flow can also be a finite homogeneous space.

Fix $n\geq 1$, and let \(\langle \bbQ,<,E_n\rangle\) be the structure consisting of the rationals with their usual ordering together with an equivalence relation $E_n$ having exactly $n$ classes, each of which is dense in $\bbQ$. Let \(A_1,\dots,A_n\) be an enumeration of the $E_n$-classes, and put \(\bff_n=\langle \bbQ,<,A_1,\dots,A_n\rangle\). We write \(G_n=\Aut(\langle \bbQ,<,E_n\rangle)\) and \(H_n=\Aut(\bff_n)\).

Thus $G_n$ preserves the equivalence relation and may permute its $n$ dense classes, while $H_n$ preserves each class individually.

\begin{theorem}
\label{thm:finite-umf-gn}
For each $n\geq 1$, the subgroup $H_n$ is a clopen normal subgroup of finite index in $G_n$, and the natural action of $G_n$ on \(G_n/H_n\) is the universal minimal flow of $G_n$.
\end{theorem}

\begin{proof}
Each element of $G_n$ permutes the $n$ equivalence classes of $E_n$, so there is a natural continuous homomorphism \(\pi\colon G_n\to S_n\) sending an automorphism to the permutation it induces on \(\{A_1,\dots,A_n\}\).
Its kernel is exactly $H_n$. Hence $H_n$ is a normal subgroup of $G_n$ of finite index at most $n!$. Since the quotient is finite and discrete, $H_n$ is clopen.

We now show that $H_n$ is extremely amenable. Let $\mathcal{OP}_n$ be the class of all finite structures of the form
\[
\langle A,<^A,P_1^A,\dots,P_n^A\rangle,
\]
where $<^A$ is a linear ordering of $A$ and \(P_1^A,\dots,P_n^A\) form a partition of $A$ into pairwise disjoint unary predicates. This is a Fra\"iss\'e order class, and its Fra\"iss\'e limit is precisely \(\bff_n=\langle \bbQ,<,A_1,\dots,A_n\rangle\).
Moreover, $\mathcal{OP}_n$ has the Ramsey property; this is a special case of the general Ramsey theorem for finite ordered relational structures. Therefore, by Corollary~\ref{cor:fraisse-order-class-ea}, \(H_n=\Aut(\bff_n)\) is extremely amenable.

Let now $Y$ be a minimal compact $G_n$-flow. Restricting the action to the subgroup $H_n$, extreme amenability gives a point \(y_0\in Y\) fixed by every element of $H_n$. The orbit map \(g\mapsto g\cdot y_0\) is therefore constant on left cosets of $H_n$, so it factors through a continuous $G_n$-map \(\varphi\colon G_n/H_n\to Y\).
Since $Y$ is minimal, the orbit \(G_n\cdot y_0=\varphi(G_n/H_n)\) is dense in $Y$. But $G_n/H_n$ is finite, so its image is finite and hence closed. Therefore \(\varphi(G_n/H_n)=Y\),
and $\varphi$ is surjective. Thus every minimal compact $G_n$-flow is a factor of $G_n/H_n$.

Finally, the action of $G_n$ on $G_n/H_n$ is transitive, hence minimal. It follows that \(G_n/H_n\) is the universal minimal flow of $G_n$.
\end{proof}

For $n=1$, the quotient $G_1/H_1$ is trivial, so this recovers the extreme amenability of $\Aut(\bbQ,<)$. For $n\geq 2$, the universal minimal flow is finite and nontrivial. This contrasts sharply with the earlier examples, where the universal minimal flow was an infinite compact space of orderings.

\chapter{Universal Minimal Flows and the Ordering Property}

In the previous chapter we studied reasonable ordered expansions and the passage from an ordered Fra\"iss\'{e} class to its reduct. We now return to the dynamical side of the Kechris--Pestov--Todorcevic correspondence. Starting from a reasonable Fra\"iss\'{e} order class $\k$, we associate to it a compact flow $X_\k$ of admissible orderings on the Fra\"iss\'{e} limit of the reduct class $\k_0$. The aim of this chapter is to show that the ordering property characterizes the minimality of this flow, and that together with the Ramsey property it yields the universal minimal flow.

\section{Admissible orderings and the flow $X_\k$}

Let $L\supseteq\{<\}$ be a signature, and put
\[
    L_0=L\setminus\{<\}.
\]
Let $\k$ be a reasonable Fra\"iss\'{e} order class in the language $L$, and let
\[
    \k_0=\k|L_0
\]
be its reduct to the language $L_0$. By the results of Chapter~7, the class $\k_0$ is a Fra\"iss\'{e} class, and if
\[
    \bff={\rm Flim}(\k),
    \qquad
    \bff_0={\rm Flim}(\k_0),
\]
then
\[
    \bff_0=\bff|L_0.
\]
Let
\[
    <^{\bff}=\prec_0,
\]
and put
\[
    G_0={\rm Aut}(\bff_0).
\]

Let $\mathrm{LO}(F_0)$ denote the compact space of all linear orderings on the underlying set $F_0$ of $\bff_0$, equipped with the logic action of $G_0$. We define
\[
    X_\k=\overline{G_0\cdot \prec_0}\subseteq \mathrm{LO}(F_0).
\]
Thus $X_\k$ is the orbit closure of the distinguished ordering $\prec_0$ induced from the Fra\"iss\'{e} limit $\bff$.

The first point is that membership in $X_\k$ can be recognized purely in terms of finite substructures.

\begin{proposition}
    \label{prop:admissible-orderings-characterization}
    Let $\prec$ be a linear ordering on $F_0$. Then $\prec\in X_\k$ if and only if for every finite substructure $\bfb_0$ of $\bff_0$, the ordered structure
    \[
        \langle \bfb_0,\prec|B_0\rangle
    \]
    belongs to $\k$.
\end{proposition}

\begin{proof}
    Assume first that $\prec\in X_\k$, and let $\bfb_0$ be a finite substructure of $\bff_0$. Since
    \[
        \prec\in \overline{G_0\cdot \prec_0},
    \]
    there exists $g\in G_0$ such that
    \[
        \prec|B_0=(g\cdot \prec_0)|B_0.
    \]
    Let
    \[
        \bfa_0=g^{-1}(\bfb_0),
    \]
    which is again a finite substructure of $\bff_0$, and let
    \[
        \bfa=\langle \bfa_0,\prec_0|A_0\rangle.
    \]
    Since $\bfa$ is a finite substructure of $\bff$, we have $\bfa\in\k$. Moreover, the map
    \[
        g|A_0\colon A_0\to B_0
    \]
    is an isomorphism from $\bfa$ onto
    \[
        \bfb=\langle \bfb_0,\prec|B_0\rangle.
    \]
    Hence $\bfb\in\k$.

    Conversely, assume that for every finite substructure $\bfb_0$ of $\bff_0$, the ordered structure
    \[
        \bfb=\langle \bfb_0,\prec|B_0\rangle
    \]
    belongs to $\k$. We show that $\prec\in X_\k$.

    Fix a finite substructure $\bfb_0$ of $\bff_0$. Since $\bfb\in\k=\Age(\bff)$, there exists an embedding
    \[
        \pi\colon \bfb\to \bff.
    \]
    Let
    \[
        \bfa=\pi(\bfb),
        \qquad
        \bfa_0=\bfa|L_0.
    \]
    Then $\pi$ is an isomorphism from $\bfb$ onto $\bfa$, and in particular an isomorphism from $\bfb_0$ onto $\bfa_0$. Since $\bfa_0$ and $\bfb_0$ are finite substructures of the ultrahomogeneous structure $\bff_0$, there exists
    \[
        g\in {\rm Aut}(\bff_0)=G_0
    \]
    extending $\pi^{-1}$. It follows that
    \[
        \prec|B_0=(g\cdot \prec_0)|B_0.
    \]
    As $\bfb_0$ was arbitrary, this shows that every finite restriction of $\prec$ agrees with a finite restriction of some translate of $\prec_0$. Therefore
    \[
        \prec\in \overline{G_0\cdot \prec_0}=X_\k.
    \]
\end{proof}

\begin{definition}
    \label{def:k-admissible-ordering}
    A linear ordering $\prec$ on $F_0$ is called \emph{$\k$-admissible} if
    \[
        \prec\in X_\k.
    \]
    Equivalently, $\prec$ is $\k$-admissible if every finite substructure of $\bff_0$, equipped with the induced ordering from $\prec$, belongs to $\k$.
\end{definition}

In other words, the points of $X_\k$ are exactly the linear orderings on $F_0$ whose finite restrictions are compatible with the class $\k$. Since $X_\k$ is a closed $G_0$-invariant subspace of the compact space $\mathrm{LO}(F_0)$, it is a compact $G_0$-flow. In the next section we determine exactly when this flow is minimal.

\section{The ordering property and minimality of $X_\k$}

We now introduce the combinatorial condition that characterizes the minimality of the flow $X_\k$.

\begin{definition}
    \label{def:ordering-property}
    We say that $\k$ has the \emph{ordering property} if for every
    \[
        \bfa_0\in\k_0
    \]
    there exists
    \[
        \bfb_0\in\k_0
    \]
    such that whenever
    \[
        \bfa=\langle \bfa_0,\prec\rangle\in\k
        \qquad\text{and}\qquad
        \bfb=\langle \bfb_0,\prec'\rangle\in\k,
    \]
    one has
    \[
        \bfa\leq \bfb.
    \]
\end{definition}

Thus the ordering property says that for every finite reduct $\bfa_0$, there is a larger finite reduct $\bfb_0$ such that every admissible ordering of $\bfb_0$ contains every admissible ordering of $\bfa_0$.

The classes considered in Chapter~6 provide both positive and negative examples. For instance, the freely ordered graph and hypergraph classes satisfy the ordering property, whereas the class of all linearly ordered finite equivalence relations does not.

The key step is the following finite reformulation of the condition
\[
    \prec_0\in\overline{G_0\cdot \prec}.
\]

\begin{lemma}
    \label{lem:orbit-closure-characterization}
    Let $\prec$ be a linear ordering on $F_0$. Then the following are equivalent.
    \begin{enumerate}
        \item[(i)]
              \[
                  \prec_0\in\overline{G_0\cdot \prec}.
              \]
        \item[(ii)] For every
              \[
                  \bfa\in\k,
              \]
              there exists a finite substructure $\bfc_0$ of $\bff_0$ such that
              \[
                  \langle \bfc_0,\prec|C_0\rangle\in\k
                  \qquad\text{and}\qquad
                  \bfa\leq \langle \bfc_0,\prec|C_0\rangle.
              \]
    \end{enumerate}
    In particular, if $\prec\in X_\k$, then condition {\rm(ii)} says exactly that every member of $\k$ embeds into a finite substructure of the ordered structure
    \[
        \langle F_0,\prec\rangle.
    \]
\end{lemma}

\begin{proof}
    Assume first that
    \[
        \prec_0\in\overline{G_0\cdot \prec},
    \]
    and let
    \[
        \bfa\in\k.
    \]
    Since $\k=\Age(\bff)$, we may identify $\bfa$ with a finite substructure of $\bff$, so that
    \[
        \bfa=\langle \bfa_0,\prec_0|A_0\rangle
    \]
    for some finite substructure $\bfa_0$ of $\bff_0$. By the assumption on $\prec$, there exists
    \[
        g\in G_0
    \]
    such that
    \[
        (g\cdot \prec)|A_0=\prec_0|A_0.
    \]
    Let
    \[
        \bfc_0=g^{-1}(\bfa_0).
    \]
    Then
    \[
        g\colon \langle \bfc_0,\prec|C_0\rangle\to \bfa
    \]
    is an isomorphism. Hence
    \[
        \langle \bfc_0,\prec|C_0\rangle\in\k
        \qquad\text{and}\qquad
        \bfa\leq \langle \bfc_0,\prec|C_0\rangle.
    \]

    Conversely, assume (ii). To show that
    \[
        \prec_0\in\overline{G_0\cdot \prec},
    \]
    it suffices to verify that for every finite substructure $\bfa_0$ of $\bff_0$, there exists
    \[
        g\in G_0
    \]
    such that
    \[
        (g\cdot \prec)|A_0=\prec_0|A_0.
    \]
    Fix such a finite substructure $\bfa_0$, and let
    \[
        \bfa=\langle \bfa_0,\prec_0|A_0\rangle\in\k.
    \]
    By (ii), there exists a finite substructure $\bfc_0$ of $\bff_0$ such that
    \[
        \bfc=\langle \bfc_0,\prec|C_0\rangle\in\k
        \qquad\text{and}\qquad
        \bfa\leq \bfc.
    \]
    Let
    \[
        \pi\colon \bfa\to \bfc
    \]
    be an embedding. Forgetting the order, $\pi$ is an embedding
    \[
        \pi\colon \bfa_0\to \bff_0.
    \]
    Since $\bff_0$ is ultrahomogeneous, $\pi^{-1}$ extends to some
    \[
        g\in G_0.
    \]
    For $x,y\in A_0$ we then have
    \[
        x\,(g\cdot \prec)\,y
        \iff
        g^{-1}(x)\prec g^{-1}(y)
        \iff
        \pi(x)\prec \pi(y)
        \iff
        x\prec_0 y,
    \]
    because $\pi$ is an embedding of ordered structures. Thus
    \[
        (g\cdot \prec)|A_0=\prec_0|A_0,
    \]
    as required.
\end{proof}

We can now characterize the minimality of $X_\k$.

\begin{theorem}
    \label{thm:minimality-ordering-property}
    The flow $X_\k$ is minimal if and only if $\k$ has the ordering property.
\end{theorem}

\begin{proof}
    Assume first that $\k$ has the ordering property. Let
    \[
        \prec\in X_\k.
    \]
    We show that
    \[
        \prec_0\in\overline{G_0\cdot \prec}.
    \]
    By Lemma~\ref{lem:orbit-closure-characterization}, it is enough to verify condition (ii) of that lemma.

    So let
    \[
        \bfa\in\k.
    \]
    Write
    \[
        \bfa=\langle \bfa_0,\prec_A\rangle.
    \]
    By the ordering property, there exists
    \[
        \bfb_0\in\k_0
    \]
    such that whenever
    \[
        \bfb=\langle \bfb_0,\prec_B\rangle\in\k,
    \]
    one has
    \[
        \bfa\leq \bfb.
    \]
    Since
    \[
        \Age(\bff_0)=\k_0,
    \]
    we may realize $\bfb_0$ as a finite substructure of $\bff_0$. Because $\prec\in X_\k$, Proposition~\ref{prop:admissible-orderings-characterization} implies that
    \[
        \langle \bfb_0,\prec|B_0\rangle\in\k.
    \]
    Hence
    \[
        \bfa\leq \langle \bfb_0,\prec|B_0\rangle.
    \]
    By Lemma~\ref{lem:orbit-closure-characterization},
    \[
        \prec_0\in\overline{G_0\cdot \prec}.
    \]

    Now let
    \[
        \prec\in X_\k.
    \]
    The orbit closure $\overline{G_0\cdot \prec}$ is a closed $G_0$-invariant subset of $X_\k$ containing $\prec_0$. Since
    \[
        X_\k=\overline{G_0\cdot \prec_0},
    \]
    it follows that
    \[
        X_\k\subseteq \overline{G_0\cdot \prec}.
    \]
    The reverse inclusion is automatic, so
    \[
        \overline{G_0\cdot \prec}=X_\k.
    \]
    Thus every orbit is dense, and $X_\k$ is minimal.

    Conversely, assume that $X_\k$ is minimal. Let
    \[
        \bfa_0\in\k_0.
    \]
    We must find
    \[
        \bfb_0\in\k_0
    \]
    witnessing the ordering property for $\bfa_0$.

    There are only finitely many linear orderings on the finite set $A_0$, hence only finitely many ordered structures of the form
    \[
        \langle \bfa_0,\prec\rangle
    \]
    that belong to $\k$. List them as
    \[
        \bfa^1,\dots,\bfa^m.
    \]

    Fix
    \[
        \prec\in X_\k.
    \]
    Since $X_\k$ is minimal,
    \[
        \prec_0\in\overline{G_0\cdot \prec}.
    \]
    Applying Lemma~\ref{lem:orbit-closure-characterization} to each $\bfa^j$, we obtain finite substructures
    \[
        \bfc^{\,j}_{\prec,0}\leq \bff_0
        \qquad (j=1,\dots,m)
    \]
    such that
    \[
        \bfc^{\,j}_{\prec}
        =
        \langle \bfc^{\,j}_{\prec,0},\prec|C^{\,j}_{\prec,0}\rangle\in\k
    \]
    and
    \[
        \bfa^j\leq \bfc^{\,j}_{\prec}.
    \]
    Let
    \[
        C_{\prec,0}=C^{\,1}_{\prec,0}\cup\cdots\cup C^{\,m}_{\prec,0},
    \]
    viewed as the finite substructure of $\bff_0$ generated by this union. Since $\prec\in X_\k$, Proposition~\ref{prop:admissible-orderings-characterization} gives
    \[
        \langle C_{\prec,0},\prec|C_{\prec,0}\rangle\in\k,
    \]
    and by construction every $\bfa^j$ embeds into this ordered structure.

    Now define
    \[
        U_{\prec}
        =
        \{\prec'\in X_\k : \prec'|C_{\prec,0}=\prec|C_{\prec,0}\}.
    \]
    This is an open neighbourhood of $\prec$ in $X_\k$. As $\prec$ ranges over $X_\k$, the sets $U_{\prec}$ form an open cover of the compact space $X_\k$. Hence there exist
    \[
        \prec_1,\dots,\prec_n\in X_\k
    \]
    such that
    \[
        X_\k=U_{\prec_1}\cup\cdots\cup U_{\prec_n}.
    \]

    Choose a finite substructure
    \[
        \bfb_0\leq \bff_0
    \]
    containing each of the finite substructures
    \[
        C_{\prec_1,0},\dots,C_{\prec_n,0}.
    \]
    We claim that $\bfb_0$ witnesses the ordering property for $\bfa_0$.

    Let
    \[
        \bfa=\langle \bfa_0,\prec_A\rangle\in\k
        \qquad\text{and}\qquad
        \bfb=\langle \bfb_0,\prec_B\rangle\in\k.
    \]
    Since $\bfb\in\Age(\bff)=\k$, there exists a finite ordered substructure
    \[
        \widetilde{\bfb}\leq \bff
    \]
    isomorphic to $\bfb$. Let
    \[
        h\colon \widetilde{\bfb}\to \bfb
    \]
    be an isomorphism, and let
    \[
        g\in G_0
    \]
    be an automorphism of $\bff_0$ extending
    \[
        h^{-1}\colon \bfb_0\to \widetilde{B}_0.
    \]
    Then
    \[
        \widehat{\prec}=g\cdot \prec_0
    \]
    belongs to $X_\k$, and its restriction to $B_0$ is exactly $\prec_B$.

    Since
    \[
        X_\k=U_{\prec_1}\cup\cdots\cup U_{\prec_n},
    \]
    there exists $i\in\{1,\dots,n\}$ such that
    \[
        \widehat{\prec}\in U_{\prec_i}.
    \]
    Hence
    \[
        \widehat{\prec}|C_{\prec_i,0}=\prec_i|C_{\prec_i,0}.
    \]
    By construction of $C_{\prec_i,0}$, every ordered expansion of $\bfa_0$ belonging to $\k$ embeds into
    \[
        \langle C_{\prec_i,0},\prec_i|C_{\prec_i,0}\rangle.
    \]
    In particular,
    \[
        \bfa\leq \langle C_{\prec_i,0},\prec_i|C_{\prec_i,0}\rangle
        =
        \langle C_{\prec_i,0},\widehat{\prec}|C_{\prec_i,0}\rangle.
    \]
    Since
    \[
        C_{\prec_i,0}\subseteq B_0
    \]
    and $\widehat{\prec}|B_0=\prec_B$, the latter is a substructure of $\bfb$. Therefore
    \[
        \bfa\leq \bfb.
    \]
    This proves that $\k$ has the ordering property.
\end{proof}

The theorem shows that the minimality of the flow $X_\k$ is controlled by a purely finite combinatorial condition on the class $\k$. In the next section we add the Ramsey property and obtain the universal minimal flow.

\section{Ramsey property and universal minimal flows}

We now add the Ramsey property and obtain the main dynamical consequence of the Kechris--Pestov--Todorcevic correspondence.

Let
\[
    G={\rm Aut}(\bff).
\]
Since $\bff$ is an ordered expansion of $\bff_0$, we have
\[
    G\leq G_0,
\]
and every element of $G$ fixes the distinguished ordering $\prec_0$.

The first result shows that, under the Ramsey property, the pointed flow
\[
    (X_\k,\prec_0)
\]
is universal among $G_0$-ambits whose distinguished point is fixed by $G$.

\begin{lemma}
    \label{lem:graph-uniqueness}
    Assume that $\k$ has the Ramsey property. Let $(X,x_0)$ be a $G_0$-ambit such that
    \[
        g\cdot x_0=x_0
        \qquad\text{for all } g\in G.
    \]
    Let
    \[
        \Phi
        =
        \overline{\{(g\cdot \prec_0,g\cdot x_0):g\in G_0\}}
        \subseteq X_\k\times X.
    \]
    If
    \[
        (\prec,x),(\prec,x')\in \Phi,
    \]
    then
    \[
        x=x'.
    \]
\end{lemma}

\begin{proof}
    Suppose that
    \[
        (\prec,x),(\prec,x')\in \Phi
        \qquad\text{and}\qquad
        x\neq x'.
    \]
    Choose disjoint open neighbourhoods
    \[
        U\ni x,
        \qquad
        U'\ni x'.
    \]
    Since $X$ is compact Hausdorff, there exists an open neighbourhood $D$ of the diagonal in $X\times X$ such that
    \[
        D[U]\cap U'=\emptyset,
    \]
    where
    \[
        D[U]=\{y\in X:\text{ there exists } z\in U \text{ with } (z,y)\in D\}.
    \]
    By replacing $D$ with
    \[
        D\cap D^{-1},
    \]
    we may assume that $D$ is symmetric.

    Because the action of $G_0$ on the compact space $X$ is continuous, there exists an open neighbourhood $V$ of the identity in $G_0$ such that
    \[
        (v\cdot z,z)\in D
        \qquad\text{for all } z\in X \text{ and all } v\in V.
    \]
    Since $G_0\leq S_\infty$ is non-Archimedean, we may choose a finite substructure
    \[
        \bfa_0\leq \bff_0
    \]
    such that the pointwise stabilizer
    \[
        G_{0,(\bfa_0)}\subseteq V.
    \]

    Because
    \[
        (\prec,x),(\prec,x')\in \Phi,
    \]
    there exist
    \[
        g,h\in G_0
    \]
    such that
    \[
        g\cdot x_0\in U,
        \qquad
        h\cdot x_0\in U',
    \]
    and
    \[
        (g\cdot \prec_0)|A_0=(h\cdot \prec_0)|A_0=\prec|A_0.
    \]
    Consider the map
    \[
        \theta=h^{-1}g\colon g^{-1}(A_0)\to h^{-1}(A_0).
    \]
    Since
    \[
        (g\cdot \prec_0)|A_0=(h\cdot \prec_0)|A_0,
    \]
    the map $\theta$ is an isomorphism of the finite ordered structures
    \[
        \langle g^{-1}(\bfa_0),\prec_0\rangle
        \quad\text{and}\quad
        \langle h^{-1}(\bfa_0),\prec_0\rangle.
    \]
    Because
    \[
        \Age(\bff)=\k
    \]
    and $\bff$ is ultrahomogeneous, $\theta$ extends to some
    \[
        \gamma\in G={\rm Aut}(\bff).
    \]
    Then for every $a\in A_0$,
    \[
        h\gamma g^{-1}(a)=a,
    \]
    so
    \[
        v:=h\gamma g^{-1}\in G_{0,(\bfa_0)}\subseteq V.
    \]
    Since $\gamma\in G$ fixes $x_0$, we obtain
    \[
        h\cdot x_0=v\cdot g\cdot \gamma^{-1}\cdot x_0=v\cdot g\cdot x_0.
    \]
    Hence
    \[
        (h\cdot x_0,g\cdot x_0)\in D.
    \]
    Because $D$ is symmetric, this implies
    \[
        h\cdot x_0\in D[U].
    \]
    But
    \[
        g\cdot x_0\in U
        \qquad\text{and}\qquad
        h\cdot x_0\in U',
    \]
    so
    \[
        h\cdot x_0\in D[U]\cap U',
    \]
    contrary to the choice of $D$. Therefore
    \[
        x=x'.
    \]
\end{proof}

\begin{theorem}
    \label{thm:ramsey-universal-ambit}
    Assume that $\k$ has the Ramsey property. Let $(X,x_0)$ be a $G_0$-ambit such that
    \[
        g\cdot x_0=x_0
        \qquad\text{for all } g\in G.
    \]
    Then there exists a unique $G_0$-map
    \[
        \pi\colon X_\k\to X
    \]
    such that
    \[
        \pi(\prec_0)=x_0.
    \]
    In particular, $(X_\k,\prec_0)$ is universal among $G_0$-ambits whose distinguished point is fixed by $G$.
\end{theorem}

\begin{proof}
    Let
    \[
        \Phi
        =
        \overline{\{(g\cdot \prec_0,g\cdot x_0):g\in G_0\}}
        \subseteq X_\k\times X.
    \]
    Since the projection
    \[
        p_1\colon X_\k\times X\to X_\k
    \]
    is a closed map and
    \[
        p_1(\Phi)
    \]
    contains the dense set
    \[
        G_0\cdot \prec_0,
    \]
    we have
    \[
        p_1(\Phi)=X_\k.
    \]
    Thus for every
    \[
        \prec\in X_\k
    \]
    there exists
    \[
        x\in X
    \]
    such that
    \[
        (\prec,x)\in \Phi.
    \]
    By Lemma~\ref{lem:graph-uniqueness}, this $x$ is unique. Hence $\Phi$ is the graph of a function
    \[
        \pi\colon X_\k\to X.
    \]

    The map
    \[
        p_1|\Phi\colon \Phi\to X_\k
    \]
    is a continuous bijection from a compact space to a Hausdorff space, hence a homeomorphism. Therefore
    \[
        \pi=p_2\circ (p_1|\Phi)^{-1}
    \]
    is continuous.

    The set
    \[
        \{(g\cdot \prec_0,g\cdot x_0):g\in G_0\}
    \]
    is $G_0$-invariant, and so is its closure $\Phi$. It follows that $\pi$ is $G_0$-equivariant. Also,
    \[
        (\prec_0,x_0)\in \Phi,
    \]
    so
    \[
        \pi(\prec_0)=x_0.
    \]

    To prove uniqueness, let
    \[
        \pi'\colon X_\k\to X
    \]
    be another $G_0$-map with
    \[
        \pi'(\prec_0)=x_0.
    \]
    For every $g\in G_0$ we then have
    \[
        \pi'(g\cdot \prec_0)=g\cdot \pi'(\prec_0)=g\cdot x_0=\pi(g\cdot \prec_0).
    \]
    Since
    \[
        G_0\cdot \prec_0
    \]
    is dense in $X_\k$ and both maps are continuous, it follows that
    \[
        \pi'=\pi.
    \]

    Finally, because $(X,x_0)$ is an ambit, the orbit
    \[
        G_0\cdot x_0
    \]
    is dense in $X$. Since $\pi$ is continuous and $G_0$-equivariant, its image is a closed $G_0$-invariant subset of $X$ containing
    \[
        \pi(G_0\cdot \prec_0)=G_0\cdot x_0.
    \]
    Hence
    \[
        \pi(X_\k)=X.
    \]
    Thus $(X_\k,\prec_0)$ is universal among such ambits.
\end{proof}

We can now combine this with the results of Chapters~5 and~7 and the minimality criterion from the previous section.

\begin{corollary}
    \label{cor:ramsey-ordering-umf}
    Assume that $\k$ has both the Ramsey property and the ordering property. Then $X_\k$ is the universal minimal flow of $G_0={\rm Aut}(\bff_0)$.
\end{corollary}

\begin{proof}
    By Corollary~\ref{cor:fraisse-order-class-ea}, the Ramsey property of the Fra\"iss\'e order class $\k$ implies that
    \[
        G={\rm Aut}(\bff)
    \]
    is extremely amenable. Let $Y$ be a minimal compact $G_0$-flow. Restricting the action to the subgroup $G$, we obtain a compact $G$-flow, so extreme amenability yields a point
    \[
        y_0\in Y
    \]
    fixed by every element of $G$. Since $Y$ is minimal as a $G_0$-flow, the orbit
    \[
        G_0\cdot y_0
    \]
    is dense in $Y$, and therefore $(Y,y_0)$ is a $G_0$-ambit with $G$-fixed distinguished point.

    By Theorem~\ref{thm:ramsey-universal-ambit}, there exists a surjective $G_0$-map
    \[
        \pi\colon X_\k\to Y.
    \]
    On the other hand, by Theorem~\ref{thm:minimality-ordering-property}, the ordering property implies that $X_\k$ is minimal. Hence every minimal compact $G_0$-flow is a factor of $X_\k$, which is therefore the universal minimal flow of $G_0$.
\end{proof}

The preceding theorem shows how the Ramsey property and admissible orderings determine the universal minimal flow of the automorphism group of the reduct. In the following chapters we turn to the converse direction, uniqueness questions for such expansions, and the role of Ramsey degrees in the computation of universal minimal flows.
% !TEX root = main.tex
\chapter{Ramsey Degrees and Converse Results}
\label{chap:ramsey-degrees}

In the previous chapters we studied classes with the Ramsey property and saw how ordered Ramsey expansions determine the universal minimal flow of the automorphism group of the Fraïssé limit. However, many natural Fraïssé classes do not themselves satisfy the Ramsey property. The purpose of this chapter is to explain that, even in that situation, there is still a finite combinatorial invariant measuring the extent to which the Ramsey property fails.

This invariant is the \emph{Ramsey degree} of a finite structure. Roughly speaking, if \(\bfa_0\in\cK_0\), then the Ramsey degree of \(\bfa_0\) is the least number \(t\) such that, after passing to a sufficiently large copy of any ambient structure, every colouring of copies of \(\bfa_0\) can be reduced on some substructure to at most \(t\) colours. When \(\cK_0\) admits a reasonable ordered expansion with the Ramsey and ordering properties, these degrees can be computed explicitly in terms of the admissible orderings on \(\bfa_0\).

The chapter has two aims. First, we define these degrees and show how they arise from ordered Ramsey expansions. Second, we prove converse results showing that the combinatorial bounds given by Ramsey degrees recover both the Ramsey property and, ultimately, the universal minimal flow description from Chapter~\ref{chap:examples-umf}.

\section{Patterns and Ramsey degrees}

Throughout this chapter, \(L_0\) is a signature, \(\cK_0\) is a Fraïssé class in \(L_0\), and \(\cK\) is a hereditary order class in the expanded signature \(L=L_0\cup\{<\}\) such that \(\cK\!\restriction\!L_0=\cK_0\).
Thus \(\cK\) is an order expansion of \(\cK_0\).

Let \(\bfa_0\in\cK_0\). We write
\[
    \Adm_{\cK}(\bfa_0)
    =
    \{\prec \mid \prec \text{ is a linear ordering on }A_0 \text{ and } \langle \bfa_0,\prec\rangle\in\cK\}.
\]
So \(\Adm_{\cK}(\bfa_0)\) is the set of all \(\cK\)-admissible linear orderings on the underlying set of \(\bfa_0\).

The group \(\Aut(\bfa_0)\) acts on \(\Adm_{\cK}(\bfa_0)\) in the natural way: if \(g\in\Aut(\bfa_0)\) and \(\prec\in\Adm_{\cK}(\bfa_0)\), then \(g\cdot\prec\) is the linear order on \(A_0\) defined by
\[
    x\,(g\cdot\prec)\,y
    \quad\Longleftrightarrow\quad
    g^{-1}(x)\prec g^{-1}(y).
\]
Since \(\cK\) is closed under isomorphism, we have \(g\cdot\prec\in\Adm_{\cK}(\bfa_0)\) whenever \(\prec\in\Adm_{\cK}(\bfa_0)\).

\begin{definition}
    Let \(\prec\in\Adm_{\cK}(\bfa_0)\). The \emph{\(\bfa_0\)-pattern} of \(\prec\) (with respect to \(\cK\)) is its orbit under the action of \(\Aut(\bfa_0)\) on \(\Adm_{\cK}(\bfa_0)\).
\end{definition}

Equivalently, two orderings \(\prec,\prec'\in\Adm_{\cK}(\bfa_0)\) have the same pattern if and only if \(\langle \bfa_0,\prec\rangle \cong \langle \bfa_0,\prec'\rangle\).
Thus patterns record the isomorphism types of admissible ordered expansions of \(\bfa_0\).
For a fixed \(\bfa_0\), we write
\(\pat_{\cK,\bfa_0}(\prec)\) for the pattern of \(\prec\).

We write
\[
    t_\cK(\bfa_0)
    =
    \left|\Adm_{\cK}(\bfa_0)/\Aut(\bfa_0)\right|
\]
for the number of \(\bfa_0\)-patterns. Since the action of \(\Aut(\bfa_0)\) on \(\Adm_{\cK}(\bfa_0)\) is free, each orbit has cardinality \(|\Aut(\bfa_0)|\). Hence
\[
    t_\cK(\bfa_0)
    =
    \frac{\left|\Adm_{\cK}(\bfa_0)\right|}{\left|\Aut(\bfa_0)\right|}.
\]

To state the relevant partition relation, let
\(\bfa_0\hookrightarrow\bfb_0\hookrightarrow\bfc_0\) in \(\cK_0\), and let
\(k,t\geq 1\). We write
\[
    \bfc_0\to(\bfb_0)^{\bfa_0}_{k,t}
\]
to mean that for every colouring
\[
    c\colon\begin{pmatrix}\bfc_0\\ \bfa_0\end{pmatrix}\to\{1,\dots,k\},
\]
there is some
\[
    \bfb_0'\in\begin{pmatrix}\bfc_0\\ \bfb_0\end{pmatrix}
\]
such that the restriction of \(c\) to
\[
    \begin{pmatrix}\bfb_0'\\ \bfa_0\end{pmatrix}
\]
takes at most \(t\) values. In particular,
\[
    \bfc_0\to(\bfb_0)^{\bfa_0}_{k,1}
    \quad\Longleftrightarrow\quad
    \bfc_0\to(\bfb_0)^{\bfa_0}_{k}.
\]

For later use, suppose \(\bfa_0\hookrightarrow\bfb_0\) in \(\cK_0\) and fix an
ordered expansion \(\langle \bfb_0,\prec^{\bfb_0}\rangle\in\cK\). We write
\[
    t_\cK(\bfa_0,\bfb_0,\prec^{\bfb_0})
\]
for the number of \(\bfa_0\)-patterns represented by orderings
\(\prec\in\Adm_{\cK}(\bfa_0)\) such that
\(\langle \bfa_0,\prec\rangle\hookrightarrow
\langle \bfb_0,\prec^{\bfb_0}\rangle\).
Thus this number counts the patterns of copies of \(\bfa_0\) that actually occur
inside the fixed ordered expansion, and it is at most \(t_\cK(\bfa_0)\).

We can now define the main combinatorial invariant of the chapter.

\begin{definition}
    Let \(\bfa_0\in\cK_0\). The \emph{Ramsey degree} of \(\bfa_0\) in \(\cK_0\),
    denoted \(t(\bfa_0,\cK_0)\), is the least integer \(t\), if it exists, such
    that for every \(\bfb_0\in\cK_0\) with
    \(\bfa_0\hookrightarrow\bfb_0\) and every integer \(k\geq 2\), there is
    \(\bfc_0\in\cK_0\) with \(\bfb_0\hookrightarrow\bfc_0\) and
    \[
        \bfc_0\to(\bfb_0)^{\bfa_0}_{k,t}.
    \]
    If no such integer exists, we put \(t(\bfa_0,\cK_0)=\infty\).
\end{definition}

Thus \(t(\bfa_0,\cK_0)\) measures how far \(\bfa_0\) is from satisfying the usual Ramsey property inside the class \(\cK_0\). The case \(t(\bfa_0,\cK_0)=1\) is exactly the ordinary Ramsey statement for copies of \(\bfa_0\).

The guiding principle of this chapter is that, when \(\cK_0\) admits a suitable ordered Ramsey expansion \(\cK\), the Ramsey degree of \(\bfa_0\) is determined by the finite set of admissible patterns on \(\bfa_0\). In the next section we make this precise.

A basic example is given by the class \(\mathcal{GR}\) of finite graphs. This class does not have the Ramsey property, but its standard ordered expansion by linear orders does. If \(\bfa_0\in\mathcal{GR}\), then every linear ordering of \(A_0\) is admissible, and two such orderings have the same pattern exactly when they differ by an automorphism of \(\bfa_0\). Thus
\[
    t(\bfa_0,\mathcal{GR})
    =
    \frac{|A_0|!}{|\Aut(\bfa_0)|}.
\]
This illustrates the general picture: Ramsey degrees arise from counting the admissible order-types of a structure up to automorphism.

\section{Ramsey degrees from Ramsey expansions}

We now show that if the ordered class \(\cK\) has the Ramsey property, then the number of admissible patterns gives an upper bound on the number of colours that must remain on a large substructure. If \(\cK\) also has the ordering property, then this bound is optimal. In particular, the Ramsey degree of \(\bfa_0\) is exactly the number of \(\cK\)-patterns on \(\bfa_0\).

\begin{proposition}
    \label{prop:ramsey-expansion-upper-bound}
    Assume that \(\cK\) has the Ramsey property. Let
    \(\bfa_0\hookrightarrow\bfb_0\) in \(\cK_0\), let
    \(\langle \bfb_0,\prec^{\bfb_0}\rangle\in\cK\), and let \(k\geq 2\). Then there
    is \(\bfc_0\) in \(\cK_0\) with \(\bfb_0\hookrightarrow\bfc_0\) such that
    \[
        \bfc_0\to(\bfb_0)^{\bfa_0}_{k,\,t_\cK(\bfa_0,\bfb_0,\prec^{\bfb_0})}.
    \]
    In particular, for every hereditary class \(\cK'\subseteq\cK\) with
    \(\cK'\!\restriction\!L_0=\cK_0\), there is \(\bfc_0\) in \(\cK_0\) with
    \(\bfb_0\hookrightarrow\bfc_0\) such that
    \[
        \bfc_0\to(\bfb_0)^{\bfa_0}_{k,\,t_{\cK'}(\bfa_0)}.
    \]
\end{proposition}

\begin{proof}
    Put \(t=t_\cK(\bfa_0,\bfb_0,\prec^{\bfb_0})\). Choose representatives \(\prec_1,\dots,\prec_t\) for the \(\bfa_0\)-patterns realised inside \(\langle\bfb_0,\prec^{\bfb_0}\rangle\).

    We define inductively ordered structures
    \[
        \langle\bfd_0^{(i)},\prec^{\bfd_0^{(i)}}\rangle\in\cK
        \qquad (0\leq i\leq t)
    \]
    as follows. Let
    \(\langle\bfd_0^{(0)},\prec^{\bfd_0^{(0)}}\rangle
      =\langle\bfb_0,\prec^{\bfb_0}\rangle\), and, for each
    \(1\leq i\leq t\), choose
    \(\langle\bfd_0^{(i)},\prec^{\bfd_0^{(i)}}\rangle\in\cK\)
    such that
    \[
        \langle\bfd_0^{(i)},\prec^{\bfd_0^{(i)}}\rangle\to
        \bigl(\langle\bfd_0^{(i-1)},\prec^{\bfd_0^{(i-1)}}\rangle\bigr)^{\langle\bfa_0,\prec_i\rangle}_{k}.
    \]
    Set \(\bfc_0=\bfd_0^{(t)}\) and write
    \(\prec^{\bfc_0}=\prec^{\bfd_0^{(t)}}\).

    We claim that \(\bfc_0\) has the required property. Let
    \[
        c\colon\begin{pmatrix}\bfc_0\\ \bfa_0\end{pmatrix}\to\{1,\dots,k\}
    \]
    be any colouring of copies of \(\bfa_0\) in \(\bfc_0\).

    We now work downward from
    \(\langle\bfd_0^{(t)},\prec^{\bfd_0^{(t)}}\rangle\). For the pattern
    represented by \(\prec_t\), define a colouring
    \[
        c_t\colon
        \begin{pmatrix}
            \langle\bfd_0^{(t)},\prec^{\bfd_0^{(t)}}\rangle \\
            \langle\bfa_0,\prec_t\rangle
        \end{pmatrix}
        \to
        \{1,\dots,k\}
    \]
    by sending a copy
    \[
        \langle\bfa_0',\prec^{\bfd_0^{(t)}}\restriction A_0'\rangle\cong\langle\bfa_0,\prec_t\rangle
    \]
    to the colour \(c(\bfa_0')\) of its reduct. By the choice of
    \(\langle\bfd_0^{(t)},\prec^{\bfd_0^{(t)}}\rangle\), there is a copy
    \[
        \langle\widetilde{\bfd}_0^{(t-1)},
        \prec^{\widetilde{\bfd}_0^{(t-1)}}\rangle\cong
        \langle\bfd_0^{(t-1)},\prec^{\bfd_0^{(t-1)}}\rangle
    \]
    inside \(\langle\bfd_0^{(t)},\prec^{\bfd_0^{(t)}}\rangle\) on which every
    copy of \(\langle\bfa_0,\prec_t\rangle\) has the same \(c_t\)-colour,
    say \(k_t\).

    Next define
    \[
        c_{t-1}\colon
        \begin{pmatrix}
            \langle\widetilde{\bfd}_0^{(t-1)},
            \prec^{\widetilde{\bfd}_0^{(t-1)}}\rangle \\
            \langle\bfa_0,\prec_{t-1}\rangle
        \end{pmatrix}
        \to
        \{1,\dots,k\}
    \]
    in the same way, by pulling back the colouring \(c\) from reducts. Again by the Ramsey property, there is a copy
    \[
        \langle\widetilde{\bfd}_0^{(t-2)},
        \prec^{\widetilde{\bfd}_0^{(t-2)}}\rangle\cong
        \langle\bfd_0^{(t-2)},\prec^{\bfd_0^{(t-2)}}\rangle
    \]
    inside \(\langle\widetilde{\bfd}_0^{(t-1)},
    \prec^{\widetilde{\bfd}_0^{(t-1)}}\rangle\) on which every copy of
    \(\langle\bfa_0,\prec_{t-1}\rangle\) has the same colour, say \(k_{t-1}\).

    Continuing in this way, after \(t\) steps we obtain a copy
    \[
        \langle\bfb_0',\prec^{\bfb_0'}\rangle\cong
        \langle\bfb_0,\prec^{\bfb_0}\rangle
    \]
    inside \(\langle\bfc_0,\prec^{\bfc_0}\rangle\) together with colours
    \(k_1,\dots,k_t\in\{1,\dots,k\}\) such that, for each
    \(i\in\{1,\dots,t\}\), every copy of \(\langle\bfa_0,\prec_i\rangle\)
    inside \(\langle\bfb_0',\prec^{\bfb_0'}\rangle\) has colour \(k_i\).

    Now let
    \[
        \bfa_0'\in\begin{pmatrix}\bfb_0'\\ \bfa_0\end{pmatrix}.
    \]
    The induced ordering \(\prec^{\bfb_0'}\restriction A_0'\) realises one of the patterns represented by some \(\prec_i\). Hence \(c(\bfa_0')=k_i\). Therefore the colouring \(c\) takes at most the \(t\) values \(k_1,\dots,k_t\)
    on
    \[
        \begin{pmatrix}\bfb_0'\\ \bfa_0\end{pmatrix}.
    \]
    This proves
    \[
        \bfc_0\to(\bfb_0)^{\bfa_0}_{k,t}.
    \]

    For the final claim, note that \(t_\cK(\bfa_0,\bfb_0,\prec^{\bfb_0})\leq t_{\cK'}(\bfa_0)\),
    since every pattern realised in \(\langle\bfb_0,\prec^{\bfb_0}\rangle\) with respect to \(\cK'\) is also a \(\cK'\)-pattern on \(\bfa_0\). Hence the same \(\bfc_0\) also satisfies
    \[
        \bfc_0\to(\bfb_0)^{\bfa_0}_{k,\,t_{\cK'}(\bfa_0)}.
    \]
\end{proof}

The proof yields a slightly stronger conclusion: one may choose a linear ordering on \(\bfc_0\) such that, on a suitable copy of \(\bfb_0\), the colour of a copy of \(\bfa_0\) depends only on the pattern induced by that ordering. We will not isolate this as a separate proposition, since the weaker statement above is all that we need.

\begin{proposition}
    \label{prop:ordering-property-lower-bound}
    Assume that \(\cK\) has both the Ramsey property and the ordering property.
    Then for every \(\bfa_0\in\cK_0\) there exists \(\bfb_0\in\cK_0\) with
    \(\bfa_0\hookrightarrow\bfb_0\) such that for every \(\bfc_0\in\cK_0\)
    with \(\bfb_0\hookrightarrow\bfc_0\) there is a colouring
    \[
        c\colon\begin{pmatrix}\bfc_0\\ \bfa_0\end{pmatrix}\to\{1,\dots,t_\cK(\bfa_0)\}
    \]
    with the following property: for every
    \[
        \bfb_0'\in\begin{pmatrix}\bfc_0\\ \bfb_0\end{pmatrix},
    \]
    the colouring \(c\) takes all the values \(1,\dots,t_\cK(\bfa_0)\) on
    \[
        \begin{pmatrix}\bfb_0'\\ \bfa_0\end{pmatrix}.
    \]
\end{proposition}

\begin{proof}
    By the ordering property, there is \(\bfb_0\in\cK_0\) with
    \(\bfa_0\hookrightarrow\bfb_0\) such that whenever
    \[
        \langle\bfa_0,\prec\rangle,\langle\bfb_0,\prec'\rangle\in\cK,
    \]
    we have
    \[
        \langle\bfa_0,\prec\rangle\hookrightarrow\langle\bfb_0,\prec'\rangle.
    \]
    Equivalently, every admissible \(\bfa_0\)-pattern is realised by some copy of \(\bfa_0\) inside every ordered expansion of \(\bfb_0\) belonging to \(\cK\).

    Now let \(\bfc_0\in\cK_0\) with \(\bfb_0\hookrightarrow\bfc_0\), and choose a linear ordering \(\prec^{\bfc_0}\) on \(C_0\) such that \(\langle\bfc_0,\prec^{\bfc_0}\rangle\in\cK\). Enumerate the \(\bfa_0\)-patterns as \(p_1,\dots,p_{t_\cK(\bfa_0)}\).
    Define
    \[
        c\colon\begin{pmatrix}\bfc_0\\ \bfa_0\end{pmatrix}\to\{1,\dots,t_\cK(\bfa_0)\}
    \]
    by declaring that, for
    \[
        \bfa_0'\in\begin{pmatrix}\bfc_0\\ \bfa_0\end{pmatrix},
    \]
    we put \(c(\bfa_0')=i\) if the ordered structure \(\langle\bfa_0',\prec^{\bfc_0}\restriction A_0'\rangle\) has pattern \(p_i\).

    Let
    \[
        \bfb_0'\in\begin{pmatrix}\bfc_0\\ \bfb_0\end{pmatrix}.
    \]
    Then \(\langle\bfb_0',\prec^{\bfc_0}\restriction B_0'\rangle\in\cK\), and by the choice of \(\bfb_0'\), every admissible \(\bfa_0\)-pattern is realised inside this ordered structure. Therefore the colouring \(c\) takes every value \(1,\dots,t_\cK(\bfa_0)\)
    on
    \[
        \begin{pmatrix}\bfb_0'\\ \bfa_0\end{pmatrix}.
    \]
\end{proof}

\begin{corollary}
    \label{cor:ramsey-degree-pattern-count}
    Assume that \(\cK\) has the Ramsey property and the ordering property. Then for every \(\bfa_0\in\cK_0\),
    \[
        t(\bfa_0,\cK_0)=t_\cK(\bfa_0).
    \]
    Equivalently, \(t_\cK(\bfa_0)\) is the least integer \(t\) such that for every \(\bfb_0\in\cK_0\) with \(\bfa_0\hookrightarrow\bfb_0\) and every \(k\geq 2\), there exists \(\bfc_0\in\cK_0\) with \(\bfb_0\hookrightarrow\bfc_0\) and
    \[
        \bfc_0\to(\bfb_0)^{\bfa_0}_{k,t}.
    \]
\end{corollary}

\begin{proof}
    Proposition~\ref{prop:ramsey-expansion-upper-bound} shows that \(t_\cK(\bfa_0)\) has the required partition property, so
    \[
        t(\bfa_0,\cK_0)\leq t_\cK(\bfa_0).
    \]
    On the other hand, Proposition~\ref{prop:ordering-property-lower-bound} shows that no smaller number can work. Hence
    \[
        t(\bfa_0,\cK_0)=t_\cK(\bfa_0).
    \]
\end{proof}

In particular, whenever \(\cK\) has the Ramsey and ordering properties, the number of admissible patterns on \(\bfa_0\) depends only on the reduct class \(\cK_0\), not on the particular expansion \(\cK\). Thus ordered Ramsey expansions do not merely imply the finiteness of Ramsey degrees; they compute those degrees exactly.

\section{Converse criteria}

We now prove converses to the results of the previous section. The first shows that, under the ordering property, the partition relation with bound \(t_\cK(\bfa_0)\) already forces the Ramsey property. We then isolate a local form of the ordering property and show that, assuming the Ramsey property, it is equivalent to the equality \(t(\bfa_0,\cK_0)=t_\cK(\bfa_0)\).

\begin{proposition}
    \label{prop:partition-bound-implies-ramsey}
    Suppose that \(\cK\) is reasonable and has the ordering property. Assume
    moreover that for every \(\bfa_0\hookrightarrow\bfb_0\) in \(\cK_0\) and
    every \(k\geq 2\), there exists \(\bfc_0\in\cK_0\) with
    \(\bfb_0\hookrightarrow\bfc_0\) such that
    \[
        \bfc_0\to(\bfb_0)^{\bfa_0}_{k,\,t_\cK(\bfa_0)}.
    \]
    Then \(\cK\) has the Ramsey property.

    In particular, any reasonable order class \(\cK\) with the ordering property that is contained in an expansion of \(\cK_0\) with the Ramsey property also has the Ramsey property.
\end{proposition}

\begin{proof}
    Let
    \[
        \langle\bfa_0,\prec\rangle,\langle\bfb_0,\prec'\rangle\in\cK
        \qquad\text{with}\qquad
        \langle\bfa_0,\prec\rangle\hookrightarrow\langle\bfb_0,\prec'\rangle.
    \]
    We must find \(\langle\bfc_0,\prec''\rangle\in\cK\) such that
    \[
        \langle\bfc_0,\prec''\rangle\to
        \bigl(\langle\bfb_0,\prec'\rangle\bigr)^{\langle\bfa_0,\prec\rangle}_{2}.
    \]

    By the ordering property, there exists \(\bfb_1\in\cK_0\) with
    \(\bfb_0\hookrightarrow\bfb_1\) such that whenever
    \[
        \langle\bfb_0,\prec^{\bfb_0}\rangle,\langle\bfb_1,\prec^{\bfb_1}\rangle\in\cK,
    \]
    we have
    \[
        \langle\bfb_0,\prec^{\bfb_0}\rangle\hookrightarrow\langle\bfb_1,\prec^{\bfb_1}\rangle.
    \]
    Since \(\cK\) is reasonable and \(\langle\bfa_0,\prec\rangle\hookrightarrow\langle\bfb_0,\prec'\rangle\), it follows that for every ordering \(\prec^{\bfb_1}\) on \(B_1\) with \(\langle\bfb_1,\prec^{\bfb_1}\rangle\in\cK\), we also have
    \[
        \langle\bfa_0,\prec\rangle\hookrightarrow\langle\bfb_1,\prec^{\bfb_1}\rangle.
    \]

    Now apply the hypothesis with \(\bfa_0\), \(\bfb_1\), and \(k=t_\cK(\bfa_0)+1\).
    Thus there exists \(\bfc_0\in\cK_0\) with
    \(\bfb_1\hookrightarrow\bfc_0\) such that
    \[
        \bfc_0\to(\bfb_1)^{\bfa_0}_{\,t_\cK(\bfa_0)+1,\;t_\cK(\bfa_0)}.
    \]
    Choose an ordering \(\prec''\) on \(C_0\) such that \(\langle\bfc_0,\prec''\rangle\in\cK\).
    We claim that this ordered structure has the required Ramsey property.

    Let
    \[
        c\colon
        \begin{pmatrix}
            \langle\bfc_0,\prec''\rangle \\
            \langle\bfa_0,\prec\rangle
        \end{pmatrix}
        \to
        \{1,2\}
    \]
    be a colouring of copies of \(\langle\bfa_0,\prec\rangle\) in \(\langle\bfc_0,\prec''\rangle\). Enumerate the \(\bfa_0\)-patterns as \(p_1,\dots,p_{t_\cK(\bfa_0)}\).
    Using \(c\), define a colouring
    \[
        c'\colon
        \begin{pmatrix}\bfc_0\\ \bfa_0\end{pmatrix}
        \to
        \{0,1,\dots,t_\cK(\bfa_0)\}
    \]
    as follows. For
    \[
        \bfa_0'\in\begin{pmatrix}\bfc_0\\ \bfa_0\end{pmatrix},
    \]
    let
    \[
        q(\bfa_0')
        =
        \pat_{\cK,\bfa_0}(\prec''\restriction A_0')
    \]
    be the pattern of the induced ordering on \(A_0'\). Set
    \[
        c'(\bfa_0')=
        \begin{cases}
            0, & \text{if }q(\bfa_0')=\pat_{\cK,\bfa_0}(\prec)
            \text{ and }c(\langle\bfa_0',\prec''\restriction A_0'\rangle)=2,  \\[1ex]
            i, & \text{if }q(\bfa_0')=p_i\text{ and either }
            p_i\neq\pat_{\cK,\bfa_0}(\prec)
            \text{ or }c(\langle\bfa_0',\prec''\restriction A_0'\rangle)=1.
        \end{cases}
    \]
    This uses at most \(t_\cK(\bfa_0)+1\) colours.

    Hence there exists
    \[
        \bfb_1'\in\begin{pmatrix}\bfc_0\\ \bfb_1\end{pmatrix}
    \]
    such that \(c'\) takes at most \(t_\cK(\bfa_0)\) values on
    \[
        \begin{pmatrix}\bfb_1'\\ \bfa_0\end{pmatrix}.
    \]
    Since \(\langle\bfb_1',\prec''\restriction B_1'\rangle\in\cK\) and \(\langle\bfa_0,\prec\rangle\) embeds into every ordered expansion of \(\bfb_1\) in \(\cK\), the pattern of \(\prec\) is realised by some copy of \(\bfa_0\) inside \(\bfb_1'\). Moreover, every \(\bfa_0\)-pattern realised in \(\langle\bfb_1',\prec''\restriction B_1'\rangle\) contributes at least one value of \(c'\) on
    \[
        \begin{pmatrix}\bfb_1'\\ \bfa_0\end{pmatrix}.
    \]
    Since \(c'\) takes at most \(t_\cK(\bfa_0)\) values there, the value \(0\) must occur and, for the pattern of \(\prec\), no additional value can occur. It follows that all copies of \(\langle\bfa_0,\prec\rangle\) inside \(\langle\bfb_1',\prec''\restriction B_1'\rangle\) have the same \(c\)-colour, namely \(2\).

    Finally, by the choice of \(\bfb_1\) we have
    \[
        \langle\bfb_0,\prec'\rangle\hookrightarrow\langle\bfb_1',\prec''\restriction B_1'\rangle,
    \]
    so
    \[
        \langle\bfc_0,\prec''\rangle\to
        \bigl(\langle\bfb_0,\prec'\rangle\bigr)^{\langle\bfa_0,\prec\rangle}_{2}.
    \]
    Thus \(\cK\) has the Ramsey property.

    For the final assertion, if \(\cK\) is contained in an expansion of \(\cK_0\) with the Ramsey property, then Proposition~\ref{prop:ramsey-expansion-upper-bound} yields exactly the partition relation assumed above. Hence the first part applies.
\end{proof}

To formulate the converse to Corollary~\ref{cor:ramsey-degree-pattern-count}, it is useful to isolate a local form of the ordering property.

\begin{definition}
    Let \(\bfa_0\in\cK_0\). We say that \(\cK\) has the \emph{ordering property at} \(\bfa_0\) if there exists \(\bfb_0\in\cK_0\) with \(\bfa_0\hookrightarrow\bfb_0\) such that for every pair of orderings \(\prec\) on \(A_0\) and \(\prec'\) on \(B_0\) with
    \[
        \langle\bfa_0,\prec\rangle,\langle\bfb_0,\prec'\rangle\in\cK,
    \]
    we have
    \[
        \langle\bfa_0,\prec\rangle\hookrightarrow\langle\bfb_0,\prec'\rangle.
    \]
\end{definition}

Thus \(\cK\) has the ordering property if and only if it has the ordering property at every \(\bfa_0\in\cK_0\).

\begin{proposition}
    \label{prop:local-ordering-property-ramsey-degree}
    Assume that \(\cK\) has the Ramsey property. Then, for each \(\bfa_0\in\cK_0\), the following are equivalent:

    \begin{enumerate}
        \item[(i)] \(\cK\) has the ordering property at \(\bfa_0\);
        \item[(ii)] \(t(\bfa_0,\cK_0)=t_\cK(\bfa_0)\).
    \end{enumerate}
\end{proposition}

\begin{proof}
    Assume first that \(\cK\) has the ordering property at \(\bfa_0\). Then the proof of Proposition~\ref{prop:ordering-property-lower-bound} shows that there exists \(\bfb_0\) with \(\bfa_0\hookrightarrow\bfb_0\) such that every admissible \(\bfa_0\)-pattern is realised inside every ordered expansion of \(\bfb_0\) belonging to \(\cK\). Combining this with Proposition~\ref{prop:ramsey-expansion-upper-bound}, exactly as in the proof of Corollary~\ref{cor:ramsey-degree-pattern-count}, we obtain \(t(\bfa_0,\cK_0)=t_\cK(\bfa_0)\).

    Conversely, suppose that \(t(\bfa_0,\cK_0)=t_\cK(\bfa_0)\)
    but that \(\cK\) does not have the ordering property at \(\bfa_0\). Then for every \(\bfb_0\in\cK_0\) with \(\bfa_0\hookrightarrow\bfb_0\) there exists an ordering \(\prec'\) on \(B_0\) such that \(\langle\bfb_0,\prec'\rangle\in\cK\) and some admissible ordering \(\prec\) on \(A_0\) with \(\langle\bfa_0,\prec\rangle\in\cK\)
    fails to embed into \(\langle\bfb_0,\prec'\rangle\). Hence the number of \(\bfa_0\)-patterns realised inside \(\langle\bfb_0,\prec'\rangle\) satisfies
    \[
        t_\cK(\bfa_0,\bfb_0,\prec')\leq t_\cK(\bfa_0)-1.
    \]
    Applying Proposition~\ref{prop:ramsey-expansion-upper-bound} with this ordered structure, we obtain, for every \(k\geq 2\), some \(\bfc_0\) with \(\bfb_0\hookrightarrow\bfc_0\) such that
    \[
        \bfc_0\to(\bfb_0)^{\bfa_0}_{k,\,t_\cK(\bfa_0)-1}.
    \]
    This implies \(t(\bfa_0,\cK_0)\leq t_\cK(\bfa_0)-1\),
    contrary to the assumption. Therefore \(\cK\) has the ordering property at \(\bfa_0\).
\end{proof}

\begin{corollary}
    \label{cor:ordering-property-ramsey-degree}
    Assume that \(\cK\) has the Ramsey property. Then the following are equivalent:

    \begin{enumerate}
        \item[(i)] \(\cK\) has the ordering property;
        \item[(ii)] for every \(\bfa_0\in\cK_0\),
              \[
                  t(\bfa_0,\cK_0)=t_\cK(\bfa_0).
              \]
    \end{enumerate}
\end{corollary}

\begin{proof}
    This follows immediately from Proposition~\ref{prop:local-ordering-property-ramsey-degree}, since the ordering property is equivalent to having the ordering property at each \(\bfa_0\in\cK_0\).
\end{proof}


\section{Recovering Ramsey expansions}

We now return to the dynamical viewpoint of the previous chapters. The results of this section show that the combinatorial information encoded by Ramsey degrees is strong enough to recover ordered expansions with the Ramsey and ordering properties, and hence to recover the universal minimal flow description from Chapter~\ref{chap:examples-umf}.

\begin{theorem}
    \label{thm:recover-ramsey-ordering-subexpansion}
    Let \(\cK_0\) be a Fraïssé class in a signature \(L_0\), and let \(\cK\) be a reasonable Fraïssé order class in the expanded signature \(L=L_0\cup\{<\}\) such that \(\cK\!\restriction\!L_0=\cK_0\). Assume that \(\cK\) has the Ramsey property. Then there exists a reasonable Fraïssé order class \(\cK'\subseteq\cK\) such that \(\cK'\!\restriction\!L_0=\cK_0\) and \(\cK'\) has both the Ramsey property and the ordering property.
\end{theorem}

\begin{proof}
    Let \(\bff_0=\Flim(\cK_0)\) and \(G_0=\Aut(\bff_0)\). Consider the \(G_0\)-flow \(X_\cK\) of all \(\cK\)-admissible linear orderings on \(F_0\). Choose a minimal subflow \(X'\subseteq X_\cK\), and fix some \(\prec'\in X'\). Let \(\cK'=\Age(\langle \bff_0,\prec'\rangle)\).
    Then \(\cK'\subseteq\cK\), and since every finite \(L_0\)-substructure of \(\bff_0\) lies in \(\cK_0\), we have \(\cK'\!\restriction\!L_0=\cK_0\).
    Moreover, \(\cK'\) is hereditary and has the joint embedding property, and it is easy to verify that \(\cK'\) is reasonable.

    We claim that
    \[
        X'=X_{\cK'}.
    \]
    Indeed, every finite substructure of \(\langle\bff_0,\prec'\rangle\) belongs to \(\cK'\), so every ordering in the orbit closure \(\overline{G_0\cdot \prec'}=X'\) is \(\cK'\)-admissible. Conversely, if an ordering on \(F_0\) is \(\cK'\)-admissible, then every finite restriction occurs in \(\Age(\langle\bff_0,\prec'\rangle)\), hence belongs to the orbit closure of \(\prec'\). Thus the two spaces coincide.

    Since \(X'=X_{\cK'}\) is minimal, it follows from
    Theorem~\ref{thm:minimality-ordering-property} that \(\cK'\) has the ordering
    property.

    It remains to show that \(\cK'\) has the Ramsey property. By Proposition~\ref{prop:partition-bound-implies-ramsey}, it is enough to verify that for every \(\bfa_0\hookrightarrow\bfb_0\) in \(\cK_0\) and every \(k\geq 2\), there exists \(\bfc_0\in\cK_0\) with \(\bfb_0\hookrightarrow\bfc_0\) such that
    \[
        \bfc_0\to(\bfb_0)^{\bfa_0}_{k,\,t_{\cK'}(\bfa_0)}.
    \]
    But this follows from Proposition~\ref{prop:ramsey-expansion-upper-bound}, since \(\cK\) has the Ramsey property and \(\cK'\subseteq\cK\) is a hereditary expansion of \(\cK_0\).

    Finally, because \(\cK'\) is reasonable and has the Ramsey and ordering properties, Corollary~\ref{cor:ordering-property-ramsey-degree} implies that the numbers \(t_{\cK'}(\bfa_0)\) are the Ramsey degrees in \(\cK_0\). In particular, Proposition~\ref{prop:partition-bound-implies-ramsey} yields the Ramsey property for \(\cK'\), and hence \(\cK'\) is a reasonable Fraïssé order class with the required properties.
\end{proof}

The preceding theorem shows that, once one has a reasonable Ramsey expansion, one may pass to a smaller expansion in which minimality is built in. The next result is the converse direction of the universal-minimal-flow theorem proved earlier.

\begin{theorem}
    \label{thm:umf-iff-ramsey-ordering}
    Let \(L\supseteq\{<\}\) be a signature, let \(L_0=L\setminus\{<\}\), and let \(\cK\) be a reasonable Fraïssé order class in \(L\). Put
    \[
        \cK_0=\cK\!\restriction\!L_0,
        \qquad
        \bff=\Flim(\cK),
        \qquad
        \bff_0=\bff\!\restriction\!L_0=\Flim(\cK_0),
    \]
    and let \(G_0=\Aut(\bff_0)\) and \(G=\Aut(\bff)\).
    Let \(X_\cK\) be the compact space of all \(\cK\)-admissible linear orderings on the underlying set \(F_0\) of \(\bff_0\). Then the following are equivalent:

    \begin{enumerate}
        \item[(i)] \(\cK\) has the Ramsey property and the ordering property;
        \item[(ii)] \(X_\cK\) is the universal minimal flow of \(G_0\).
    \end{enumerate}
\end{theorem}

\begin{proof}
    The implication \((i)\Rightarrow(ii)\) is
    Corollary~\ref{cor:ramsey-ordering-umf}.

    Now assume that \(X_\cK\) is the universal minimal flow of \(G_0\). Since it is
    in particular minimal, Theorem~\ref{thm:minimality-ordering-property}
    implies that \(\cK\) has the ordering property. It therefore remains to prove
    that \(\cK\) has the Ramsey property.

    By Proposition~\ref{prop:partition-bound-implies-ramsey}, it is enough to show that for every \(\bfa_0\hookrightarrow\bfb_0\) in \(\cK_0\),
    every \(k\geq 2\), and every colouring
    \[
        c\colon\begin{pmatrix}\bff_0\\ \bfa_0\end{pmatrix}\to\{1,\dots,k\},
    \]
    there exists
    \[
        \bfb_0'\in\begin{pmatrix}\bff_0\\ \bfb_0\end{pmatrix}
    \]
    such that the restriction of \(c\) to
    \[
        \begin{pmatrix}\bfb_0'\\ \bfa_0\end{pmatrix}
    \]
    takes at most \(t_\cK(\bfa_0)\) values.

    Consider the compact \(G_0\)-flow
    \[
        \{1,\dots,k\}^{\begin{pmatrix}\bff_0\\ \bfa_0\end{pmatrix}},
    \]
    where \(G_0\) acts by
    \[
        (g\cdot \gamma)(\bfa_0')=\gamma(g^{-1}(\bfa_0'))
        \qquad
        \left(g\in G_0,\ \gamma\in \{1,\dots,k\}^{\begin{pmatrix}\bff_0\\ \bfa_0\end{pmatrix}}\right).
    \]
    Let \(X=\overline{G_0\cdot c}\) be the orbit closure of the colouring \(c\). Since \(X_\cK\) is universal minimal, there is a \(G_0\)-map \(\pi\colon X_\cK\to X\). Let \(\prec_0\) be the distinguished linear ordering on \(F_0\) coming from the Fraïssé limit \(\bff=\langle\bff_0,\prec_0\rangle\), and put \(c_0=\pi(\prec_0)\).

    Because every element of \(G=\Aut(\bff)\) fixes the ordering \(\prec_0\), and \(\pi\) is \(G_0\)-equivariant, every element of \(G\) also fixes \(c_0\). Hence if
    \[
        \bfa_0',\bfa_0''\in\begin{pmatrix}\bff_0\\ \bfa_0\end{pmatrix}
    \]
    have the same \(\cK\)-pattern with respect to the ordering \(\prec_0\), then some element of \(G\) sends \(\bfa_0'\) to \(\bfa_0''\), and therefore \(c_0(\bfa_0')=c_0(\bfa_0'')\). Thus the value of \(c_0\) on a copy of \(\bfa_0\) depends only on its pattern. It follows that \(c_0\) takes at most \(t_\cK(\bfa_0)\) values.

    Now fix any
    \[
        \bfb_0''\in\begin{pmatrix}\bff_0\\ \bfb_0\end{pmatrix}.
    \]
    Since \(c_0\in \overline{G_0\cdot c}\), there exists \(g\in G_0\) such that
    \[
        (g\cdot c)\restriction\begin{pmatrix}\bfb_0''\\ \bfa_0\end{pmatrix}
        =
        c_0\restriction\begin{pmatrix}\bfb_0''\\ \bfa_0\end{pmatrix}.
    \]
    Let \(\bfb_0'=g^{-1}(\bfb_0'')\).
    Then the restriction of \(c\) to
    \[
        \begin{pmatrix}\bfb_0'\\ \bfa_0\end{pmatrix}
    \]
    has the same set of values as the restriction of \(c_0\) to
    \[
        \begin{pmatrix}\bfb_0''\\ \bfa_0\end{pmatrix},
    \]
    and therefore takes at most \(t_\cK(\bfa_0)\) values. This is exactly the partition relation required in Proposition~\ref{prop:partition-bound-implies-ramsey}.

    Hence \(\cK\) has the Ramsey property. Since we already know that \(\cK\) has the ordering property, this proves \((ii)\Rightarrow(i)\).
\end{proof}

The preceding two theorems complete the converse direction of the Kechris--Pestov--Todorcevic (KPT) correspondence~\cite{KPT05} developed in the previous chapters. Ordered Ramsey expansions not only produce universal minimal flows, but can also be recovered from the dynamical structure of those flows. In this sense, Ramsey degrees measure the residual combinatorial complexity of the reduct class, while the Ramsey and ordering properties together characterise the expansions that realise the universal minimal flow.



\backmatter
% !TEX root = main.tex
\begin{thebibliography}{100}

\bibitem{AH78}
F.G. Abramson and L.A. Harrington, Models without indiscernibles,
{\it J. Symbolic Logic} {\bf 43} (1978), 572--600.

\bibitem{AS}
S. Adams and G. Stuck, Splitting of nonnegatively curved leaves in
minimal sets of foliations, {\it Duke Math. J.} (1993), {\bf 71}(1),
71--92.

\bibitem{A}
J. Auslander, {\it Minimal Flows and Their Extensions}, North Holland,
1988.

\bibitem{BK}
H. Becker and A.S. Kechris, {\it The Descriptive Set Theory of
Polish Group Actions}, London Math. Soc. Lecture Note Series, {\bf 232},
Cambridge Univ. Press, 1996.

\bibitem{BL} Y. Benyamini and J. Lindenstrauss,
{\it Geometric Nonlinear Functional Analysis,} Vol. 1,
Colloquium Publications {\bf 48}, AMS,
Providence, RI, xi+488 p., 2000.

\bibitem{Be}
S.K. Berberian, {\it Lectures in Functional Analysis and Operator
Theory}, Springer-Verlag, 1974.

\bibitem{Bo}
S.A. Bogatyi, Metrically homogeneous spaces, {\it Russian Math. Surveys}
{\bf 57}(2) (2002), 221--240.

\bibitem{Ca}
P.J. Cameron, {\it Oligomorphic Permutation Groups}, London Math.
Society Lecture Note Series, {\bf 152}, 1990.

\bibitem{Ch}
G.L. Cherlin, {\it The Classification of Countable Homogeneous
Directed Graphs and Countable Homogeneous \(n\)-tournaments}, Memoirs
of the Amer. Math. Soc., {\bf 131}, No. 621, 1998.

\bibitem{Co}
C. Constantinescu, {\it \(C^*\)-algebras, Vol. 2:  Banach Algebras
and Compact Operators}, North Holland, 2001.

\bibitem{De}
D. Devlin, {\it Some partition theorems and ultrafilters on \(\omega\)},
Ph.D. Thesis, Dartmouth College, 1979.

\bibitem{dV}
J. de Vries, {\it Elements of Topological Dynamics}, Kluwer, 1993.

\bibitem{Di} J. Dieudonn\'e, Sur la completion des groupes topologiques,
{\it C. R. Acad. Sci. Paris} {\bf 218} (1944), 774--776.

\bibitem{E60}
R. Ellis, Universal minimal sets, {\it Proc. Amer. Math. Soc.} {\bf
11} (1960), 540--543.

\bibitem{E69}
R. Ellis, {\it Lectures on Topological Dynamics}, W.A. Benjamin, 1969.

\bibitem{Eng}
R. Engelking, {\it General Topology}, Math. Monographs, {\bf 60},
PWN - Polish Scient. Publishers, Warsaw, 1977.

\bibitem{EH}
P. Erd\"{o}s and A. Hajnal,
Unsolved and solved problems in set theory,
{\it Proc. Symp. Pure Math.} {\bf 25} (1974), 269--287.

\bibitem{EHR65}
P. Erd\"{o}s, A. Hajnal and R. Rado, Partition relations for cardinal numbers, {\it Acta Math. Acad. Sci. Hung.}, {\bf 16}, (1965), 93--196.

\bibitem{ER56}
P. Erd\"{o}s and R. Rado, A partition calculus in set theory, {\it Bull. Amer. Math. Soc.}, {\bf 62} (1956), 427--489.

\bibitem{Fo97}
W.L. Fouch\'{e}, Symmetry and the Ramsey degree of posets, {\it Discrete Math.}
{\bf 167/168} (1997), 309--315.

\bibitem{Fo98}
W.L. Fouch\'{e}, Symmetries in Ramsey theory, {\it East-West J.
Math.}, {\bf 1} (1998), 43--60.

\bibitem{Fo99}
W.L. Fouch\'{e}, Symmetry and the Ramsey degrees of finite relational
structures, {\it J. Comb. Theory Ser. A} {\bf 85}(2) (1999), 135--147.

\bibitem{Fo99a}
W.L. Fouch\'{e}, Symmetries and Ramsey properties of trees, {\it
Discrete Math.} {\bf 197/198} (1999), 325--330.

\bibitem{Fr}
R. Fra\"{\i}ss\'{e}, Sur l'extension aux relations de quelques propriet\'es
des ordres, {\it Ann. Sci. \'Ecole Norm. Sup.} {\bf 71} (1954), 363--388.

\bibitem{GaP} F. Galvin and K. Prikry,
Borel sets and Ramsey's theorem, {\it J. Symb. Logic} {\bf 38} (1973),
193--198.

\bibitem{GiP}
T. Giordano and V. Pestov, Some extremely amenable groups, {\it C.R. Math.
Acad. Sci. Paris} {\bf 334}(4) (2002), 273--278.

\bibitem{GiP2}
T. Giordano and V. Pestov, Some extremely amenable groups related to
operator algebras and ergodic theory, ArXiv e-print math.OA/0405288,
May 2004, 37 pp.

\bibitem{Gl98}
E. Glasner, On minimal actions of Polish groups, {\it Top. Appl.}
{\bf 85} (1998), 119--125.

\bibitem{Gl00}
E. Glasner, Structure theory as a tool in topological dynamics, {\it Descriptive
Set Theory and Dynamical Systems}, M. Foreman et al. (Eds.), London Math.
Society Lecture Note Series, {\bf 277}, 173--209, Cambridge Univ. Press, 2000.

\bibitem{GlW02}
E. Glasner and B. Weiss, Minimal actions of the group \(S(\Z)\) of
permutations of the integers, {\it Geom. and Funct. Anal.} {\bf 12} (2002),
964--988.

\bibitem{GlW03}
E. Glasner and B. Weiss, The universal minimal system for the group
of homeomorphisms of the Cantor set, {\it Fundam. Math.}
{\bf 176} (2003), 277--289.

\bibitem{GLR}
R.L. Graham, K. Leeb, and B.L. Rothschild, Ramsey's theorem for a
class of categories, {\it Adv. in Math} {\bf 8} (1972), 417--433.

\bibitem{GR}
R.L. Graham and B.L. Rothschild, Ramsey's theorem for \(n\)-parameter
sets, {\it Trans. Amer. Math. Soc.} {\bf 159} (1971), 257--292.

\bibitem{GRS}
R.L. Graham, B.L. Rothschild, and J.H. Spencer, {\it Ramsey Theory},
2nd Edition, Wiley, 1990.

\bibitem{Gra} E. Granirer,
Extremely amenable semigroups I,II,
{\it Math. Scand.} {\bf 17} (1965), 177--179; {\bf 20} (1967), 93--113.

\bibitem{GraL}
E. Granirer and A.T. Lau, Invariant means on locally compact groups,
{\it Ill. J. Math}. {\bf 15} (1971), 249--257.

\bibitem{Gro}
M. Gromov, Filling Riemannian manifolds, {\it J. Diff. Geometry} {\bf 18} (1983),
1--147.

\bibitem{GroM}
M. Gromov and V.D. Milman, A topological application of the isoperimetric
inequality, {\it Amer. J. Math.} {\bf 105} (1983), 843--854.

\bibitem{HC}
W. Herer and J.P.R. Christensen, On the existence of pathological
submeasures and the construction of exotic topological groups,
{\it Math. Ann.} {\bf 213} (1975), 203--210.

\bibitem{H}
W. Hodges, {\it Model Theory}, Cambridge Univ. Press, 1993.

\bibitem{KPT05}
A.S. Kechris, V.G. Pestov, and S. Todorcevic, Fra\"{\i}ss\'{e} limits,
Ramsey theory, and topological dynamics of automorphism groups,
{\it Geom. Funct. Anal.} {\bf 15} (2005), 106--189.

\bibitem{KB}
D.E. Knuth and T.J. Buckholtz, Computation of Tangent, Euler, and Bernoulli
numbers, {\it Math. of Computation} {\bf 21}(100) (1967), 663--688.

\bibitem{LW}
A.H. Lachlan and R. Woodrow, Countable ultrahomogeneous undirected
graphs, {\it Trans. Amer. Math. Soc.} {\bf 262} (1980), 51--94.

\bibitem{LMP}
A.T.-M. Lau, P. Milnes, and J. Pym, On the structure of minimal
left ideals in the largest compactification of a locally compact group,
{\it J. London Math. Soc} {\bf 59}(2) (1999), 133--152.

\bibitem{L}
M. Ledoux, {\it The concentration of measure phenomenon}. Math. Surveys and
Monographs, {\bf 89}, Amer. Math. Soc., 2001.

\bibitem{Mill}
K. Milliken, A Ramsey theorem for trees, {\it J. Comb. Theory} {\bf 26} (1979),
215--237.

\bibitem{M67}
V.D. Milman, Infinite-dimensional geometry of the unit sphere in Banach space,
{\it Sov. Math., Dokl.} {\bf 8} (1967), 1440--1444.

\bibitem{M88}
V.D. Milman, The heritage of P. L\'{e}vy in geometrical functional analysis,
{\it Ast\'{e}risque} {\bf 157--158} (1988), 273--301.

\bibitem{MS}
V.D. Milman and G. Schechtman, {\it Asymptotic theory of finite-dimensional
normed spaces (with an Appendix by M. Gromov)},
Lecture Notes in Math., {\bf 1200}, Springer, 1986.

\bibitem{Mit} T. Mitchell,
Fixed points and multiplicative left invariant means,
{\it Trans. Am. Math. Soc.} {\bf 122} (1966), 195--202.

\bibitem{N89}
J. Ne\v{s}et\v{r}il, For graphs there are only four types of
hereditary Ramsey classes, {\it J. Comb. Theory, Series B} {\bf 46}(2)
(1989), 127--132.

\bibitem{N95}
J. Ne\v{s}et\v{r}il, Ramsey Theory, {\it Handbook of Combinatorics},
R. Graham et al. (Eds.), 1331--1403, Elsevier, 1995.

\bibitem{N03}
J. Ne\v{s}et\v{r}il, Private communication, 2003.

\bibitem{Nes04a}
J. Ne\v{s}et\v{r}il, Ramsey classes and homogeneous structures,
{\it Combin. Probab. Comput.} {\bf 14} (2005), 171--189.

\bibitem{N04}
J. Ne\v{s}et\v{r}il, Metric spaces are Ramsey,
{\it European J. Combin.} {\bf 28} (2007), 457--468.

\bibitem{NPRV}
J. Ne\v{s}et\v{r}il, H.J. Pr\"omel, V. R\"odl and B. Voigt, Canonical
ordering theorems, a first attempt, Proc. of the 10th Winter School on
Abstract Analysis (Srn\'\i, 1982), {\it Rend. Circ. Math. Palermo (2)} (1982),
193--197.

\bibitem{NR77}
J. Ne\v{s}et\v{r}il and V. R\"odl, Partitions of finite relational and
set systems, {\it J. Comb. Theory} {\bf 22}(3) (1977), 289--312.

\bibitem{NR78}
J. Ne\v{s}et\v{r}il and V. R\"odl, On a probabilistic graph-theoretical
method, {\it Proc. Amer. Math. Soc.} {\bf 72}(2) (1978), 417--421.

\bibitem{NR83}
J. Ne\v{s}et\v{r}il and V. R\"odl, Ramsey classes of set systems,
{\it J. Comb. Theory, Series A} {\bf 34}(2) (1983), 183--201.

\bibitem{NR84}
J. Ne\v{s}et\v{r}il and V. R\"odl, Combinatorial partitions of finite posets
and lattices---Ramsey lattices, {\it Algebra Universalis} {\bf 19} (1984),
106--119.

\bibitem{NR90}
J. Ne\v{s}et\v{r}il and V. R\"odl, {\it Mathematics of Ramsey Theory},
Springer, 1990.

\bibitem{Ng}
L. Nguyen Van Th\'{e}, Ramsey degrees of finite ultrametric spaces,
ultrametric Urysohn spaces and dynamics of their isometry groups,
{\it European J. Combin.} {\bf 30} (2009), 934--945.

\bibitem{OS}
E. Odell and T. Schlumprecht, The distortion problem,
{\it Acta Mathematica} {\bf 173} (1994), 259--281.

\bibitem{P98}
V. Pestov, Some universal constructions in abstract topological dynamics,
{\it Contemporary Math.} {\bf 215} (1998), 83--99.

\bibitem{P98a}
V. Pestov, On free actions, minimal flows and a problem by Ellis, {\it
Trans. Amer. Math. Soc.} {\bf 350}(10) (1998), 4149--4165.

\bibitem{P99}
V. Pestov, Topological groups:  where to from here?, {\it Topology
Proceed.} {\bf 24} (1999), 421--502.

\bibitem{P02}
V. Pestov, Ramsey-Milman phenomenon, Urysohn metric spaces, and extremely
amenable groups, {\it Israel J. Math.} {\bf 127} (2002), 317--357.

\bibitem{P02a}
V. Pestov, Remarks on actions on compacta by some infinite-dimensional
groups, {\it Infinite dimensional Lie groups in geometry and representation
theory, (Washington, DC 2000)}, World Sci., 145--163, 2002.

\bibitem{P02b}
V. Pestov, \(MM\)-spaces and group actions, {\it L' Enseign. Math.}
{\bf 48} (2002), 209--236.

\bibitem{PS}
M. Pouzet and N. Sauer, Edge partitions of the Rado graph,
{\it Combinatorica} {\bf 16}(4) (1996), 505--520.

\bibitem{Pr}
H.J. Pr\"omel, Some remarks on natural orders for combinatorial
cubes, {\it Discrete Math}. {\bf 73} (1989), 189--198.

\bibitem{Py}
J. Pym, A note on \(G^{\mathcal L\mathcal U\mathcal C}\) and Veech's Theorem, {\it Semigroup
Forum} {\bf 59} (1999), 171--174.

\bibitem{Ra}
R. Rado, Direct decomposition of partitions, {\it J. London Math. Soc.}
{\bf 29} (1954), 71--83.

\bibitem{RD} W. Roelcke and S. Dierolf, {\it Uniform Structures on Topological
Groups and Their Quotients,} McGraw-Hill, 1981.

\bibitem{Ru}
W. Rudin, {\it Functional Analysis}, McGraw-Hill, 1973.

\bibitem{S}
N. Sauer, Coloring finite substructures of countable structures,
Paul Erd\"{o}s and his Mathematics, II, {\it Bolyai Math. Studies,} {\bf 11}
(2002), 525--553.

\bibitem{Th}
S. Thomas, Groups acting on infinite dimensional projective spaces,
{\it J. London Math. Soc.} {\bf 34}(2) (1986), 265--273.

\bibitem{T}
S. Todorcevic, {\it Introduction to Ramsey Spaces}, Annals of Mathematics
Studies, vol.~174, Princeton University Press, 2010.

\bibitem{Tu}
S. Turek, Universal minimal dynamical system for reals, {\it Comment.
Math. Univ. Carolin.} {\bf 36}(2) (1995), 371--375.

\bibitem{U}
P. Urysohn, Sur un espace m\'etrique universel, {\it Bull. Sci.
Math.} {\bf 51} (1927), 43--64, 74--90.

\bibitem{Usp90}
V. Uspenskij, On the group of isometries of the Urysohn universal
metric space, {\it Comment. Univ. Carolinae} {\bf 31}(1) (1990), 181--182.

\bibitem{Usp00}
V. Uspenskij, On universal minimal compact \(G\)-spaces, {\it Topology
Proc.} {\bf 25} (2000), 301--308.

\bibitem{Usp02}
V. Uspenskij, Compactification of topological groups, {\it Proc. of the 9th
Prague Topol. Symp. (Prague 2001)}, 331--346 (electronic), Topol. Atlas,
Toronto, 2002.

\bibitem{V}
W. Veech, Topological dynamics, {\it Bull. Amer. Math. Soc.}
{\bf 83}(5) (1977), 775--830.

\end{thebibliography}

\end{document}
